\chapter{Multiletter Hankel Approximation}\label{sec:hankel}

\section{Introduction and the approximation question}
An intelligent system that receives observations over time must decide
what to retain from its history. The retained state should support future
predictions, remain consistent after new events, and fit the resources
available for its execution. These requirements arise in knowledge
representation and reasoning under uncertainty whenever an inference
procedure replaces an event history by a compact sequential model.
The relevant mathematical question is which distinctions between
histories must survive so that the specified future responses remain
accurate after further updates.

Predictive representations of state express this idea through predictions
of future observations~\citep{littman2001}. Weighted automata provide a
finite linear model of word-dependent responses, with a classical
equivalence between finite Hankel rank and minimal realization
dimension~\citep{carlyle1971,fliess1974}. Spectral methods connect these
representations to data, including learning from a single trajectory
under finite-state policies~\citep{balle2017}. Once a model is available,
reducing its state raises further questions about which output errors are
controlled, whether the reduced state defines a valid probability law,
and how many bits its recurrent implementation requires.

Recent work on agent memory makes the operational questions concrete.
Martin~\citep{martin2026memory} describes persistent external records
organized as transcripts, notes, linked representations, and context
loaded into later sessions. Offline consolidation can check notes against
the source history and revise their organization. This notion of memory
is distinct from updating model parameters or retaining activations
within an inference run. FrontierMemory~\citep{letta2026frontiermemoryv1}
evaluates whether memory derived from prior trajectories improves later
tasks, distinguishing task conventions, user preferences, and knowledge
of the working environment. Its first release studies offline memory
formation. These systems motivate a mathematical question: which
history distinctions must be retained to preserve specified future
answers through subsequent updates? Empirical task utility and a
uniform predictive certificate are different evaluation objects; their
relationship requires an explicit task model.

The distinction matters even in a simple sequence of revisions.
An event that grants a permission followed by one that revokes it has a
different effect from the reverse order. A representation intended to
answer later permission queries must preserve the distinction relevant
to those queries. The multiletter formalism keeps each event transition
and the order of its composition. For a scalar word response $f$, its
ordinary word-Hankel operator has entries $H_f[u,v]=f(uv)$, so a row
records the continuations associated with a past word $u$. The choice of
responses, transition model, and error metric determines which histories
may be merged.

This paper develops three levels of certified reduction. The first
concerns linear realizations and structured operator error. The second
requires a proper probabilistic model and controls discrepancies between
laws of completed words. The third requires an autonomous executable
state with an explicit update, readout, and storage contract.
These levels are connected by additional assumptions that are stated
where they are used. In particular, conditioning on a rare history and
implementing a real coordinate with finitely many bits each introduce
requirements beyond a small linear approximation error.

The contributions are organized into three sections.
{\clubpenalty=10000
\begin{enumerate}[label=(\roman*)]
\item \begin{samepage}Section~\ref{sec:hankel} gives an explicit binary family whose
singular-value tail vanishes while its best fixed-rank word-Hankel error
stays positive. It identifies the exact boundary limit, solves the
finite-parameter rank-one problem, and constructs matching first-order
responses. Joint Stein certificates yield finite algebraic descriptions.
For rational mean-square-stable sources and real competitors, decision
is polynomial in input bit length when alphabet size and target rank
are fixed.
\end{samepage}
\item \begin{samepage}Section~\ref{sec:probability} establishes sharp one-state
Kullback--Leibler approximation rates under exponential length tails
and higher moments. A mean bound alone permits a discontinuity, with a
separate construction for every fixed positive-state budget. Occupation
flows turn a supplied probability dictionary into a proper decoded law
with an explicit total-variation certificate. The acquisition results
specify the positivity, conditioning, sampling, and closure premises
under which that certificate can be obtained.
\end{samepage}
\item \begin{samepage}Section~\ref{sec:memory} characterizes residual state behind a
supplied suffix window and proves exact persistent-bit optima for
explicit systems. Its Fibonacci state count applies to the specified
initial leaf union. Mode checkpoints support replay through irreversible
transitions. Pressure bounds and finite-precision execution results
include the separation, observability, and arithmetic hypotheses needed
for their claims.
\end{samepage}
\end{enumerate}}

For artificial intelligence, these results give conditions under which a
compressed sequential representation supports a stated inference or
memory contract. Operator certificates concern the modeled continuation
responses. Probability certificates concern complete-word laws, with
additional conditions for rare-prefix predictions. Executable-memory
certificates account for persistent state and distinguish it from program
size, workspace, and update time. The dictionary and interface are part of
the mathematical input. The results are theoretical, and the final
recurrent-state application states a separate requirement for measured
gains in complete encoded length.

\paragraph{Research programme and deployment context.}
This work forms part of Datar Consulting's research on memory and
compression for large language models (LLMs), with sovereign model
deployments as an application objective. The programme studies
persistent agent memory, numerical recurrent state---including Gated
DeltaNet (GDN)~\citep{yang2025gated}---and exact residual encodings.
Here a sovereign deployment means that the operator controls the
infrastructure, retained data, and execution policy within a chosen
deployment boundary. The research objective is to reduce storage and
execution costs while preserving specified predictive behavior within
that boundary. An agent's external records and a decoder's recurrent
matrix state have different interfaces and correctness criteria.
The results below supply contracts for those interfaces under explicit
hypotheses; practical compression gains and complete-system behavior
remain implementation and evaluation questions.

For the operator problem, a weighted finite automaton represents a word
series by $f(w)=c^*A_wb$, where
$A_{a_1\cdots a_t}=A_{a_1}\cdots A_{a_t}$.
We use counting-measure $\ell^2(\Sigma^*)$ and define
\[
 \delta_k^{\mathbb F}(H_f)=
 \inf_{\substack{H_g\text{ bounded word-Hankel over }\mathbb F\\
                  \operatorname{rank}H_g\le k}}
                  \|H_f-H_g\|,\qquad \mathbb F\in\{\mathbb R,\mathbb C\}.
\]
The unconstrained lower bound is $\sigma_{k+1}(H_f)$ for compact $H_f$.
Classical one-letter AAK theory, including its automata formulation
\citep{lacroce2024}, gives exact optimal reduction under its hypotheses.
Singular value automata give canonical realizations and truncation bounds
\citep{sva2019}. A multiletter noncommutative approach was proposed in
\citep{lacroce2022}. Here the target is the ordinary counting-measure
word-Hankel structure. Subsection~\ref{I:U:sec:free-aak-obstruction}
checks the row-contraction premise needed for that cited approach.

The binary obstruction shows that a vanishing spectral tail alone cannot
control the structured optimum uniformly at fixed rank. Its stability
margin closes at the boundary, so uniform stability assumptions matter
when comparing it with other approximation guarantees. The exact
boundary constant has a scalar rational-approximation characterization.
Its identification with a reciprocal Hardy coefficient norm uses
published low-degree extremals through budget four and the stated
external outer-extremal input at general budget. The proofs identify
these inputs and keep the real and complex competitor classes explicit.

\section{A binary obstruction with an exact limit}
For distinct letters $a,b$, $k\ge1$ and $r\in(0,1)$, put $q=1-r^2$ and
\begin{equation}\label{I:eq:binary-source}
 f_{r,k}(a^j b^k)=qr^j\quad(j\ge0),\qquad
 f_{r,k}(w)=0\text{ otherwise}.
\end{equation}
The same notation is used throughout the proofs. Let
\[
 E_m^{\mathbb F}=\inf_{\substack{h\in H^2\text{ over }\mathbb F\\
                 \operatorname{rank}(h_{i+j})_{i,j\ge0}\le m}}
                 \|z^m-h\|_{H^2},\qquad E_0^{\mathbb F}=1,
\]
where circle measure is normalized. Write $E_m=E_m^{\mathbb C}$.

\begin{theorem}[Structural separation and boundary reduction]
\label{I:I:main:separation}
The source in \eqref{I:eq:binary-source} has rank $k+1$, norm one, and nonzero
singular values $1,\sqrt q,\ldots,\sqrt q$, with $k$ copies of $\sqrt q$.
For either field and every $0<r<1$,
\[
 \delta_k^{\mathbb F}(H_{f_{r,k}})\ge E_k>0.
\]
For each fixed budget and each field separately,
\[
 \lim_{r\uparrow1}\delta_k^{\mathbb F}(H_{f_{r,k}})=E_{k-1}^{\mathbb F}>0.
\]
Consequently no finite factor depending only on $(k,|\Sigma|)$ bounds
$\delta_k^{\mathbb F}/\sigma_{k+1}$ on all bounded word-Hankel sources.
\end{theorem}

The singular decomposition is explicit. Let
$z_j=\sqrt q\sum_{n\ge0}r^ne_{a^nb^j}$ for $0\le j\le k$.
These vectors are orthonormal and direct enumeration of the cuts of a word gives
\[
 H_{f_{r,k}}=z_0z_k^*+\sqrt q\sum_{j=0}^{k-1}z_{k-j}e_{b^j}^*.
\]
Both input and output systems in this decomposition are orthonormal.
The singular-value statement follows, and every usual finite-rank
spectral-tail norm tends to zero with $q$.

For the uniform floor, the row space of a rank-$k$ word-Hankel is
invariant under each backward letter shift. Its overlap with $z_k$ is
bounded by a scalar backward-shift invariant subspace problem. The
constant previously denoted by $\kappa_k$ has the single interpretation
\[
 \kappa_k=\Gamma_k^2=1-E_k^2,
 \qquad
 \Gamma_k=\max_{D^*D+l^*l\preceq I_k}\|lD^k\|.
\]
Thus the uniform floor at budget $k$ is $E_k$, while the boundary
limit is $E_{k-1}$. Full proofs are in
Theorems~\ref{I:U:thm:separation}, \ref{I:M:rev13:thm:uniform} and
\ref{I:B:thm:main}.

\subsection{Why one dimension is lost at the boundary}
Let $A,B$ be the backward shifts for $a,b$ restricted to a competitor's
row space, and let $lx=x[\varepsilon]$. Then
$A^*A+B^*B+l^*l\preceq I$. The exact identity
\[
 (I-rA^*)(I-rA)-q(I-A^*A)=(A-rI)^*(A-rI)
\]
implies $\|\sqrt q\,B(I-rA)^{-1}\|\le1$. By compactness, a sequence
of passive tuples has a limit. If $I-A_0$ is invertible, the relevant
singular-vector overlap tends to zero. Otherwise the fixed subspace of
the contraction $A_0$ is reducing, and passivity forces both $B_0$ and
$l_0$ to vanish on it. At least one dimension is therefore unavailable
to the remaining scalar extremal problem. This gives the lower bound
$E_{k-1}^{\mathbb F}$. A block realization lifting a scalar minimizer
provides the matching upper bound. The argument is over all bounded
rank-$k$ competitors and does not impose a common stability margin.

\section{The Hardy coefficient bridge and its inputs}
For $0<p<\infty$, define
$C(m,p)=\sup_{\|F\|_{H^p}=1}|[z^m]F|$.
The scalar rational class includes polynomial parts; its rank convention
is the rank of the coefficient Hankel $(h_{i+j})$, not a rational-degree
convention that discards those parts.

\begin{theorem}[Coefficient-norm identification]
\label{I:I:main:hardy}
For every $m\ge0$, $E_m^{\mathbb F}\ge1/C(m,2/3)$.
Equality over both fields for $0\le m\le3$ follows from the published
low-degree extremals in~\citet{brevig2020}. For general $m$, equality
over both fields uses the outer-extremal theorem of
\citet{outer2026}, together with the polynomial structure of coefficient
extremals in~\citet{brevig2020}.
\end{theorem}

The lower inequality has a short proof independent of that external
outer-extremal input. A degree-$m$ model space is $K_\Theta$, with
$\Theta=D_m^*/D$. Let $b=P_{\le m}\Theta$. Taylor coefficient extraction
and H\"older's inequality give
\[
 \|D\|_2^2=[z^m](bD^2)
 \le C(m,2/3)\|bD^2\|_{2/3}
 \le C(m,2/3)\|b\|_2\|D\|_2^2.
\]
The squared projection error for $z^m$ is
$\sum_{j=0}^m|\theta_j|^2=\|b\|_2^2$, proving the bound. Smaller
model spaces may be enlarged, so the same inequality covers all ranks
at most $m$.

For the upper bound, an outer coefficient extremal has the form
$P^3$, normalized by $\|P\|_2=1$ and $[z^m]P^3=C$.
Stationarity on the polynomial unit sphere yields
$P^2=CP_m^*+O(z^{m+1})$. Therefore
$\Theta=P_m^*/P=P/C+O(z^{m+1})$, and the projection error is $1/C$.
Nonnegative real coefficient extremals give a real attaining lift.
Theorem~\ref{I:B:thm:hardy-bridge} includes the full construction and the
precise external input; no numerical optimization is used as a proof.

\begin{center}
\begin{tabular}{cll}\toprule
Budget $k$ & Boundary error & Source of scalar input\\\midrule
1 & $1$ & Elementary\\
2 & $\sqrt3/2$ & Published degree-one extremal\\
3 & $3/4$ & Published degree-two extremal\\
4 & $\sqrt{(1103-33\sqrt{33})/2048}$ & Published degree-three extremal\\
General $k$ & $1/C(k-1,2/3)$ & Outer-extremal theorem~\cite{outer2026}\\\bottomrule
\end{tabular}
\end{center}
The classical model-space formulation is also used in rational
approximation~\citep{rational2026}. A numerical coincidence with a
constant in a different rational-degree problem is not an identification
of the approximation classes.

\section{Finite-parameter optimality and first-order response}
Write $D_k^{\mathbb F}(q)=\delta_k^{\mathbb F}(H_{f_{r,k}})^2$.
At rank one, let $u(q)$ be the unique root in $(q,\sqrt q)$ of
\[
 q^2(1-u^2)(u-q)^2=(1-q)^3u^2(q-u^2).
\]
The exact optimum is
\[
 D_1^{\mathbb R}(q)=D_1^{\mathbb C}(q)
 =\frac{q^2+(1-2q)u(q)^2}{q(1-u(q)^2)}.
\]
Theorem~\ref{I:F:thm:F1-exact} constructs and classifies every optimizer.
The phase reduction is an explicit two-sided unitary conjugation; the
real and complex equality here is proved, rather than inferred from
real-valued input data.

At every fixed budget $k\ge2$, the general response theorem gives
\[
 D_k^{\mathbb F}(q)=(E_{k-1}^{\mathbb F})^2-b_k^{\mathbb F}q+o(q).
\]
The coefficient is the attained finite variational maximum in
\eqref{I:R:eq:main-response-coefficient}; no universal closed formula or
uniform efficient optimizer is asserted. The diagonal real/complex
comparison uses the scalar-extremal classification identified in
Section~\ref{I:R:subsec:real-diagonal-orbits}. Explicit cases are
\begin{align*}
 D_2^{\mathbb F}(q)&=\frac34-\frac14q+O(q^2),\\
 D_3^{\mathbb F}(q)&=\frac9{16}-\frac17q+O(q^2),\\
 D_4^{\mathbb F}(q)&=\frac{1103-33\sqrt{33}}{2048}
       -\frac{855+103\sqrt{33}}{15872}q+o(q).
\end{align*}
The theorem also permits independent state budget and terminal delay,
and extensions to arbitrary terminal languages.

The proof separates two tasks. A passive-envelope calculation bounds
all structured competitors from below. A lossless construction then
produces an actual word-Hankel operator attaining that response to first
order. Passivity alone is insufficient for this second task. Joint
stability of the constructed realization follows from a finite
observation Gramian, and the improved $O(q^2)$ upper remainder follows
from an exact sign-conjugation symmetry of the analytic trial. These
steps are proved in Theorem~\ref{I:R:thm:v4-matching}. The global lower
remainder is only $o(q)$ without the additional structure established
in the explicit low-order cases.

\section{Finite certificates and computational scope}
Given a stable finite realization, a candidate's error realization has
joint reachability and observability Gramians. Their Stein equations and
a positive-semidefinite norm constraint give an exact finite certificate
for $\delta_k(H_f)\le t$. Attainment is essential for the non-strict
threshold. Factoring positive-semidefinite matrices yields a degree-three
existential real formula and algebraic optimal witnesses
(Theorem~\ref{I:U:thm:exact} and Corollary~\ref{I:U:cor:algebraic}).
This is an application of established realization, matrix-positivity and
real-algebraic tools, not a hardness result.

For rational mean-square-stable sources and real competitors, when
$d=|\Sigma|$ and $k$ are fixed, source-dependent Gramian variables
can be eliminated before applying real-algebraic decision algorithms.
There remain
$p=dk^2+2k+k(k+1)/2$ candidate variables; the number and degrees of the
polynomials grow polynomially with the source dimension. The resulting
bit-time bound has a polynomial exponent $O(p)$ for fixed parameters.
This is an XP result in $(d,k)$, not a fixed-parameter tractability claim
and not joint polynomiality in all dimensions
(Theorem~\ref{I:U:thm:fixed}). The constructive global net in
Theorem~\ref{I:F:thm:g1-net} gives a separate certified approximation scheme
with explicit but generally large costs.

\section{Limits and organization of the proofs}
The target space and norm are part of every claim. Ordinary counting-measure
backward letter shifts have row norm $\sqrt d$, so the row-contraction
hypothesis required by the cited noncommutative AAK route is not verified
for those shifts. This observation does not identify that model with all
Fock-space or bilateral-shift formulations.

Three further questions remain open here: useful complexity dependence
on varying rank; finite-parameter optimizers and field comparison beyond
the solved examples; and a uniform multiletter balanced-truncation tail
bound under explicit geometric hypotheses. The Stein-witness truncation
bounds proved below do not establish the classical
$2\sum_{j>k}\sigma_j$ estimate.

The remaining subsections contain complete proofs and constructions, including the
operator preliminaries, obstruction family, exact boundary reduction,
Hardy bridge, finite-parameter and response analysis, structured examples,
and numerical certificates. Repeated foundational results are cited
within the article rather than reproved under different names.
\label{I:I:main-end}

\section*{Detailed proofs and supporting results}

\section{Notation and finite stable sources}\label{I:U:sec:preliminaries}
Let $\Sigma$ be a finite alphabet of size $d\ge1$, let $\Sigma^*$ include
the empty word $\varepsilon$, and let $\UFF$ be either $\mathbb R$ or
$\mathbb C$. All operators act on the corresponding Hilbert space
$\UHH$ unless another space is specified. Its orthonormal basis is
$(e_w)_{w\in\Sigma^*}$. A word-Hankel operator has entries
\[
H_f[u,v]=f(uv),\qquad u,v\in\Sigma^*.
\]
The real and complex optimization problems are treated separately; no
equality between their optimal values is assumed. The notation $X^*$
means transpose in the real case and conjugate transpose in the complex
case. Positive semidefinite order is denoted by $\preceq$.

A finite realization of dimension $n$ consists of matrices
$A_a\in\UFF^{n\times n}$ and vectors $b_f,c_f\in\UFF^n$ such that
\begin{equation}\label{I:U:eq:realization}
f(w)=c_f^*A_wb_f,
\qquad A_{a_1\cdots a_j}=A_{a_1}\cdots A_{a_j},
\qquad A_\varepsilon=\UId_n.
\end{equation}
It is \emph{mean-square stable} when the positive linear map
\[
\Phi_A(X)=\sum_{a\in\Sigma}A_aXA_a^*
\]
has spectral radius less than one. This is stability of a finite linear
map, not the generally different condition on the joint spectral radius
of arbitrary products of the $A_a$.

Under this assumption the observation and reachability maps, with rows
$c_f^*A_u$ and columns $A_vb_f$, respectively, are bounded and satisfy
\begin{align}
P_f&=\sum_{w\in\Sigma^*}A_wb_fb_f^*A_w^*,
&P_f&=b_fb_f^*+\sum_aA_aP_fA_a^*,\label{I:U:eq:sourceP}\\
Q_f&=\sum_{w\in\Sigma^*}A_w^*c_fc_f^*A_w,
&Q_f&=c_fc_f^*+\sum_aA_a^*Q_fA_a.\label{I:U:eq:sourceQ}
\end{align}
These are convergent series of finite matrices; the Stein equations have
unique solutions. Writing the maps as $O_f,C_f$ gives
$H_f=O_fC_f$, $P_f=C_fC_f^*$, and $Q_f=O_f^*O_f$.
The nonzero squared singular values of $H_f$ are the nonzero eigenvalues
of $P_f^{1/2}Q_fP_f^{1/2}$. In particular,
\begin{equation}\label{I:U:eq:trace}
\Unorm{H_f}^2\le\Unorm{H_f}_{\mathrm{HS}}^2
=\Utr(P_fQ_f).
\end{equation}
These identities do not require invertible Gramians.

Throughout the algorithmic results, the source is supplied by a finite
stable realization, not by a table of finitely many coefficients or an
arbitrary infinite-dimensional oracle. Algebraic input entries are given
by effective exact representations. Bit-complexity statements concern
rational input, with numerators and denominators encoded in binary and
total description length denoted by $L$. For complex rational data, both
real and imaginary parts have this encoding. The threshold is rational
and included in $L$. A requested rank $k\ge n$ is trivial since the
source itself is feasible; $k=0$ means that the only competitor is zero.

For later use define the left and right creation operators by
\begin{equation}\label{I:U:eq:shifts}
L_ae_w=e_{aw},\qquad R_ae_w=e_{wa}.
\end{equation}
They obey
$\sum_aL_aL_a^*=\sum_aR_aR_a^*=\UId-e_\varepsilon e_\varepsilon^*$
and $R_a^*R_a=L_a^*L_a=\UId$. The word-Hankel identities are equivalently
\begin{equation}\label{I:U:eq:intertwining}
R_a^*H_f=H_fL_a\qquad(a\in\Sigma).
\end{equation}
Taking adjoints yields the dual identity
\begin{equation}\label{I:U:eq:adjoint-intertwining}
L_a^*H_f^*=H_f^*R_a\qquad(a\in\Sigma).
\end{equation}
In particular, the prefix residual $v\mapsto p(uv)$ of a real probability
law is the column $H_p^*e_u$, rather than $H_pe_u$.
For a complex series, $(H_f^*e_u)[v]=\overline{f(uv)}$ instead.


The structured approximation problem, with scalar field fixed, is
\begin{equation}\label{I:U:eq:problem}
 \delta_k(H_f)=\inf\{\|H_f-H_g\|:H_g\text{ bounded word-Hankel},\ \Urank H_g\le k\}.
\end{equation}
The empty word is denoted by $\varepsilon$; a positive error tolerance is
written $\epsilon$ or explicitly named when both occur.

\section{A binary obstruction at every rank budget}\label{I:U:sec:obstruction}
Use the source in \eqref{I:eq:binary-source}, with $q=1-r^2$.
The source changes with $r$, but the alphabet and target rank stay fixed.

\begin{theorem}[No uniform scalar AAK factor]\label{I:U:thm:separation}
The source has rank $k+1$ and singular values $1,\sqrt q,\ldots,\sqrt q$,
where $\sqrt q$ occurs $k$ times. Let $C_k=\binom{2k}{k}$. Over either
scalar field,
\[
 \delta_k(H_{f_{r,k}})\ge
 \max\left\{\sqrt q,\sqrt{\frac{1+q^2(C_k-1)}{C_k}}\right\}.
\]
In particular, $\sigma_{k+1}\to0$ but
$\liminf_{r\uparrow1}\delta_k\ge C_k^{-1/2}>0$. No constant depending
only on $k$ and the alphabet size can bound $\delta_k$ by that singular value.
\end{theorem}
\begin{proof}
Let $z_j=\sqrt q\sum_{n\ge0}r^n e_{a^n b^j}$ for $0\le j\le k$.
These are orthonormal. Directly sorting the possible cuts of $a^n b^k$ gives
\[
 H_{f_{r,k}}=z_0z_k^*+\sqrt q\sum_{j=0}^{k-1}z_{k-j}e_{b^j}^*.
\]
The right vectors $z_k,e_\varepsilon,e_b,\ldots,e_{b^{k-1}}$ are
orthonormal, as are the corresponding left vectors. This proves the rank and
spectrum, and the usual unstructured rank lower bound gives $\sqrt q$.

Consider any bounded rank-$s$ word-Hankel $G$, where $s\le k$. Its finite
range is invariant under the right backward shifts. The canonical realization
on that range (proved in Lemma~\ref{I:U:lem:compact}) has letter matrices of norm
at most one. The characteristic polynomial of its $b$ matrix, multiplied
by $z^{k-s}$, is a monic annihilating polynomial
$p(z)=z^k+\sum_{j=0}^{k-1}c_jz^j$. Its roots have modulus at most one.
The elementary symmetric coefficient bounds and Vandermonde's identity yield
\[
 S:=\sum_{j=0}^{k-1}|c_j|^2\le\sum_{j=0}^{k-1}\binom{k}{j}^2=C_k-1.
\]
For each $n$, the columns indexed by $a^n b^j$ obey the annihilating
relation. Boundedness permits summation of the convergent vectors $z_j$,
and therefore $Gx=0$ for $x=z_k+\sum_{j<k}c_jz_j$.
The displayed source decomposition gives
\[
 \|x\|^2=1+S,\qquad
 H_{f_{r,k}}x=z_0+q\sum_{j<k}c_jz_{k-j},\qquad
 \|H_{f_{r,k}}x\|^2=1+q^2S.
\]
Since $(1+q^2S)/(1+S)$ decreases with $S\ge0$, its value is at least
$(1+q^2(C_k-1))/C_k$. Taking the infimum over $G$ proves the claim.
The zero competitor is included by the polynomial $z^k$.
\end{proof}


\begin{remark}[Scope of the obstruction]
For rational $r$, the target has a nonnegative rational realization: a waiting
state loops on $a$ with weight $r$, followed by a chain of $k$ transitions on
$b$, with output scale $q$. It is minimal by the spectrum above and is
mean-square stable. Its stability margin tends to zero as $r\uparrow1$.
Thus a positive theorem depending on a fixed stability margin or additional
geometry is not excluded. The target is not a normalized word probability
law. The theorem is neither a computational hardness statement nor a
persistent-bit lower bound.
\end{remark}
\subsection{The row-contraction hypothesis in noncommutative AAK theory}\label{I:U:sec:free-aak-obstruction}

The preceding separation does not contradict Popescu's multivariable
AAK theorem \cite{popescu2003}.  The required hypothesis is stronger than
the Hankel intertwining identity alone.  The following calculation makes
the distinction explicit.

\begin{proposition}[Target-shift obstruction]\label{I:U:prop:rowobstruction}
On counting-measure word space, let
$L_ae_w=e_{aw}$ and $R_ae_w=e_{wa}$.  Every bounded word-Hankel satisfies
\[
 H_fL_a=R_a^*H_f.
\]
For an alphabet of size $d$, however,
\begin{equation}\label{I:U:eq:row-shift-norm}
 \Unorm{[R_1^*\ \cdots\ R_d^*]}=\sqrt d.
\end{equation}
Thus the tuple $(R_a^*)_{a\in\Sigma}$ is not a row contraction when
$d\ge2$.
\end{proposition}
\begin{proof}
The intertwining identity follows by evaluating both sides at row $u$
and column $v$: each entry is $f(uav)$.  Each $R_a$ is an isometry, so
\[
 [R_1^*\ \cdots\ R_d^*]
 [R_1^*\ \cdots\ R_d^*]^*
 =\sum_aR_a^*R_a=d\UId.
\]
This proves \eqref{I:U:eq:row-shift-norm}.  Equivalently, taking
$h_a=R_ax$ with $x\ne0$ gives
$\Unorm{\sum_aR_a^*h_a}^2=d^2\Unorm{x}^2$, whereas
$\sum_a\Unorm{h_a}^2=d\Unorm{x}^2$.
\end{proof}


Popescu's hypotheses require a row-expansive source tuple $X$ and a
row-contractive target tuple $T$, namely
\[
 \Unorm{\sum_aX_ay_a}^2\ge\sum_a\Unorm{y_a}^2,
 \qquad
 \Unorm{\sum_aT_ah_a}^2\le\sum_a\Unorm{h_a}^2.
\]
The resulting theorem is stated in terms of restrictions to
$X$-invariant subspaces of bounded codimension and, equivalently,
multipliers on those subspaces
\cite[Theorems~5.3--5.4]{popescu2003}.
The left creators satisfy the first inequality, but
Proposition~\ref{I:U:prop:rowobstruction} rules out the second for the
ordinary backward target shifts.  Unitary word reversal preserves
the norm in \eqref{I:U:eq:row-shift-norm}; a dilation cannot fix this by
compression either, since compressions of row contractions remain
row contractions.

The comparison with the noncommutative framework of
\cite[Theorems~7--10]{lacroce2022} must keep its operator conventions
separate from the counting-measure model used here. For the latter model,
the hypotheses of Popescu's theorem are not verified by the ordinary
backward target shifts: their row norm is $\sqrt d$, and their ranges
are not pairwise orthogonal. Proposition~\ref{I:U:prop:rowobstruction}
and Theorem~\ref{I:U:thm:separation} are the precise conclusions needed
here. They do not address a different Fock-space or bilateral-shift
formulation without a separate identification of its norm and rank classes.

\section{Attainment and finite-dimensional realizations}\label{I:U:sec:realization}
We first separate existence from computation. Existence uses only the operator topology; the realization results subsequently turn finite-rank competitors into finite certificates.

\begin{theorem}[Attainment for every bounded source]\label{I:U:thm:attainment}\label{I:S:thm:attainment}
Let $H_f$ be any bounded operator on word space (in particular a bounded word-Hankel operator), and let $k\ge0$. There exists a bounded word-Hankel operator $G_\star$ with $\Urank G_\star\le k$ such that
\[
 \Unorm{H_f-G_\star}=\delta_k(H_f).
\]
The assertion holds separately over $\mathbb R$ and $\mathbb C$, without any finite-rank or stability assumption on the source.
\end{theorem}
\begin{proof}
Choose a minimizing sequence $(G_j)$ of admissible operators with
$\Unorm{H_f-G_j}\le\delta_k(H_f)+1$. Since the zero operator is admissible,
$\delta_k(H_f)\le\Unorm{H_f}$ and hence
$\Unorm{G_j}\le 2\Unorm{H_f}+1=:M$.
The closed ball of radius $M$ in $B(\UHH)$ is compact in the weak operator topology. In the present separable space, a subsequence can be obtained directly by diagonal extraction of all matrix coefficients: the uniform norm bound extends the limiting coefficients to a bounded operator $G$, and then finite-support approximation gives weak operator convergence on arbitrary pairs of vectors.

The word-Hankel identities are equalities between matrix coefficients and therefore pass to this limit. Every $(k+1)\times(k+1)$ minor also passes to the limit and remains zero, so $\Urank G\le k$. Indeed, rank greater than $k$ would give $k+1$ independent columns and then a nonzero coordinate minor of that order. Finally, for any unit vectors $x,y$,
\[
 |\langle(H_f-G)x,y\rangle|
 =\lim_j|\langle(H_f-G_j)x,y\rangle|
 \le\liminf_j\Unorm{H_f-G_j}.
\]
Taking the supremum over $x,y$ proves the result. For real operators, perform the same argument in the real Hilbert space.
\end{proof}


The following normalization packages bounded realizations into a compact parameter set. It uses the familiar reachability and observability factorization of a rational Hankel operator; see \cite{sva2019} for the weighted-automata setting. Its formulation with non-strict inequalities also admits singular and redundant realizations.

For $M>0$ and $k\ge1$, let $\mathcal K_{k,M}$ consist of tuples
$\theta=((T_a)_{a\in\Sigma},b,c,P)$, with $T_a\in\mathbb F^{k\times k}$,
$b,c\in\mathbb F^k$ and $P=P^*$, satisfying
\begin{align}
 \sum_a T_a^*T_a+cc^*&\preceq\UId_k,\label{I:U:eq:normal-observation}\\
 \sum_a T_aPT_a^*+bb^*&\preceq P,\label{I:U:eq:normal-reachability}\\
 0\preceq P&\preceq M^2\UId_k.\label{I:U:eq:normal-bound}
\end{align}
Here $\mathbb F$ is the chosen scalar field. Define $g_\theta(w)=c^*T_wb$.

\begin{lemma}[Compact and complete parametrization]\label{I:U:lem:compact}
The set $\mathcal K_{k,M}$ is compact, and
\[
 \{H_{g_\theta}:\theta\in\mathcal K_{k,M}\}
 =\{G:G\text{ is bounded word-Hankel},\ \Urank G\le k,\ \Unorm G\le M\}.
\]
\end{lemma}
\begin{proof}
For the forward inclusion, introduce the positive maps
\[
 \Psi(X)=\sum_aT_a^*XT_a,
 \qquad \Phi_T(X)=\sum_aT_aXT_a^*.
\]
Iterating \eqref{I:U:eq:normal-observation} and
\eqref{I:U:eq:normal-reachability} gives the explicit partial-sum bounds
\begin{align*}
 \sum_{j=0}^{N}\Psi^j(cc^*)
 &\preceq \UId_k-\Psi^{N+1}(\UId_k)\preceq\UId_k,\\
 \sum_{j=0}^{N}\Phi_T^j(bb^*)
 &\preceq P-\Phi_T^{N+1}(P)\preceq P.
\end{align*}
The summands are respectively
$\sum_{|u|=j}T_u^*cc^*T_u$ and
$\sum_{|v|=j}T_vbb^*T_v^*$. Thus the matrix with rows $c^*T_u$ defines a bounded observation operator $O:\mathbb F^k\to\UHH$, and the matrix with columns $T_vb$ defines a bounded reachability operator $C:\UHH\to\mathbb F^k$, with
\[
 O^*O\preceq\UId_k,\qquad CC^*\preceq P.
\]
For $C$, one may first define $C^*x=(b^*T_v^*x)_v$ and use the second bound, avoiding any assumption of absolute convergence of an arbitrary input sum. Since
$(OC)[u,v]=c^*T_uT_vb=g_\theta(uv)$, the product is the desired bounded Hankel, and its rank is at most $k$ and its norm at most $M$.

Conversely, let $G$ be a nonzero bounded word-Hankel of rank $s\le k$. Its range $S$ is finite-dimensional and closed. Choose an isometry
$O:\mathbb F^s\to\UHH$ onto $S$. Define the left and right creation operators by
\[
 L_ae_v=e_{av},\qquad R_ae_u=e_{ua}.
\]
The Hankel identity is $R_a^*G=GL_a$, so $S$ is invariant under $R_a^*$. Set
\[
 T_a=O^*R_a^*O,\quad c=O^*e_\varepsilon,
 \quad C=O^*G,\quad b=Ce_\varepsilon.
\]
Then $OT_a=R_a^*O$ and $T_aC=CL_a$. To check the word order explicitly, if $o_u$ denotes row $u$ of $O$, then $o_{ua}=o_uT_a$ and $o_\varepsilon=c^*$. Thus $o_u=c^*T_u$. Similarly, column $av$ of $C$ equals $T_a$ times column $v$, so column $v$ is $T_vb$. Hence $G=OC$ has the prescribed realization.

The ranges of the $R_a$ partition the nonempty words. Therefore
\[
 \sum_aT_a^*T_a+cc^*
 =O^*\left(\sum_aR_aR_a^*+e_\varepsilon e_\varepsilon^*\right)O
 =\UId_s.
\]
Taking $P=CC^*$ and partitioning the columns of $C$ by their first letter gives
\[
 P=bb^*+\sum_aT_aPT_a^*,\qquad
 P\preceq\Unorm{G}^2\UId_s\preceq M^2\UId_s.
\]
Pad all state variables by zeros to reach dimension $k$. The zero Hankel is represented by the zero tuple.

Finally, the constraints imply
$\Unorm{T_a}\le1$, $\Unorm c\le1$, $\Unorm b\le M$, and
$\Unorm P\le M^2$. They are closed conditions in a finite-dimensional real vector space, including after separating real and imaginary parts. Hence the parameter set is compact.
\end{proof}


Although the compact set permits boundary modes, every output it represents also has a strictly stable realization of the same allowed dimension.

\begin{lemma}[Stable realization of a bounded finite-rank Hankel]\label{I:U:lem:stable}
Every bounded word-Hankel of rank at most $k$ has a realization of dimension $k$ satisfying $\rho(\Phi_T)<1$. A nonzero operator of rank $s$ has such a realization of dimension $s$.
\end{lemma}
\begin{proof}
Use the canonical $s$-dimensional realization in the converse proof of Lemma~\ref{I:U:lem:compact}. Its observation operator is an isometry, so
\[
 \sum_{w\in\Sigma^*}T_w^*cc^*T_w=\UId_s.
\]
The identity $\UId_s=cc^*+\Psi(\UId_s)$ then yields
\[
 \Psi^N(\UId_s)=\UId_s-\sum_{|w|<N}T_w^*cc^*T_w\longrightarrow0
\]
in matrix norm. For a Hermitian matrix $X$, positivity and
$-\Unorm X\UId_s\preceq X\preceq\Unorm X\UId_s$ imply
$\Unorm{\Psi^N(X)}\le\Unorm X\,\Unorm{\Psi^N(\UId_s)}$.
Splitting a general complex matrix into Hermitian real and imaginary parts shows that the powers of $\Psi$ tend to zero as linear maps on the entire matrix space. In finite dimension this is equivalent to $\rho(\Psi)<1$.
For real realizations, apply the same argument to the complexification of the real matrix map.
The map $\Phi_T$ is the adjoint of $\Psi$ for the Hilbert--Schmidt inner product, and hence has the same spectral radius. Zero padding preserves strict stability; the zero operator has a zero-transition realization.
\end{proof}


\begin{remark}
Theorem~\ref{I:U:thm:attainment} is a compactness argument, not an algorithm. Lemma~\ref{I:U:lem:stable} concerns the existence of a stable representation, not the stability of every redundant representation of the same series. In particular, the closed constraints defining $\mathcal K_{k,M}$ intentionally allow unreachable or unobservable boundary modes.
\end{remark}

\section{An exact finite certificate}\label{I:U:sec:exact}
Assume now that the source is supplied by the finite mean-square-stable realization
$f(w)=c_f^*A_wb_f$ of dimension $n$. Let the candidate have $k$ states and put $m=n+k$. Its joint error realization is
\begin{equation}\label{I:U:eq:error-realization}
 B_a=\begin{pmatrix}A_a&0\\0&T_a\end{pmatrix},\qquad
 \beta=\begin{pmatrix}b_f\\b\end{pmatrix},\qquad
 \gamma=\begin{pmatrix}c_f\\-c\end{pmatrix}.
\end{equation}
Thus $\gamma^*B_w\beta=f(w)-g(w)$. No damping is applied to either block.

The next theorem is a Hankel-norm certificate built from reachability and observability Stein inequalities. The two ingredients have standard precedents in weighted-automata Gramians \cite{sva2019} and realization theory over free semigroups \cite{ball2005}. We prove the equivalence directly, including singular Gramians, rather than identifying it with a transfer-function bounded-real lemma.

\begin{theorem}[Exact joint Stein certificate]\label{I:U:thm:exact}
For $t\ge0$, the inequality $\delta_k(H_f)\le t$ holds if and only if there exist candidate parameters $(T_a,b,c)$, Hermitian matrices $P,Q\in\mathbb F^{m\times m}$ and $W\in\mathbb F^{m\times m}$ such that
\begin{align}
 P&\succeq0, & Q&\succeq0,\label{I:U:eq:exact-positive}\\
 P-\beta\beta^*-\sum_aB_aPB_a^*&\succeq0,\label{I:U:eq:exact-reach}\\
 Q-\gamma\gamma^*-\sum_aB_a^*QB_a&\succeq0,\label{I:U:eq:exact-observe}\\
 P&=WW^*,\label{I:U:eq:exact-factor}\\
 t^2\UId_m-W^*QW&\succeq0.\label{I:U:eq:exact-norm}
\end{align}
Every feasible tuple defines a bounded word-Hankel $H_g$ of rank at most $k$ with $\Unorm{H_f-H_g}\le t$. At $t=\delta_k(H_f)$ a feasible tuple exists.
\end{theorem}
\begin{proof}
Suppose first that a tuple is feasible. Iterating
\eqref{I:U:eq:exact-reach}--\eqref{I:U:eq:exact-observe} gives
\[
 \sum_{|w|\le N}B_w\beta\beta^*B_w^*\preceq P,
 \qquad
 \sum_{|w|\le N}B_w^*\gamma\gamma^*B_w\preceq Q
\]
for every $N$. Exactly as in Lemma~\ref{I:U:lem:compact}, these inequalities define bounded reachability and observation maps $C_E$ and $O_E$, with exact Gramians
\[
 P_E=C_EC_E^*\preceq P,\qquad Q_E=O_E^*O_E\preceq Q.
\]
Restriction to the candidate state coordinates also gives bounded candidate maps $C_g,O_g$: the lower diagonal blocks of the partial sums are precisely their partial Gramians. Thus $H_g=O_gC_g$ is bounded and has rank at most $k$, and
$E:=H_f-H_g=O_EC_E$.

For bounded operators factoring through an $m$-dimensional space,
\begin{equation}\label{I:U:eq:gram-norm}
 \Unorm E^2=\lambda_{\max}(Q_E^{1/2}P_EQ_E^{1/2}).
\end{equation}
Indeed, $E^*E=C_E^*Q_EC_E$ has the same nonzero eigenvalues as $Q_E^{1/2}C_EC_E^*Q_E^{1/2}$; the equality also holds when both operators are zero. Congruence preserves the positive-semidefinite order, so
\begin{align*}
 \Unorm E^2
 &\le\lambda_{\max}(Q_E^{1/2}P Q_E^{1/2})\\
 &=\lambda_{\max}(W^*Q_EW)
 \le\lambda_{\max}(W^*QW)\le t^2.
\end{align*}
The middle equality follows from $P=WW^*$ and equality of the nonzero spectra of $ZZ^*$ and $Z^*Z$, with $Z=Q_E^{1/2}W$. No inverse or positive definiteness is used.

Conversely, Theorem~\ref{I:U:thm:attainment} supplies an optimal bounded rank-$\le k$ Hankel. Choose the canonical realization of Lemma~\ref{I:U:lem:compact}, padded with zeros, and form \eqref{I:U:eq:error-realization}. The source maps are bounded by its mean-square stability; the candidate maps are bounded by the lemma. Consequently
\[
 C_E=\begin{pmatrix}C_f\\ C_g\end{pmatrix},\qquad
 O_E=\begin{pmatrix}O_f&-O_g\end{pmatrix}
\]
are bounded. Their finite-dimensional Gramians $P_E,Q_E$ satisfy the exact Stein equations obtained by splitting words at the first and last letter, respectively. Set $P=P_E$, $Q=Q_E$ and choose $W=P_E^{1/2}$. Equation~\eqref{I:U:eq:gram-norm} then implies
$W^*QW\preceq\delta_k(H_f)^2\UId_m\preceq t^2\UId_m$.
This proves necessity and feasibility at the optimum.
\end{proof}


\begin{corollary}[Algebraic optimization and an existential formulation]\label{I:U:cor:algebraic}\label{I:S:thm:etr}
For a finite mean-square-stable source with algebraic coefficients, $\delta_k(H_f)$ is a nonnegative real algebraic number, and an optimal candidate exists with algebraic coefficients. Both can be computed by real quantifier elimination and algebraic sample-point extraction.
For rational source data and a rational threshold, the decision problem $\delta_k(H_f)\le t$ has a polynomial-size existential formula over the reals, of degree at most three, using
\[
 O\bigl(dk^2+(n+k)^2\bigr)
\]
real scalar variables. In particular, the rational-input decision problem belongs to $\exists\mathbb R$.
\end{corollary}
Here $\exists\mathbb R$ denotes the class of decision problems admitting
a polynomial-time many-one reduction to existential feasibility of
polynomial equalities and inequalities over the reals; see
\citet{schaefer2024}. Membership is an upper bound, not a hardness claim.
\begin{proof}
Replace every matrix inequality $Y\succeq0$ in Theorem~\ref{I:U:thm:exact} by an equation $Y=UU^*$ with a fresh square matrix $U$ of the appropriate size; Hermitian variables can be represented by their independent real coordinates. The condition $P\succeq0$ is already implied by $P=WW^*$. For complex data, separate real and imaginary parts. All resulting scalar polynomial equations and inequalities have degree at most three. The cubic terms are $B_aPB_a^*$, $B_a^*QB_a$, and $W^*QW$; all PSD factorizations themselves are quadratic. The number of matrices of order $m$ is constant, yielding the stated variable count. Expanding the sums and products produces polynomially many terms in $d,n,k$ and in the binary size of rational input coefficients.

By Theorem~\ref{I:U:thm:exact}, the projection of this feasible set onto $t$ is exactly the closed ray $[\delta_k(H_f),\infty)$. Quantifier elimination over the real algebraic numbers gives a univariate semialgebraic description of this ray, whose finite endpoint is algebraic. Applying sample-point extraction at that endpoint produces an algebraic feasible tuple and hence an optimal algebraic realization. These are standard effective operations over real closed fields \cite{basu2006,renegar1992}.

For rational input the threshold formula is an existential real formula of polynomial size, which proves membership in $\exists\mathbb R$. The stipulated source stability can itself be checked in polynomial bit time: solve $Y-\sum_aA_aYA_a^*=\UId_n$ over the rationals and test $Y\succ0$. A singular linear system is rejected; a positive-definite solution implies stability by conjugating the positive map with $Y^{1/2}$, and stability gives a unique such solution. Thus invalid source encodings can be rejected before the reduction. If $k$ is supplied in binary, the case $k\ge n$ is trivial for $t\ge0$, because the source itself is admissible. Otherwise $k<n$, so the construction is polynomial in the input size. Negative thresholds are rejected before constructing the formula. The case $k=0$ has the sole candidate $G=0$ and is handled by the same Gramian norm test without candidate variables.
\end{proof}


\begin{remark}\label{I:U:rem:exact-scope}
The existential formulation is an exact decision and synthesis description, not a claim of polynomial-time solvability when all parameters vary. It also does not assert equality between the real and complex optima: the chosen scalar field is fixed throughout each instance. Algebraic output describes an exact optimizer; rounding it to a prescribed finite precision requires a separate error analysis. No finite-horizon or numerical optimization oracle is assumed in the certificate.
\end{remark}

\section{Polynomial-time decision for fixed alphabet and target rank}
\label{I:U:sec:complexity}

The cubic certificate of Theorem~\ref{I:U:thm:exact} uses a number of
variables that grows with the source dimension.  For a fixed alphabet and
target rank, one can eliminate the Gramian variables before applying real
algebraic decision algorithms.  The resulting formulation has a fixed number
of variables and polynomially growing degrees.  Throughout this section the
source, threshold, and requested accuracy are rational, and competitors are
real.  The complexity model counts bit operations on binary encodings of
rational numbers.

\begin{theorem}[Fixed-parameter slices]\label{I:U:thm:fixed}
Fix integers $d\ge1$ and $k\ge0$.  Given a rational, mean-square-stable
$n$-state realization of $f$ and a rational threshold $t$, deciding whether
$\delta_k(H_f)\le t$ takes deterministic time polynomial in the binary
input length.  The exponent may depend on $d$ and $k$.
\end{theorem}

The theorem concerns polynomial time on each slice with $d,k$ fixed.
The construction below gives an XP algorithm with respect to $(d,k)$;
it does not give a fixed-parameter tractable algorithm or a running time
polynomial jointly in $d,k,n$. With $p=dk^2+2k+k(k+1)/2$ candidate variables, the semialgebraic decision step has exponent $O(p)=O(dk^2)$ in the polynomial-size input parameters; constants and prefactors may depend on $(d,k)$.

\begin{lemma}[A real-rooted polynomial test]\label{I:U:lem:coefftest}
Suppose a real $m\times m$ matrix $R$ has only real eigenvalues.
All of its eigenvalues are nonnegative if and only if all nonleading
coefficients of $\det(z\UId_m+R)$ are nonnegative.
\end{lemma}
\begin{proof}
Write the eigenvalues, with algebraic multiplicity, as
$\mu_1,\ldots,\mu_m$.  If each $\mu_i\ge0$, the coefficients of
$\prod_i(z+\mu_i)$ are nonnegative.  Conversely, if some $\mu_i<0$,
this polynomial has the positive root $-\mu_i$.  A monic polynomial
whose remaining coefficients are nonnegative is strictly positive at
every positive real number.  This is a contradiction.  Diagonalizability
of $R$ is not needed.
\end{proof}


\begin{proof}[Proof of Theorem~\ref{I:U:thm:fixed}]
Negative thresholds are immediately rejected.  We give the proof for
$k\ge1$; for $k=0$ the sole competitor is zero and the same calculation
below has no candidate variables.

\paragraph{A fixed number of candidate variables.}
By Theorem~\ref{I:U:thm:attainment} and Lemma~\ref{I:U:lem:stable}, the inequality
$\delta_k(H_f)\le t$ is equivalent to the existence of a stable
$k$-state realization $g(w)=c^\top T_wb$ with error at most $t$.
A smaller stable realization can be padded with zero states.
Introduce the entries of $T_1,\ldots,T_d,b,c$, and a symmetric
$k\times k$ matrix $X$ as variables.  Impose
\begin{equation}\label{I:U:eq:fixed-stability}
 X\succ0,\qquad X-\sum_{a\in\Sigma}T_aXT_a^\top=\UId_k.
\end{equation}
These conditions are equivalent to mean-square stability of the tuple
$T$.  For the forward implication, positivity and
$\UId_k\succeq \Unorm{X}^{-1}X$ give a strict order contraction in the
order-unit norm determined by $X$; the spectral radius is therefore less
than one.  For the reverse implication, take
$X=\sum_{j\ge0}\Phi_T^j(\UId_k)$.
Positive definiteness in~\eqref{I:U:eq:fixed-stability} is expressed by
positivity of the $k$ leading principal minors of $X$.
The number of real variables is
\begin{equation}\label{I:U:eq:fixed-varcount}
 p=dk^2+2k+\frac{k(k+1)}2,
\end{equation}
which is independent of $n$.

\paragraph{Eliminating the joint Gramians.}
Let $m=n+k$, and form
\[
 B_a=\begin{pmatrix}A_a&0\\0&T_a\end{pmatrix},\qquad
 \beta=\begin{pmatrix}b_f\\b\end{pmatrix},\qquad
 \gamma=\begin{pmatrix}c_f\\-c\end{pmatrix}.
\]
The joint tuple $B$ is stable: a positive definite Stein solution for
the source and the matrix $X$ in~\eqref{I:U:eq:fixed-stability}, placed on
the diagonal, give a positive definite Stein solution for $B$.
Hence its exact reachability and observability Gramians are uniquely
defined by
\[
 P=\beta\beta^\top+\sum_aB_aPB_a^\top,\qquad
 Q=\gamma\gamma^\top+\sum_aB_a^\top QB_a.
\]
These are the usual finite-dimensional Gramian equations for weighted
automata; see also~\cite{sva2019}.
Using column vectorization, put
\begin{align}
 \mathcal L&=\UId_{m^2}-\sum_a B_a\otimes B_a,
 &\Delta&=\det\mathcal L,\label{I:U:eq:fixed-L}\\
 \widehat P
 &=\operatorname{unvec}\!\left(
       \operatorname{adj}(\mathcal L)\operatorname{vec}(\beta\beta^\top)
     \right),&&\nonumber\\
 \widehat Q
 &=\operatorname{unvec}\!\left(
       \operatorname{adj}(\mathcal L^\top)\operatorname{vec}(\gamma\gamma^\top)
     \right).&&
 \label{I:U:eq:fixed-adj}
\end{align}
Stability guarantees $\Delta\ne0$ throughout the feasible set, and
\begin{equation}\label{I:U:eq:fixed-rational-grams}
 P=\widehat P/\Delta,\qquad Q=\widehat Q/\Delta.
\end{equation}
Thus the Gramians are rational functions of the fixed collection of
variables; they need not be introduced as additional variables.

\paragraph{A polynomial norm constraint without square roots.}
The squared norm of the error Hankel is
$\lambda_{\max}(P^{1/2}QP^{1/2})$.
The matrices $PQ$ and $P^{1/2}QP^{1/2}$ have the same characteristic
polynomial: apply
$\det(z\UId-UV)=\det(z\UId-VU)$ with
$U=P^{1/2}$ and $V=P^{1/2}Q$.
In particular, all eigenvalues of $PQ$ are real and nonnegative,
including when $P$ is singular.
Define the polynomial matrix
\begin{equation}\label{I:U:eq:fixed-R}
 R=t^2\Delta^2\UId_m-\widehat P\widehat Q
   =\Delta^2(t^2\UId_m-PQ).
\end{equation}
Its eigenvalues are real.  Since $\Delta^2>0$, the error norm is at
most $t$ if and only if these eigenvalues are nonnegative.  By
Lemma~\ref{I:U:lem:coefftest}, this is equivalent to
\begin{equation}\label{I:U:eq:fixed-coeff}
 \det(z\UId_m+R)=z^m+a_1z^{m-1}+\cdots+a_m,
 \qquad a_1\ge0,\ldots,a_m\ge0.
\end{equation}
Each $a_i$ is a polynomial in the variables counted
in~\eqref{I:U:eq:fixed-varcount}.  Squaring the common denominator avoids
any assumption about its sign.  No matrix square root, eigenvalue
variable, or list of exponentially many principal minors is used.

\paragraph{Degrees, coefficient sizes, and construction time.}
Each entry of the $m^2\times m^2$ matrix $\mathcal L$ has degree at most two.
Consequently
\[
 \deg\Delta\le2m^2,\qquad
 \deg\widehat P,\deg\widehat Q\le2m^2,\qquad
 \deg a_i\le4m^3,
\]
where $t$ is the fixed rational input threshold.
Together,~\eqref{I:U:eq:fixed-stability} and~\eqref{I:U:eq:fixed-coeff}
give $O(k^2+m)$ polynomial equalities and inequalities of degree
$O(m^3)$ in $p$ variables.

Their coefficients have bit lengths polynomial in $m$ and the binary
input length.  To see this, clear the input denominators and use the
determinant expansion: each monomial coefficient is a sum of products
of polynomially many input coefficients, and the logarithm of the
number of summands is polynomial in $m$.  The same bound applies to
adjugate entries, their products, and the final characteristic
polynomial.  This argument bounds bit length; it does not call for
enumerating the determinant expansion.

For completeness, the polynomials themselves can be constructed in
polynomial bit time when $p$ is fixed.  Their degree bounds imply only
$m^{O(p)}$ possible monomials.  Evaluate determinants, adjugates, and
characteristic polynomials over the rationals on an integer
interpolation grid of polynomial side length, and perform exact
multivariate interpolation.  The grid contains polynomially many
points for fixed $p$, and every evaluation has polynomial bit
complexity.  Adjugates are evaluated directly by cofactors or a
division-free matrix algorithm, so singular interpolation points cause
no difficulty.  Their intermediate and output bit lengths are
polynomial by the same determinant bounds.  The interpolation step
also has polynomial bit complexity.

\paragraph{Decision.}
Real algebraic feasibility with a fixed number of variables is
decidable in time polynomial in the number of polynomials, their
degrees, and their coefficient bit lengths
\cite{renegar1992,basu2006}.  Apply this result to the formula above.
The source dimension enters through polynomial degree and coefficient
size, while the number of variables is the fixed number $p$.
This proves the claimed bit-complexity bound.
\end{proof}


\begin{corollary}[Additive synthesis on a fixed slice]\label{I:U:cor:fixed-synthesis}
For fixed $d,k$, a rational stable source and a positive rational
$\varepsilon$ admit an algorithm that returns real algebraic
coefficients for a stable realization $g$ of dimension at most $k$
such that
\[
 \delta_k(H_f)\le\Unorm{H_f-H_g}
 \le\delta_k(H_f)+\varepsilon.
\]
Its bit complexity is polynomial in the binary encodings of the
source and $\varepsilon$, with an exponent depending on $d,k$.
\end{corollary}
\begin{proof}
Compute the rational source Gramians $P_f,Q_f$ by their Stein
equations.  The rational number
$F=1+\operatorname{tr}(P_fQ_f)$ has polynomial encoding length and
satisfies
$\Unorm{H_f}\le\Unorm{H_f}_{\mathrm{HS}}
=\sqrt{\operatorname{tr}(P_fQ_f)}\le F$.
The zero competitor is feasible at threshold $F$.
Bisect $[0,F]$ using Theorem~\ref{I:U:thm:fixed} until the interval width
is at most $\varepsilon$, retaining a feasible upper endpoint $u$.
The number of decisions is
$O(1+\log_+(F/\varepsilon))$, and all intermediate thresholds have
polynomial encoding length.
Standard real algebraic sample-point extraction
\cite{basu2006} applied to the feasible formula at $u$ returns
algebraic $T_a,b,c,X$ in polynomial time for fixed $p$.
They define a stable candidate with error at most
$u\le\delta_k(H_f)+\varepsilon$.
\end{proof}


\begin{remark}[Output model and scope]
The corollary returns exact real algebraic data, for example through
isolating intervals and algebraic sample-point encodings.  It does not
assert that independently rounding the coefficients to floating point
preserves stability or the error certificate.  Nor does it make a
practical running-time claim for the resulting algebraic algorithms.
For Gaussian-rational sources and complex competitors, the same
fixed-slice argument applies after writing real and imaginary parts
as separate variables and expressing the two Stein equations in real
coordinates on Hermitian matrices.  Their denominators can be
combined and squared in the norm test.  The number of candidate
variables still depends only on $d,k$, and all degree and bit-size
bounds remain polynomial in the source input length.
\end{remark}

\section{Sharper bounds for the binary obstruction family}\label{I:M:rev13:sec:sep}

Theorem~\ref{I:U:thm:separation} gives, for the source
$f_{r,k}(w)=q\,r^j$ on $w=a^jb^k$ ($q=1-r^2$),
\begin{equation}\label{I:M:rev13:eq:old-bound}
\delta_k(H_{f_{r,k}})\ge\max\Bigl\{\sqrt q,\ \sqrt{\tfrac{1+q^2(C_k-1)}{C_k}}\Bigr\},\qquad
C_k=\tbinom{2k}k,
\end{equation}
so that $\liminf_{r\uparrow1}\delta_k\ge C_k^{-1/2}$. We prove a bound that is uniform in
$r$ and much larger, and for $k=1$ we identify the limit. Throughout,
$z_j=\sqrt q\sum_{n\ge0}r^ne_{a^nb^j}$ ($0\le j\le k$) are the orthonormal vectors of the proof of Theorem~\ref{I:U:thm:separation} and
\[
H_f=z_0z_k^{*}+\sqrt q\sum_{j=0}^{k-1}z_{k-j}e_{b^j}^{*},\qquad
H_f^{*}z_0=z_k ,
\]
the last identity because $z_{k-j}\perp z_0$ for $j<k$.

\subsection{A model-space constant}
For a monic polynomial $\chi$ of degree $k$ with all roots in the open unit disk, let
$V_\chi\subset\ell^2(\MN)$ be the $k$-dimensional space of $\chi$-recurrent square-summable
sequences ($\sum_i c_iu_{j+i}=0$ for all $j\ge0$, $\chi=\sum_ic_iz^i$). Define
\begin{equation}\label{I:M:rev13:eq:kappa}
\kappa_k=\sup_\chi\ \sup_{0\ne u\in V_\chi}\frac{|u_k|^2}{\sum_{j\ge0}|u_j|^2}
=\sup_\chi\Mnorm{P_{V_\chi}e_k}^2 .
\end{equation}
\begin{lemma}[Blaschke form of $\kappa_k$]\label{I:M:rev13:lem:kappa}
Under $\ell^2(\MN)\cong H^2$, $V_\chi$ is the model space $K_\Theta=H^2\ominus\Theta H^2$ of
the degree-$k$ Blaschke product $\Theta$ with zeros at the conjugates of the roots of
$\chi$, and
\[
\Mnorm{P_{K_\Theta}z^k}^2=1-\sum_{m=0}^{k}|\theta_m|^2,\qquad
\Theta(z)=\sum_{m\ge0}\theta_mz^m .
\]
Hence $\kappa_k=1-\inf_\Theta\sum_{m\le k}|\theta_m|^2$ over degree-$k$ Blaschke products.
Moreover $\kappa_k\le1-1/C_k$, $\kappa_1=1/4$, and $\kappa_2=7/16$.
\end{lemma}
\begin{proof}
A $\chi$-recurrent sequence is $u_j=\sum_i p_i(j)\lambda_i^j$ over the roots $\lambda_i$
with polynomial coefficients; under $u\mapsto\sum_ju_jz^j$ the sequence $(\lambda^j)_j$
becomes $1/(1-\lambda z)=k_{\bar\lambda}$, the reproducing kernel at $\bar\lambda$, and
derivatives for multiple roots; these kernels span $K_\Theta$. Since
$P_{K_\Theta}f=f-\Theta P_+(\bar\Theta f)$ and multiplication by the inner function $\Theta$
is isometric, $\Mnorm{P_{K_\Theta}z^k}^2=1-\Mnorm{P_+(\bar\Theta z^k)}^2$, and
$\langle\bar\Theta z^k,z^j\rangle=\langle z^k,\Theta z^j\rangle=\overline{\theta_{k-j}}$
for $0\le j\le k$, zero for $j>k$. For the crude bound, the coefficients of $\chi$ satisfy
$|c_i|\le\binom ki$, so $|u_k|^2\le(\sum_{i<k}|c_i|^2)\sum_{j<k}|u_j|^2\le(C_k-1)
(\Mnorm u^2-|u_k|^2)$. For $k=1$, $\Theta=B_\lambda$ gives $\theta_0=-\lambda$,
$\theta_1=1-|\lambda|^2$ and $1-|\theta_0|^2-|\theta_1|^2=|\lambda|^2(1-|\lambda|^2)\le1/4$,
attained at $|\lambda|^2=1/2$. For $k=2$ apply two steps of the Schur algorithm to $\Theta$.
Writing $\alpha=\Theta(0)$, there are $\beta\in\mathbb D$ and $|\gamma|=1$
such that
\[
 \Theta(z)=\frac{\alpha+zF(z)}{1+\bar\alpha zF(z)},\qquad
 F(z)=\frac{\beta+z\gamma}{1+\bar\beta z\gamma}.
\]
Indeed $(\Theta-\alpha)/(1-\bar\alpha\Theta)$ is a degree-two inner
rational function vanishing at zero; dividing by $z$ leaves degree one.
Repeating this step leaves a unimodular constant. This is the elementary
finite Schur algorithm; see \citet{schur_nevanlinna}.
The first coefficients are
\[
 \theta_0=\alpha,\quad \theta_1=(1-|\alpha|^2)\beta,\quad
 \theta_2=(1-|\alpha|^2)\bigl((1-|\beta|^2)\gamma-\bar\alpha\beta^2\bigr).
\]
Set $a=|\alpha|$ and $t=|\beta|^2$. By the reverse triangle inequality,
\begin{align*}
 \sum_{j=0}^2|\theta_j|^2
 &\ge a^2+(1-a^2)^2\bigl[t+(1-(1+a)t)^2\bigr]\\
 &=\frac9{16}+(a-\tfrac12)^2(a+\tfrac34)\\
 &\quad +(1-a^2)^2(1+a)^2
 \left(t-\frac{1+2a}{2(1+a)^2}\right)^2
 \ \ge\ \frac9{16}.
\end{align*}
Equality holds for $\alpha=1/2$, $\beta=2/3$, $\gamma=1$, giving
\[
 \Theta(z)=\frac{z^2+z+1/2}{1+z+z^2/2},\qquad
 (\theta_0,\theta_1,\theta_2)=(1/2,1/2,1/4).
\]
Its zeros $(-1\pm i)/2$ lie in $\mathbb D$, so this is an admissible
degree-two Blaschke product and $\kappa_2=7/16$. No global numerical
minimisation is needed.
\end{proof}


\subsection{A uniform lower bound}
\begin{theorem}[Uniform obstruction constant]\label{I:M:rev13:thm:uniform}
For every $k\ge1$, every $0<r<1$ and either scalar field,
\[
\delta_k(H_{f_{r,k}})\ \ge\ \sqrt{1-\kappa_k}=E_k
\quad\bigl(\ge C_k^{-1/2}\bigr).
\]
Here $E_k$ is the complex scalar distance of
Lemma~\ref{I:B:lem:one-letter-projection}; that lemma gives
$\kappa_k=\Gamma_k^2=1-E_k^2$ without the outer-extremal input.
In particular $\inf_{0<r<1}\delta_1\ge\sqrt3/2$ and $\inf_{0<r<1}\delta_2\ge3/4$, while
$\sigma_{k+1}(H_{f_{r,k}})=\sqrt q$ ranges over $(0,1)$.
\end{theorem}
\begin{proof}
Let $G$ be a bounded word-Hankel operator of rank at most $k$ and $A=H_f-G$. Then
$\Mnorm A=\Mnorm{A^*}\ge\Mnorm{A^*z_0}=\Mnorm{z_k-G^*z_0}\ge\Mdist(z_k,\Mrange G^*)$.
By Lemma~\ref{I:U:lem:compact}, $G=OC$ with $O$ an
isometry of $\MFF^s$ ($s=\Mrank G\le k$), $C=O^*G$, $Ce_v=T_vb$ and $\Mnorm{T_a}\le1$; hence
$\overline{\Mrange G^*}=\overline{\Mrange C^*}=\{h_\eta:\eta\in\MFF^s\}$ with $h_\eta(v)=\eta^*T_vb$. Complex conjugation preserves norms and fixes the real vector $z_k$,
so it does not affect the desired projection estimate. Fix $\eta$ with
$h_\eta\ne0$. For each $n\ge0$ the sequence $u^{(n)}_j=\eta^*T_a^nT_b^jb$, $j\ge0$,
satisfies the recurrence of the characteristic polynomial of $T_b$ (Cayley--Hamilton),
whose roots lie in the closed unit disk, and is square-summable because
$\sum_j|u^{(n)}_j|^2\le\Mnorm{h_\eta}^2<\infty$. A square-summable sequence recurrent for a
polynomial with unimodular roots has no component on those roots, so it is recurrent for a
degree-$k$ polynomial with roots in the open disk (replace the unimodular roots by $0$).
Padding with zero states if $s<k$ changes nothing. Therefore
$|u^{(n)}_k|^2\le\kappa_k\sum_j|u^{(n)}_j|^2$ for every $n$, and
\[
|\langle z_k,h_\eta\rangle|^2\le\sum_n|h_\eta(a^nb^k)|^2
=\sum_n|u^{(n)}_k|^2\le\kappa_k\sum_n\sum_j|h_\eta(a^nb^j)|^2\le\kappa_k\Mnorm{h_\eta}^2,
\]
using Cauchy--Schwarz on the support of the unit vector $z_k$ and the fact that the words
$a^nb^j$ are distinct. Thus $\Mnorm{P_{\Mrange G^*}z_k}^2\le\kappa_k$ and
$\Mdist(z_k,\Mrange G^*)^2\ge1-\kappa_k$. The crude bound of Lemma~\ref{I:M:rev13:lem:kappa} gives the
parenthetical comparison.
\end{proof}


\begin{remark}
The proof of Theorem~\ref{I:U:thm:separation} bounds the Rayleigh quotient of $H_f$ on one kernel vector of $G$;
Theorem~\ref{I:M:rev13:thm:uniform} instead bounds how much of the top right singular vector $z_k$
can lie in the $k$-dimensional backward-shift-invariant row space of $G$, using only the
words $a^nb^j$. The resulting constant does not depend on $r$, so the obstruction holds at
every $r$. The diverging ratio to $\sigma_{k+1}=\sqrt q$ occurs as $r\uparrow1$, when the source stability margin closes; a margin-dependent positive bound is not excluded.\end{remark}

\subsection{The rank-one case: the obstruction is asymptotically total}
\begin{lemma}[Rank-one word-Hankel operators]\label{I:M:rev13:lem:rank-one}\label{I:S:lem:rankone}
Over $\MFF\in\{\MR,\mathbb C\}$, every nonzero bounded word-Hankel operator of rank one is
$G=\gamma\,mm^{\top}$ (transpose, also over $\mathbb C$), where
$m_w=\prod_{i}t_{w_i}$ for scalars $t_a\in\MFF$ with $\sum_a|t_a|^2<1$ and
$\gamma\in\MFF\setminus\{0\}$; conversely each such $G$ is a bounded rank-one word-Hankel.
\end{lemma}
\begin{proof}
Write $G=xy^*$ and $g(uv)=G[u,v]=x_u\bar y_v$. Then $g(u)=x_u\bar y_\varepsilon$,
$g(v)=x_\varepsilon\bar y_v$ and $g(\varepsilon)=x_\varepsilon\bar y_\varepsilon$. If
$g(\varepsilon)=0$ then $x_\varepsilon=0$ or $y_\varepsilon=0$, forcing $g\equiv0$.
Otherwise $g(uv)=g(u)g(v)/g(\varepsilon)$, so $w\mapsto g(w)/g(\varepsilon)$ is
multiplicative: $g(w)=\gamma\prod_it_{w_i}$ with $\gamma=g(\varepsilon)$,
$t_a=g(a)/g(\varepsilon)$, and $G[u,v]=\gamma m_um_v$. Finally
$\Mnorm G=|\gamma|\Mnorm m^2$ and $\Mnorm m^2=\sum_n(\sum_a|t_a|^2)^n$.
\end{proof}


\begin{theorem}[Sharpened rank-one obstruction]\label{I:M:rev13:thm:rank-one-sharp}
For $k=1$, $0<r<1$, $q=1-r^2$, over either scalar field,
\begin{equation}\label{I:M:rev13:eq:rank-one-sharp}
\delta_1(H_{f_{r,1}})\ \ge\ \frac{\sqrt{1+2\sqrt q}}{1+\sqrt q}\ =\ 1-\frac q{2}+O(q^{3/2}).
\end{equation}
Consequently $\lim_{r\uparrow1}\delta_1(H_{f_{r,1}})=1=\Mnorm{H_{f_{r,1}}}$ while
$\sigma_2(H_{f_{r,1}})=\sqrt q\to0$: as the stability margin closes, no rank-one
word-Hankel operator approximates the source better than the zero operator, up to
$O(q)$. The right-hand side of \eqref{I:M:rev13:eq:rank-one-sharp} is at least $\sqrt3/2$ for all
$q\in(0,1]$, so it dominates Theorem~\ref{I:M:rev13:thm:uniform} at $k=1$, and it exceeds
\eqref{I:M:rev13:eq:old-bound} for $q<0.76$ approximately.
\end{theorem}
\begin{proof}
Let $G=\gamma mm^{\top}$ be any bounded rank-one competitor (Lemma~\ref{I:M:rev13:lem:rank-one};
$G=0$ gives error $\Mnorm{H_f}=1$) and $A=H_f-G$, with
$H_f=z_0z_1^{\top}+\sqrt q\,z_1e_\varepsilon^{\top}$. Put
\[
p=m^{\top}z_0=\sqrt q\sum_nr^nt_a^n=\frac{\sqrt q}{1-rt_a},\qquad
s=m^{\top}z_1=\sqrt q\sum_nr^nt_a^nt_b=t_bp .
\]
Then $Az_1=z_0-\gamma s\,m$ and, since $H_f^{\top}z_0=z_1$ and $G^{\top}=G$,
$A^{\top}z_0=z_1-\gamma p\,m$. With $c=\gamma sp$,
\[
\Mnorm{Az_1}^2=1-2\operatorname{Re}c+|c|^2\frac{\Mnorm m^2}{|p|^2},\qquad
\Mnorm{A^{\top}z_0}^2=1-2\operatorname{Re}c+|c|^2\frac{\Mnorm m^2}{|s|^2}.
\]
Because $|s|=|t_b||p|\le|p|$, the second expression dominates, so
$\Mnorm A^2\ge1-2|c|+|c|^2\Mnorm m^2/|s|^2\ge1-|s|^2/\Mnorm m^2$ after minimising over
$|c|\ge0$ (if $t_b=0$ then $s=0$ and $\Mnorm{Az_1}=\Mnorm{z_0}=1$). Now
\[
\frac{|s|^2}{\Mnorm m^2}=\frac{q\,|t_b|^2(1-|t_a|^2-|t_b|^2)}{|1-rt_a|^2}
\le\frac{q\,\beta(1-\alpha^2-\beta)}{(1-r\alpha)^2}\le\frac q4\,\frac{(1-\alpha^2)^2}{(1-r\alpha)^2},
\]
with $\alpha=|t_a|$, $\beta=|t_b|^2\le1-\alpha^2$. The function
$\varphi(\alpha)=(1-\alpha^2)/(1-r\alpha)$ on $[0,1)$ has derivative proportional to
$r\alpha^2-2\alpha+r$ and attains its maximum $2/(1+\sqrt q)$ at
$\alpha=(1-\sqrt q)/r$ (where $1-\alpha^2=2\sqrt q/(1+\sqrt q)$ and $1-r\alpha=\sqrt q$).
Hence $|s|^2/\Mnorm m^2\le q/(1+\sqrt q)^2$ and
$\delta_1^2\ge1-q/(1+\sqrt q)^2=(1+2\sqrt q)/(1+\sqrt q)^2$. The expansion and the limit
are immediate, and $\delta_1\le\Mnorm{H_f}=1$ because $G=0$ is admissible.
\end{proof}




\begin{remark}The rank-two boundary limit is proved in Theorem~\ref{I:B:thm:main}, and its first correction is proved in Theorem~\ref{I:R:thm:main-general}.\end{remark}

\section{Exact behavioral realization and the marked algebra}\label{I:S:sec:exact}
Fix a finite nonempty alphabet $\Sigma$, $d=|\Sigma|$, and the free monoid $\Swords$ with empty word $\Seps$. A linear weighted finite automaton (WFA) over a field $\SK$ is
\[
 f(w)=\alpha^\top A_w\beta,\qquad A_{a_1\cdots a_m}=A_{a_1}\cdots A_{a_m},\quad A_\Seps=I_n.
\]
We use row states $s(w)=\alpha^\top A_w$. Here \emph{row-span minimality} means that the $s(w)$ span $\SK^{1\times n}$; \emph{column-span minimality} means that $\{A_w\beta\}$ spans $\SK^n$. A realization is minimal when both hold. In complex analytic arguments $\alpha^\top$ remains the algebraic output row; Hermitian adjoints in Gramians are written separately.

\begin{theorem}[Hankel realization and exact finite witnesses]\label{I:S:thm:realization}
The algebraic Hankel matrix $H_f[u,v]=f(uv)$ has finite rank $n$ if and only if $f$ has a finite linear realization. Its rank is the minimum linear state dimension. Two realizations of dimensions $n_1,n_2$ agree on all words if they agree on words of length at most $n_1+n_2-1$.
\end{theorem}
\begin{proof}
A realization gives $H_f=OC$, with $O[u,:]=\alpha^\top A_u$ and $C[:,v]=A_v\beta$, so its rank is at most $n$. Conversely take the span of residual rows $r_u(v)=f(uv)$. The shift $r_u\mapsto r_{ua}$ is well-defined: if $\sum c_u r_u=0$, evaluation at $av$ gives $\sum c_u f(uav)=0$ for all $v$. The initial residual is $r_\Seps$ and the output is evaluation at $\Seps$. A basis supplies a realization of dimension $\Srank H_f$.

For a realization of dimension $N$, let $V_j=\SSpan\{\alpha^\top A_w:|w|\le j\}$. Once $V_j=V_{j+1}$, it is invariant and contains all reachable rows. If the initial row is nonzero, $\dim V_0=1$, so stabilization occurs by $j=N-1$; the zero case is immediate. Vanishing output on words of length at most $N-1$ therefore implies vanishing on all words. Apply this to the direct-sum difference realization of dimension $n_1+n_2$.
\end{proof}


An explicit minimization first restricts to the generated row space and then quotients by the rows annihilating every future column. Exact basis closure terminates in finite dimension. The same argument applies to finitely many output coordinates: the future tests include both a word and an output coordinate.

\begin{proposition}[Discrete and affine memory have different costs]\label{I:S:prop:types}
The minimum number of deterministic states realizing an observable $f$ is the number of residual classes for $u\sim v$ iff $f(ut)=f(vt)$ for every $t$. It can be infinite while $\Srank H_f$ is finite. If a vector observable has residual affine-hull dimension $r_{\rm aff}$, the minimum dimension of an affine-state realization with affine readout is $r_{\rm aff}$, and
\[
 \Srank H_f-1\le r_{\rm aff}\le\Srank H_f.
\]
\end{proposition}
\begin{proof}
Equal deterministic states must give equal future outputs; the residual classes themselves form an attaining deterministic automaton because the relation is a right congruence. For $f(a^m)=m$ in characteristic zero, residuals $t\mapsto m+|t|$ are all distinct but their linear span has dimension two. For the affine statement, every future observable is affine in an affine state, so the residual affine dimension is a lower bound. Shift maps preserve affine relations by the same evaluation argument as in Theorem~\ref{I:S:thm:realization}; their restriction to the residual affine hull attains the bound. An affine hull of dimension $r$ has linear span of dimension $r$ or $r+1$, proving the inequalities.
\end{proof}


Define the syntactic ideal and the generated algebra by
\[
 I_f=\{p\in\SK\langle\Sigma\rangle:f(upv)=0\ \forall u,v\},\quad
 \ScA_f=\SK\langle\Sigma\rangle/I_f,\quad B=\Salg(I,(A_a)_a).
\]
The letters, unit, and tested observable are retained as markings. For a finite-dimensional algebra, let $J(B)$ denote its Jacobson radical and $\SLL(B)=\min\{j\ge1:J(B)^j=0\}$, with $\SLL(0)=0$.

\begin{theorem}[Universal marked quotient and Loewy dominance]\label{I:S:thm:marked}
Every finite realization induces a letter-marked surjection $B\twoheadrightarrow\ScA_f$. It is an isomorphism for a minimal realization. For any finite-dimensional algebra epimorphism $\pi:B\twoheadrightarrow A$,
\[
 \pi(J(B)^j)=J(A)^j\quad(j\ge0).
\]
Thus dimensions of the algebra and of every graded radical layer, and Loewy length, cannot increase under the quotient.
\end{theorem}
\begin{proof}
The representation kernel is contained in $I_f$: a zero matrix polynomial remains zero between every prefix row and future column. This yields the quotient. Under both span conditions, $\alpha^\top A_u p(A)A_v\beta=0$ for all $u,v$ implies $p(A)=0$, giving kernel equality. The image of $J(B)$ is a nilpotent ideal, hence contained in $J(A)$. Also $A/\pi(J(B))$ is a quotient of the semisimple algebra $B/J(B)$ and is semisimple. Its radical vanishes, giving the reverse inclusion. Taking products proves the power identities; passage to successive quotients gives the graded surjections.
\end{proof}


\begin{remark}[The charged algebra matters]
This theorem concerns a quotient of the algebra actually generated by the transitions. It is false with an arbitrary containing algebra in its place: $\SK[t]/(t^2)$ embeds in $M_2(\SK)$, which is semisimple. Loewy length is also not neural depth. A single recurrent layer with an $(m+1)$-state nilpotent shift realizes a length-$m$ impulse. For a unary matrix over $\SC$, the Chinese remainder decomposition of $\SC[t]/(\mu_A)$ shows that the Loewy length of its generated algebra is the largest Jordan-block size, not the number of network layers.
\end{remark}

\section{Bounded Hankels, strict stability, and exact norms}\label{I:S:sec:analytic}
Now $\ScH=\ell^2(\Swords)$ has its unweighted orthonormal word basis. The analytic theorems are stated over $\SR$; the Hermitian versions hold over $\SC$. Define
\[
 \Phi_A(X)=\sum_a A_a^\top X A_a,\qquad
 \Psi_A(X)=\sum_a A_a X A_a^\top
\]
on real symmetric matrices. They are adjoint for the trace pairing. The observability and controllability feature maps, in their operator-factorization convention, are
\[
 (Ox)(u)=\alpha^\top A_u x,\qquad Cz=\sum_v A_v\beta z(v).
\]

\begin{theorem}[Minimal boundedness is strict all-word stability]\label{I:S:thm:stability}
For a minimal finite real realization the following are equivalent:
\[
 H_f\text{ bounded};\quad f\in\ell^2(\Swords);\quad
 \rho(\Phi_A)<1;\quad
 \exists X\succ0:\ X-\Phi_A(X)\succ0.
\]
The Gramians are finite and positive definite and satisfy
\begin{align}
 P&=\sum_w A_w^\top\alpha\alpha^\top A_w
   =\alpha\alpha^\top+\Phi_A(P),\label{I:S:eq:P}\\
 Q&=\sum_w A_w\beta\beta^\top A_w^\top
   =\beta\beta^\top+\Psi_A(Q).\label{I:S:eq:Q}
\end{align}
Without minimality, strict stability remains sufficient but is not necessary for boundedness of the realized Hankel.
\end{theorem}
\begin{proof}
Boundedness implies $f=H_f e_\Seps\in\ell^2$. If $f\in\ell^2$, all shifted sequences $u\mapsto f(uv)$ and $v\mapsto f(uv)$, with the other word fixed, are square summable, because concatenation with a fixed word is injective. Choose a basis among $A_v\beta$; its images under $O$ are such sequences. Hence $O$ is bounded. Choose a basis among $A_u^\top\alpha$ to prove boundedness of $C^*$. Their Gramians are precisely \eqref{I:S:eq:P}--\eqref{I:S:eq:Q}; the span conditions make them positive definite.

Unique decomposition by the last letter gives the first Stein equation. Its tail identity is
\[
 \Phi_A^m(P)=\sum_{|v|=m}A_v^\top P A_v
           =\sum_{|w|\ge m}A_w^\top\alpha\alpha^\top A_w\longrightarrow0.
\]
In the order-unit norm $\Snorm{Y}_P=\Snorm{P^{-1/2}YP^{-1/2}}$, positivity implies
$\Snorm{\Phi_A^m}_{P\to P}\le\Snorm{\Phi_A^m(P)}_P$. Some power has norm less than one, so $\rho(\Phi_A)<1$.

Conversely, a positive map with spectral radius below one has $X=\sum_{j\ge0}\Phi_A^j(I)\succ0$ and $X-\Phi_A(X)=I$. If a strict Lyapunov inequality holds, choose $\mu>0$ with $\Phi_A(X)\preceq(1-\mu)X$. The same order-unit argument gives spectral radius below one. Since $\Psi_A$ is adjoint to $\Phi_A$, both Neumann sums in \eqref{I:S:eq:P}--\eqref{I:S:eq:Q} converge. They define bounded $O,C$, and $H_f=OC$ first on finitely supported vectors and then by continuity. This converse does not use minimality.
\end{proof}


\begin{corollary}[Complete contractive cover]
Every bounded finite-rank real word Hankel has a minimal realization similar to one satisfying $\sum_a A_a^\top A_a\prec I$. Any smaller rank can be embedded in a prescribed larger stable dimension by zero-padding.
\end{corollary}
\begin{proof}
Apply the realization theorem and Theorem~\ref{I:S:thm:stability}; conjugate $A_a$ to $X^{1/2}A_aX^{-1/2}$. Zero-padding adds a strictly positive identity slack on the unused coordinates.
\end{proof}


\begin{theorem}[Finite norm formula, including singular Gramians]\label{I:S:thm:norm}
For any stable, possibly nonminimal realization, $H_f=OC$ and
\[
 \Snorm{H_f}^2=\lambda_{\max}(P^{1/2}QP^{1/2})
              =\lambda_{\max}(L^\top QL),\qquad P=LL^\top.
\]
Its nonzero squared singular values are the nonzero eigenvalues of these finite matrices. For $t\ge0$,
\begin{equation}\label{I:S:eq:factorfree}
 \Snorm{H_f}\le t\quad\Longleftrightarrow\quad t^2P-PQP\succeq0.
\end{equation}
\end{theorem}
\begin{proof}
$H_f^*H_f=C^*PC$; the nonzero spectra of $T^*T$ and $TT^*$ agree for $T=P^{1/2}C$, proving the formula. The factor $L$ gives the same nonzero spectrum, as does any square factor of $P$. Put $S=P^{1/2}QP^{1/2}$. It vanishes on $\ker P$ and preserves $\Sran P$. Congruence by the invertible restriction of $P^{1/2}$ on that range identifies $S\preceq t^2I$ with $t^2P-PQP\succeq0$ there. Both latter terms vanish on the kernel, and its additional zero eigenvalues are harmless because $t\ge0$.
\end{proof}


For two models use $D_a=\Sdiag(A_a,B_a)$, $c=(\alpha,-\gamma)$, $b=(\beta,\nu)$ as columns. Their difference is $c^\top D_wb$. Its Gramians certify the error on \emph{all} words; a finite word section alone normally supplies only a lower bound.

\begin{proposition}[Elementary norm and tail bounds]\label{I:S:prop:normbounds}
For any series $h$,
\[
 \Snorm{H_h}_{\mathrm{HS}}^2=\sum_w(|w|+1)|h(w)|^2
\]
whenever either side is finite. If $h\in\ell^1$, then $\Snorm{H_h}\le\Snorm h_1$. If $h=\sum_m h_m$ with $h_m$ homogeneous and $\sum_m\Snorm{h_m}_2<\infty$, then $\Snorm{H_h}\le\sum_m\Snorm{h_m}_2$.
\end{proposition}
\begin{proof}
Every word of length $m$ has exactly $m+1$ prefix--suffix splits, giving the Hilbert--Schmidt identity by nonnegative summation. Each absolute row and column sum is at most $\Snorm h_1$; weighted Cauchy--Schwarz proves the Schur bound. A homogeneous symbol has norm $\Snorm{h_m}_2$ by the orthogonal block argument in Lemma~\ref{I:S:lem:homogeneous}; summing operators proves the last assertion and identifies all entries of the limit.
\end{proof}


\section{Existence, exact algebraicity, and certificate logic}\label{I:S:sec:global}
For a real target, $\mathbb F\in\{\SR,\SC\}$, and $k\ge0$, define
\[
 \mathscr H_k^{\mathbb F}=\{H_g:H_g\text{ bounded},\ g:\Swords\to\mathbb F,\ \Srank H_g\le k\},
 \qquad \SdeltaF{k}(H)=\inf_{G\in\mathscr H_k^{\mathbb F}}\Snorm{H-G}.
\]
Ranks are linear WFA dimensions. Competitors are not required to be commuting, nonnegative, or probabilistically normalized. For compact $H$, the unstructured lower bound is $\sigma_{k+1}(H)\le\SdeltaF{k}(H)$. Real targets satisfy $\delta_k^\SC\le\delta_k^\SR$.

Attainment over either scalar field is Theorem~\ref{I:U:thm:attainment};
its proof applies to any bounded target operator. The algebraic endpoint
and degree-three existential formulation are Corollary~\ref{I:U:cor:algebraic}.

The same conclusion holds for a bounded real-algebraic WFA after minimal
stable realization. Strict thresholds can be encoded by requiring the
factor of the norm slack to be invertible. No practical running-time
bound follows from general quantifier elimination.

\begin{remark}[Strict versus non-strict certificates]
For singular $P$, strict positive definiteness of $t^2P-PQP$ is impossible even if the error is strictly below $t$. The $t^2I-L^\top QL$ encoding handles strict thresholds correctly. Alternatively, proving error at most a smaller rational threshold suffices for a strict public upper endpoint. A feasible candidate is an upper certificate, not a certificate of its optimality; a solver timeout is not an infeasibility proof.
\end{remark}

\begin{proposition}[Perturbation transfer]\label{I:S:prop:transfer}
For one fixed competitor class, $|\SdeltaF{k}(H)-\SdeltaF{k}(\widetilde H)|\le\Snorm{H-\widetilde H}$. If $L\le\SdeltaF{k}(H_f)\le\Snorm{H_f-H_g}\le U$ and $\Snorm{H_T-H_f}\le\eta$, then
\[
 \max(0,L-\eta)\le\SdeltaF{k}(H_T)\le\Snorm{H_T-H_g}\le U+\eta.
\]
For compact operators, the spectral gap $\SdeltaF{k}-\sigma_{k+1}$ decreases by at most $2\eta$ under a perturbation of size $\eta$.
\end{proposition}
\begin{proof}
Apply the triangle inequality against each competitor and take infima in both directions. Individual singular values are 1-Lipschitz, proving the last assertion.
\end{proof}


\section{Rank-one approximation and the free-word obstruction}\label{I:S:sec:rankone}
Let $L_a e_w=e_{aw}$ and $R_a e_w=e_{wa}$. Then $H_f L_a=R_a^*H_f$. Consequently the range of a finite-rank Hankel is invariant under every $R_a^*$, and its adjoint range is invariant under every $L_a^*$. Restricting the usual subspace-width minimization to these invariant subspaces gives lower bounds on $\SdeltaF{k}$. This restriction remembers the letter action which a free matrix SVD discards.

By Lemma~\ref{I:M:rev13:lem:rank-one}, write $G=c k_zk_z^\top$,
where $k_z(w)=\prod_i z_{w_i}$, $\sum_a|z_a|^2<1$, and
$\|k_z\|^2=(1-\|z\|_2^2)^{-1}$. The coefficient transpose is
intentional also for complex parameters.

\begin{lemma}[One-corner completion]\label{I:S:lem:corner}
For a bounded operator $T$ and unit output/input vectors $u,v$,
\[
 \min_c\Snorm{T-cuv^*}
   =\max\{\Snorm{T(I-vv^*)},\Snorm{(I-uu^*)T}\}.
\]
The minimizing scalar can be real for real data.
\end{lemma}
\begin{proof}
Use the decompositions along $u,v$. Only the upper-left scalar block is free; the bottom row and right column give the lower bound $\gamma$. If $\gamma=0$, cancel the only nonzero block. Otherwise divide by $\gamma$ and denote the three fixed blocks by $b,c,d$, so their row and column norms are at most one. The inequalities give contractions $B,C$ with
$b=B(I-d^*d)^{1/2}$ and $c=(I-dd^*)^{1/2}C$. These follow by defining the contraction on the range of the defect square root and extending continuously. The block matrix with free corner $-Bd^*C$ factors as
\[
 \begin{pmatrix}B&0\\0&I\end{pmatrix}
 \begin{pmatrix}-d^*&(I-d^*d)^{1/2}\\(I-dd^*)^{1/2}&d\end{pmatrix}
 \begin{pmatrix}C&0\\0&I\end{pmatrix}.
\]
The middle matrix is unitary by the defect-operator identities; the others are contractions. Rescaling proves attainment. This is the classical Parrott completion argument~\cite{parrott}.
\end{proof}


\begin{theorem}[An exact real rank-one variational formula]\label{I:S:thm:rankone}
For a bounded real Hankel and $P_b$ the orthogonal projection onto $\SR k_b$,
\[
 \delta_1^\SR(H_f)=\min_{b\in\SR^d,\ \Snorm b<1}
 \max\{\Snorm{H_f(I-P_b)},\Snorm{(I-P_b)H_f}\}.
\]
The minimum over the open ball is attained. If $H_f$ has finite rank with singular values $\sigma_j$ and right/left singular vectors $\xi_j,\eta_j$, each squared term is a finite rank-one downdate:
\[
 \lambda_{\max}\bigl(D^{1/2}(I-c^Rc^{R\top})D^{1/2}\bigr),\qquad
 \lambda_{\max}\bigl(D^{1/2}(I-c^Lc^{L\top})D^{1/2}\bigr),
\]
where $D=\Sdiag(\sigma_j^2)$, $c^R_j=\Sip{\xi_j}{k_b/\Snorm{k_b}}$ and $c^L_j=\Sip{\eta_j}{k_b/\Snorm{k_b}}$.
\end{theorem}
\begin{proof}
Fix $b$ and apply Lemma~\ref{I:S:lem:corner} with $u=v=k_b/\Snorm{k_b}$. The zero candidate cannot improve on this minimum. The objective is continuous on the ball and bounded above by $\Snorm{H_f}$. As $\Snorm b\to1$, the normalized kernels converge weakly to zero, uniformly on every fixed finite set of coordinates: each such coordinate is bounded in absolute value by $\sqrt{1-\Snorm b^2}$. Choose a finitely supported unit vector nearly attaining $\Snorm{H_f}$ to see that $\Snorm{H_f(I-P_b)}\to\Snorm{H_f}$, uniformly in boundary direction. Apply the same argument to $H_f^*$ for the other term. If the infimum equals $\Snorm{H_f}$, every point minimizes; otherwise it lies in a compact sub-ball and is attained. Finally apply the identity between the nonzero spectra of $TT^*$ and $T^*T$ to $(I-P_b)[\xi_j]D^{1/2}$, and similarly on the left.
\end{proof}


For a real minimal stable realization the kernel evaluations are rational:
\[
 Ck_b=(I-\sum_a b_aA_a)^{-1}\beta,\qquad
 O^*k_b=(I-\sum_a b_aA_a^\top)^{-1}\alpha.
\]
The inverses exist on the ball: in the contractive gauge,
$\Snorm{\sum b_aA_a}\le\Snorm b<1$. The transpose sum has the same norm. The power series converge absolutely by Cauchy--Schwarz against the finite Gramians.

\begin{corollary}[Rank-two secular equality criterion]\label{I:S:cor:secular}
Suppose $\Srank H_f=2$ and $\sigma_1>\sigma_2>0$. A one-sided term in Theorem~\ref{I:S:thm:rankone} equals $\sigma_2$ exactly when its corresponding normalized coefficients satisfy
\[
 c_2=0,\qquad c_1^2\ge1-\sigma_2^2/\sigma_1^2.
\]
Both the left and the right requirements must hold at the same $b$ for spectral equality. If $\sigma_1=\sigma_2=s$, every one-sided term equals $s$ and this zero test is inapplicable.
\end{corollary}
\begin{proof}
Subtract $\sigma_2^2I$ from the $2\times2$ downdate. Its determinant is
$-\sigma_2^2c_2^2(\sigma_1^2-\sigma_2^2)$. Negative semidefiniteness therefore forces $c_2=0$, after which its remaining diagonal inequality is the stated condition. For a repeated value, $s^2(I-cc^\top)$ has an eigenvalue $s^2$ and is bounded above by $s^2I$.
\end{proof}


\subsection{A commuting isospectral separation}
\begin{theorem}[Spectrum and abstract algebra do not determine the optimum]\label{I:S:thm:commutative}
Let $f=\mathbf1_{\Seps}+\mathbf1_b+\mathbf1_{aa}$ over $\{a,b\}$. It has nonnegative commuting minimal transitions $A_a=J$, $A_b=J^2$, where $J$ is the $3\times3$ upper shift. Its nonzero singular values are $(2,1,1)$, and
\[
 \sqrt{727/544}\le\delta_1^\SC(H_f)\le\delta_1^\SR(H_f)<29/25.
\]
There is a unary-supported nonnegative target with exactly the same singular values and abstract syntactic algebra $\mathbb F[t]/(t^3)$ whose optimal rank-one error is exactly one.
\end{theorem}
\begin{proof}
Take $\alpha=e_1$, $\beta=(1,0,1)^\top$. The only nonzero rows and columns are indexed by $(\Seps,a,b,aa)$, and the entire nonzero operator is
\[
 H=\begin{pmatrix}1&0&1&1\\0&1&0&0\\1&0&0&0\\1&0&0&0\end{pmatrix}.
\]
Its eigenvectors $(2,0,1,1)/\sqrt6$, $(0,1,0,0)$ and $(1,0,-1,-1)/\sqrt3$ have eigenvalues $2,1,-1$. Thus the rank is three and the realization is minimal.

For a complex rank-one candidate with parameters $(p,q)$, put $x=|p|^2$, $t=|q|$, $m=1-x-t^2>0$, and $z=q+p^2$. Projecting $H$ away from its possible output range gives the finite matrix
$M=H^2-mww^*$, $w=(1+z,p,1,1)^\top$. On $\Sran H$ one eigenvalue is one; the other two, $0\le\mu_-\le1\le\mu_+\le4$, satisfy
\[
 Q_0(\lambda)=(\lambda-4)(\lambda-1)+4A(\lambda-1)+B(\lambda-4),
\]
where $A=m|2+z|^2/6$ and $B=m(x+|1-z|^2/3)$. Direct evaluation gives
\[
 Q_0(4/3)=\tfrac23\{m(4\Re z-|z|^2-4x)-4/3\}.
\]
Since $t+x<1+t-t^2\le5/4$, the reverse triangle inequality yields
$4\Re z-|z|^2-4x\le4t-(t+x)^2$. If the right side is negative the next bound is immediate; otherwise both factors decrease with $x$. Therefore
$m(4\Re z-|z|^2-4x)\le(1-t^2)(4t-t^2)$.
For $y=t-1/2\in[-1/2,1/2]$, the exact identity
\[
 4-3(1-t^2)(4t-t^2)
 =\tfrac{51}{4}(y-1/17)^2+\tfrac5{272}
   +3y^2(y+1/2)(5/2-y)
\]
makes $Q_0(4/3)\le-5/1224$. Hence
$\mu_+\ge4/3+(5/1224)/(4/3)=727/544$, proving the lower bound. Subsection~\ref{I:S:app:commupper} gives a rational candidate and its strict upper certificate.

For the comparison put $r=2^{1/3}$ and let $h(\Seps)=2-r$, $h(a)=r^2-1$, $h(aa)=r$, with all other coefficients zero. Its unary block
\[
 K=\begin{pmatrix}2-r&r^2-1&r\\r^2-1&r&0\\r&0&0\end{pmatrix}
\]
has characteristic polynomial $(\lambda-2)(\lambda-1)(\lambda+1)$. Let $b_0=r^{-1}$ and $h_*(a^n)=(2-r)b_0^n$, zero on words containing $b$. The vector $(1,b_0,b_0^2)$ is a top eigenvector, $b_0^6=1/4$, and $\Snorm{H_{h_*}}=2$. Its normalized infinite kernel has squared overlap $3/4$ with the normalized top eigenvector of $K$. The difference of the two rank-one terms of size two has norm $2\sqrt{1-3/4}=1$; the orthogonal eigenvectors of eigenvalues $1,-1$ are unchanged. Thus $\Snorm{H_h-H_{h_*}}=1$, which is optimal by the singular-value bound. Both targets have generated algebra $\mathbb F[J]\simeq\mathbb F[t]/(t^3)$, but the letter $b$ is marked by $J^2$ in one and zero in the other.
\end{proof}


The identification of $b$ with two unary steps is not an isometry: $b$ and $aa$ are distinct words with different splits. More generally, if a symmetric series depends only on its count vector $\mathbf m$, define $c(\mathbf m)=|\mathbf m|!/\prod_i m_i!$. Normalized sums over words with a fixed count give an isometry under which its nonzero block is
\[
 \widehat H_f[\mathbf m,\mathbf n]
 =f(\mathbf m+\mathbf n)\sqrt{c(\mathbf m)c(\mathbf n)}.
\]
This follows by summing the identical entries over each fiber. Dropping the multinomial factors changes the norm. In a minimal rational realization, symmetry is equivalent to commuting transitions, since $f(uabv)=f(ubav)$ for all contexts forces $A_aA_b=A_bA_a$. At rank one all competitors are symmetric; at larger ranks, restricting competitors to symmetric series can increase the optimum.

\section{Coefficient and trace bounds; a sharp homogeneous constant}\label{I:S:sec:homogeneous}
For $p,q\ge0$, write $C_{p,q}(f)=(f(uv))_{|u|=p,|v|=q}$, and let $E_k(C)^2=\sum_{j>k}\sigma_j(C)^2$.

\begin{theorem}[Coefficient lower bound across lengths]\label{I:S:thm:coefficient}
For every bounded Hankel and $N\ge0$,
\[
 \delta_k^\SC(H_f)^2\ge
 \sum_{m=0}^N\max_{0\le p\le m}E_k(C_{p,m-p}(f))^2.
\]
If $H_f$ is compact, this combines with the independent lower bound $\sigma_{k+1}(H_f)$.
\end{theorem}
\begin{proof}
Every $C_{p,q}(g)$ of a rank-$k$ competitor has rank at most $k$. Frobenius Eckart--Young gives an error of at least $E_k(C_{p,q}(f))^2$. For fixed $m=p+q$, that squared error is exactly $\sum_{|w|=m}|f(w)-g(w)|^2$, independent of the split. Maximize over the splits, and sum over different lengths. The result is bounded above by $\Snorm{(H_f-H_g)e_\Seps}^2\le\Snorm{H_f-H_g}^2$. Take infima.
\end{proof}

One must maximize, not sum, over splits of the same length; otherwise the same coefficients are counted repeatedly without justification.

\begin{lemma}[Homogeneous norm and spectrum]\label{I:S:lem:homogeneous}
If $f\ne0$ is supported at length $m$, then $\Snorm{H_f}=a:=\Snorm f_2$. If $m\ge2$, then $\sigma_1=\sigma_2=a$ and $\sigma_3\ge a/\sqrt d$.
\end{lemma}
\begin{proof}
The word-length decomposition makes each nonzero block map degree $q$ to degree $m-q$, with block $C_{m-q,q}(f)$. Their input and output spaces are mutually orthogonal, so their singular-value lists combine. Each Frobenius norm is $a$. The two endpoint blocks have singular value $a$, and every other block has norm at most $a$. The block $C_{1,m-1}$ has at most $d$ rows and Frobenius norm $a$, hence has a singular value at least $a/\sqrt d$.
\end{proof}


\begin{theorem}[A trace certificate at one length]\label{I:S:thm:trace}
Put $A=\Snorm{H_f}$ and $a_m=(\sum_{|w|=m}|f(w)|^2)^{1/2}$. Then
\[
 \delta_k^\SC(H_f)\ge
 \max\left\{0,\frac{(m+1)a_m-kA}{m+1+k}\right\}.
\]
For a homogeneous degree-$m$ target this becomes
$\delta_k^\SC(H_f)\ge A\max\{0,(m+1-k)/(m+1+k)\}$.
\end{theorem}
\begin{proof}
If $a_m=0$ the claim is immediate. Let $T=H_{f\mathbf1_{|w|=m}}$, so $\Snorm T=a_m$. For a candidate $G=H_g$ of rank at most $k$ and error $t$, trace duality and $\Snorm G_1\le k\Snorm G\le k(A+t)$ give
\[
 (m+1)\left|\sum_{|w|=m}\overline{f(w)}g(w)\right|
 =|\Sip T G_{\rm HS}|\le a_m k(A+t).
\]
The empty column, paired with $f\mathbf1_{|w|=m}/a_m$, gives $a_m-t$ as a lower bound on the real part of that normalized coefficient inner product. Combining and rearranging proves the result. The factor $m+1$ comes from the exact number of splits and is justified here by trace duality.
\end{proof}


\begin{theorem}[Unbounded alphabet-dependent gap and sharp fixed-alphabet constant]\label{I:S:thm:sharp}
Let $f_{d,m}=\sum_{i=1}^d\mathbf1_{a_i^m}$ with $m\ge2$. It has a nonnegative commuting nilpotent realization of rank $d(m-1)+2$ and spectrum
\[
 (\sqrt d,\sqrt d,\underbrace{1,\ldots,1}_{d(m-1)}).
\]
For $2\le k\le d$, at $m=2$,
\[
 \max\{1,\sqrt{d-k}\}\le\delta_k^\SC(H_{f_{d,2}})
 \le\delta_k^\SR(H_{f_{d,2}})\le\sqrt{d-k+2}.
\]
Thus no multiplicative constant independent of $d$ bounds $\delta_k/\sigma_{k+1}$ for a fixed $k\ge2$. For each fixed $d\ge1$ and $\mathbb F\in\{\SR,\SC\}$,
\begin{equation}\label{I:S:eq:sharp}
 \sup_{\substack{f\ne0\ \text{real and homogeneous}\\ \deg f\ge2}}
 \frac{\SdeltaF{2}(H_f)}{\sigma_3(H_f)}=\sqrt d.
\end{equation}
The same supremum holds within the nonnegative commuting nilpotent target class.
\end{theorem}
\begin{proof}
Give each letter a chain of $m$ edges sharing only its source and sink with the other chains. The $d(m-1)+2$ states realize the indicated series; mixed products are zero and the letter matrices commute. The endpoint Hankel blocks have singular values $\sqrt d$, while each of the $m-1$ interior blocks is an identity on its $d$ pure words. This proves the rank and spectrum. For $m=2$, $C_{1,1}=I_d$ gives $\sqrt{d-k}$ by Theorem~\ref{I:S:thm:coefficient}. Keep $k-2$ pure-square terms to obtain a rank-$k$ candidate; the remaining $d-k+2$ terms have norm $\sqrt{d-k+2}$. At $k=2$ use zero. Letting $d\to\infty$ proves the unbounded ratio.

For every homogeneous target, zero and Lemma~\ref{I:S:lem:homogeneous} give $\delta_2\le\sqrt d\,\sigma_3$. For the family above, Theorem~\ref{I:S:thm:trace} and the spectral bound give
\[
 \max\{1,\sqrt d(m-1)/(m+3)\}\le\delta_2\le\sqrt d,
 \qquad\sigma_3=1.
\]
Let $m\to\infty$ with $d$ fixed. This proves \eqref{I:S:eq:sharp}, including its restricted-class assertion.
\end{proof}


For $d=100$, tolerance two, the spectral state lower bound is only two, while Theorem~\ref{I:S:thm:coefficient} excludes all ranks at most 95. Keeping 96 terms gives a rank-98 candidate with error two. Thus the true tolerance budget lies between 96 and 98. Equation~\eqref{I:S:eq:sharp} is a supremum, not an equality for every homogeneous target. It does not establish a finite factor for all binary mixed-length targets.

\subsection{A strict binary rank-two lower bound without a solver}
Let $f_0=\mathbf1_{aa}+\mathbf1_{bb}$, $h=(e_{aa}+e_{bb})/\sqrt2$, $T=\SSpan\{e_\Seps,h\}$ and $M=\SSpan\{e_a,e_b\}$. Its exact support is $S=T\oplus M$, and
$\Snorm{H_{f_0}x}^2=\Snorm x^2+\Snorm{P_Tx}^2$ on $S$.

\begin{proposition}[Uniform binary geometric obstruction]\label{I:S:prop:binary}
Let $\tau$ be the unique root in $(1,\sqrt2)$ of
\[
 \Phi(t)=\sqrt2-t-\sqrt2(\sqrt2+t)(t^2-1).
\]
Then $\delta_2^\SC(H_{f_0})\ge\tau>2627/2500=1.0508$.
\end{proposition}
\begin{proof}
For a rank-at-most-two candidate $G$, write $t=\Snorm{H_{f_0}-G}$. Only $1\le t<\sqrt2$ needs consideration. There are two-dimensional spaces $K_R\subset\ker G\cap S$ and $K_L\subset\ker G^*\cap S$. On either, $\Snorm{P_Tx}\le q\Snorm x$, $q^2=t^2-1<1$. If the columns of $\binom XY$ are an orthonormal basis of such a space, $X^*X\preceq q^2I$ and $Y^*Y\succeq(1-q^2)I$. The square matrix $Y$ has $YY^*\succeq(1-q^2)I$, proving $\Snorm{(I-P_K)P_M}\le q$. Therefore
\[
 \Snorm{P_MGP_M}\le q^2\Snorm G\le(t^2-1)(\sqrt2+t).
\]
By the Hankel equalities, the four coefficients of $F=P_MGP_M$ are exactly the degree-two coefficients of $Ge_\Seps$. Hence
\[
 \sqrt2-t\le\Snorm F_F\le\sqrt2\Snorm F
 \le\sqrt2(t^2-1)(\sqrt2+t).
\]
Thus $\Phi(t)\le0$. On $[1,\sqrt2]$, $\Phi$ is strictly decreasing and changes sign. Exact rational squaring with checked positive signs gives
$2627/2500<\tau<105083/100000$; the endpoint inequalities follow by direct rational substitution. The bound is uniform and includes complex candidates.
\end{proof}


\section{Constructive equality and controlled-gap families}\label{I:S:sec:constructive}
\subsection{The length-only equality locus}
\begin{proposition}[Radial unary reduction]\label{I:S:prop:radial}
If $f(w)=\varphi(|w|)$ and $\psi(n)=d^{n/2}\varphi(n)$ has bounded finite-rank unary Hankel, then
$H_f\simeq H_\psi\oplus0$ and $\SdeltaF{k}(H_f)=\sigma_{k+1}(H_f)$ whenever scalar AAK applies over the chosen field.
\end{proposition}
\begin{proof}
The isometry $Ve_n=d^{-n/2}\sum_{|w|=n}e_w$ has $V^*H_fV=H_\psi$ by direct summation. $H_f$ vanishes on the complement of the length-constant subspace and has range in it. A unary rank-$k$ AAK optimizer $h$ lifts to $g(w)=d^{-|w|/2}h(|w|)$ with every letter transition equal to the unary transition divided by $\sqrt d$. The lift preserves error and rank; the singular-value lower bound supplies the reverse inequality. Only the unary optimizer is imported from classical AAK~\cite{lacroce2024}.
\end{proof}


\subsection{An exact mixed-length optimum with seven target states}
Take real parameters $s,\eta,\rho$, with $|\rho|<1$, and let
\[
 g_s(w)=s\{(1/2)^{\#a(w)}(1/4)^{\#b(w)}
             +(1/4)^{\#a(w)}(1/2)^{\#b(w)}\},\qquad s\ne0,
\]
and define $h_\rho(a^3)=-\rho$, $h_\rho(a^j)=(1-\rho^2)\rho^{j-4}$ for $j\ge4$, zero on other words, with $|\rho|<1$. At $\rho=0$ it is $\mathbf1_{a^4}$.

\begin{theorem}[A globally optimal explicit two-state model]\label{I:S:thm:mixedexact}
For $s\eta\ne0$, $f=g_s+\eta h_\rho$ has rank seven, commuting nonproportional transitions, and
\[
 \delta_2^\SC(H_f)=\delta_2^\SR(H_f)
  =\sigma_3(H_f)=\Snorm{H_f-H_{g_s}}=|\eta|.
\]
If $|\eta|<2|s|/33$, one state cannot meet tolerance $|\eta|$.
\end{theorem}
\begin{proof}
The two-state realization of $g_s$ has diagonal transitions
$B_a=\Sdiag(1/2,1/4)$, $B_b=\Sdiag(1/4,1/2)$, initial row $(s,s)$ and final column $(1,1)^\top$. The residual has $A_b=0$ and
\[
 T_\rho=\begin{pmatrix}0&1&0&0&0\\0&0&1&0&0\\0&0&0&1&0\\0&0&0&0&1\\0&0&0&0&\rho\end{pmatrix},
 \quad c=e_1,\quad v=(0,0,0,-\rho,1-\rho^2)^\top.
\]
Its exact Stein solutions are
$P=\Sdiag(1,1,1,1,(1-\rho^2)^{-1})$ and
$Q=\Sdiag(1,1,1,1,1-\rho^2)$, so $PQ=I_5$. Thus $H_{h_\rho}$ has norm one and is an isometry on a five-dimensional input subspace. For any rank-two operator $G$, the difference $H_{g_s}-G$ has rank at most four and has a unit kernel vector in that subspace. The error on that vector is $|\eta|$. The candidate $H_{g_s}$ attains it, even among unstructured rank-two operators, proving the singular-value equality as well.

The direct-sum realization has seven states. The principal section indexed by $\Seps$, $a$, $a^2$, $a^3$, $a^4$, $b$, $b^2$ has determinant
$s^2\eta^5(1-\rho^2)/1024\ne0$: eliminate the invertible $(b,b^2)$ block of the rank-two part, whose determinant is $s^2/1024$, leaving the residual five-block of determinant $1-\rho^2$. This establishes minimality. The two geometric kernels have Gram matrix $\begin{psmallmatrix}16/11&4/3\\4/3&16/11\end{psmallmatrix}$, so the smaller singular value of $H_{g_s}$ is $4|s|/33$. Weyl's bound gives $\sigma_2(H_f)\ge4|s|/33-|\eta|>|\eta|$ under the final assumption.
\end{proof}


More generally, if $R$ has norm one and is an isometry on a space of dimension $2k+1$, then every $G_0$ of rank at most $k$ is a globally optimal rank-$k$ approximation to $G_0+\eta R$, with error $|\eta|$. The preceding kernel argument proves this general plateau principle. The family is deliberately solvable; it does not solve arbitrary mixed-length approximation.

\subsection{Recognizable low-rank background with one homogeneous residual}
For $r\ge1$, $m\ge2r$, let $\mathcal C_{r,m}$ consist of bounded rational real series $f$ such that some bounded $g_0$ of rank at most $r$ satisfies $f(w)=g_0(w)$ whenever $|w|\ne m$.

\begin{theorem}[Reconstruction with a certified additive gap]\label{I:S:thm:constructgap}
The background $g_0$ is unique and reconstructible from coefficients of lengths at most $2r-1$. Given the full exact WFA, membership is decidable by finite linear algebra. With $h=f-g_0$, $U=\Snorm h_2$, one obtains
\[
 L:=U\frac{m+1-2r}{m+1+2r}\le\delta_r^\SC(H_f)
 \le\delta_r^\SR(H_f)\le\Snorm{H_f-H_{g_0}}=U,
 \qquad U-L=\frac{4rU}{m+1+2r}.
\]
\end{theorem}
\begin{proof}
Two rank-at-most-$r$ backgrounds agree on every coefficient up to length $2r-1$, since $m\ge2r$. Their difference is identically zero by Theorem~\ref{I:S:thm:realization}. For reconstruction set $V=\Sigma^{\le r-1}$ and $K=(f(pq))_{p,q\in V}$. Under the promise, $K$ has rank $s=\Srank H_{g_0}\le r$. Select an invertible $s$-square submatrix $M=H_f[P,Q]$ and form $M_a=(f(paq))_{P,Q}$. The realization
\[
 \widehat B_a=M^{-1}M_a,\qquad
 \widehat\alpha^\top=(f(q))_{q\in Q},\qquad
 \widehat\beta=M^{-1}(f(p))_{p\in P}
\]
is similar to the minimal background realization: factor $M=O_PC_Q$ and cancel the invertible factors. The zero-rank case gives zero. Reject if $\Srank K>r$.

Without the promise, form $D=F-\widehat G$. For each length less than $m$, generate its exact-length row space and test the output. Generate the row space at length $m+1$, close it under all letters, and test the output on that closure. These finite tests are equivalent to vanishing outside length $m$. If they pass, the candidate is bounded as the difference of a bounded Hankel and a polynomial Hankel. Its residual norm is $U$ by Lemma~\ref{I:S:lem:homogeneous}. For any competitor $q$ of rank at most $r$, the Hankel $H_q-H_{g_0}$ has rank at most $2r$. Apply the homogeneous trace bound at budget $2r$ to $H_h$ and take infima.
\end{proof}


The residual norm can be computed without enumerating its degree-$m$ words: for $h(w)=c^\top D_wb$, start $Q_0=bb^\top$, iterate $Q_{j+1}=\sum_aD_aQ_jD_a^\top$, and take $U^2=c^\top Q_mc$. A perturbation of the target by Hankel norm at most $e$ transfers the bracket to $[\max(0,L-e),U+e]$. For a fixed $m$, the displayed gap is not arbitrarily small merely by rerunning the construction.

\begin{proposition}[Directional trace improvement]\label{I:S:prop:directional}
If $h=f-g_0$ is homogeneous of degree $m$, $U=\Snorm h_2>0$, $A_+\ge\Snorm{H_f}$, and $c_m=|\sum_{|w|=m}\overline{h(w)}f(w)|/U$, then
\[
 \delta_r^\SC(H_f)\ge
 \max\{0,((m+1)c_m-rA_+)/(m+1+r)\}.
\]
\end{proposition}
\begin{proof}
For an error $t$, Cauchy--Schwarz on the empty column gives $U(c_m-t)\le|\Sip h q|$. The trace identity and rank bound give $(m+1)|\Sip h q|\le Ur(A_++t)$. Rearrange.
\end{proof}


For $g_0=g_s$, $s=33/736$, and $h=\mathbf1_{a^{397}}+\mathbf1_{b^{397}}$, one has $\Snorm{H_{g_0}}=1/8$, $U=\sqrt2$, and $\sigma_3(H_f)\le9/8$. Positivity gives $c_m\ge U$, so
\[
 \frac98\ge\sigma_3(H_f),\qquad
 \frac{99}{100}\sqrt2-\frac1{1600}\le\delta_2^\SC(H_f)
 \le\delta_2^\SR(H_f)\le\sqrt2.
\]
The left error endpoint exceeds $1.399445$, and the certified width is below $0.014768$. This is a different target from the fixed small target of the next section. Increasing $m$ changes the target and cannot be described as closing the interval for one fixed target.

\section{The question and the exact answer}
Let $\Sigma$ be any finite alphabet containing distinct letters $a,b$. All word spaces use counting measure. For a scalar series $f:\Sigma^*\to\BC$, put
\[
H_f[u,v]=f(uv),\qquad
\delta_k^{\BC}(H_f)=\inf\{\Bnorm{H_f-H_g}:H_g\text{ bounded word-Hankel},\ \Brank H_g\le k\}.
\]
The norm is the operator norm on $\BHH$. For real $f$, let $\delta_k^{\BR}$ denote the same infimum restricted to real coefficients. Without a superscript, $\delta_k$ means $\delta_k^{\BC}$. In this part use the source \eqref{I:eq:binary-source}, with fixed $k\ge1$, $0<r<1$, and $q=1-r^2$.
The asymptotic is always at fixed $k$, rather than uniformly over growing budgets.

For $0<p<\infty$, the classical Hardy quasi-norm is
\[
\Bnorm{F}_{H^p}=
\sup_{0<\rho<1}\left(\int_0^{2\pi}|F(\rho e^{it})|^p\frac{dt}{2\pi}\right)^{1/p},
\qquad C(m,p)=\sup_{\Bnorm{F}_{H^p}=1}|[z^m]F|.
\]
The circle measure is normalized. In particular $C(0,p)=1$. This is the convention of Brevig--Saksman~\cite{brevig2020}; it is fixed independently of the limit we wish to prove.

\begin{theorem}[Exact boundary reduction and its Hardy identification]
\label{I:B:thm:main}
For $\mathbb F\in\{\mathbb R,\mathbb C\}$ define
\[
 E_m^{\mathbb F}=\inf\{\|z^m-h\|_{H^2}:h\in H^2\text{ over }\mathbb F,
                    \ \operatorname{rank}(h_{i+j})_{i,j\ge0}\le m\}.
\]
For every fixed integer $k\ge1$, the self-contained reduction gives
\[
 \boxed{\displaystyle\lim_{r\uparrow1}\delta_k^{\mathbb F}(H_{f_{r,k}})
       =E_{k-1}^{\mathbb F}\ge\frac1{C(k-1,2/3)}>0.}
\]
For $k\le4$, the published low-degree extremals of
Brevig--Saksman~\cite{brevig2020} give
$E_{k-1}^{\mathbb R}=E_{k-1}^{\mathbb C}=1/C(k-1,2/3)$.
For arbitrary $k$, the same equality uses the outer-extremal theorem
of Brevig--S\o vik~\cite{outer2026}, together with the classical
polynomial structure of coefficient extremals. This is an external
theorem input; the reduction to $E_{k-1}^{\mathbb F}$ does not use it.
Both infima use ordinary word-Hankel rank. No equality of the real and
complex optima at finite $r$ is asserted.
\end{theorem}

The proof has two independent reductions. First, with $m=k-1$,
\[
\lim_{r\uparrow1}\delta_k^{\BC}(H_{f_{r,k}})
=E_m:=\inf_{\Brank H_h\le m}\Bnorm{z^m-h}_{H^2},
\]
where $H_h[i,j]=h_{i+j}$ is a one-letter coefficient Hankel and $h\in H^2$. Second,
\[
E_m=\frac1{C(m,2/3)}.
\]
The first reduction follows from finite-dimensional shift geometry and an explicit construction. The second uses classical finite model spaces and the outer polynomial extremal supplied by~\cite{outer2026}. The proof below distinguishes those imported ingredients from the identities established here.

The first four limits are obtained from the primary low-order constants in~\cite{brevig2020}:
\begin{center}
\begin{tabular}{cl}
\toprule Budget $k$ & Exact limiting structured error\\\midrule
1 & $1$\\
2 & $\sqrt3/2$\\
3 & $3/4$\\
4 & $\sqrt{1153/[2(1103+33\sqrt{33})]}$\\\bottomrule
\end{tabular}
\end{center}
Remark 2 of~\cite{rational2026} concerns degree-one rational approximation of $z^N$. Its value at $N=2$ also equals $\sqrt3/2$, but this is a different approximation problem from $E_1$, where the target is $z$. The contribution here is the precise word-Hankel boundary reduction and coefficient-norm identification.

\section{Exact geometry of the target}
Let $s=\sqrt q$. Define vectors in separate input and output copies of $\BHH$ by
\[
u_j=s\sum_{n\ge0}r^n e_{a^n b^j},\quad 0\le j\le k,
\qquad v_{r,k}=s\sum_{n\ge0}r^n e_{a^n b^k}.
\]
All these normalized vectors with distinct $j$ have disjoint supports. Splitting the possible column words gives
\begin{equation}\label{I:B:eq:target-svd}
H_{f_{r,k}}=u_0v_{r,k}^*+s\sum_{j=0}^{k-1}u_{k-j}e_{b^j}^*.
\end{equation}
The input vectors $v_{r,k},e_\Beps,e_b,\ldots,e_{b^{k-1}}$ are orthonormal, as are the corresponding output vectors. Thus
\[
\Bnorm{H_{f_{r,k}}}=1,\quad \Brank H_{f_{r,k}}=k+1,\quad
\sigma_1=1,\quad\sigma_2=\cdots=\sigma_{k+1}=s,
\quad u_0^*H_{f_{r,k}}=v_{r,k}^*.
\]
Equation~\eqref{I:B:eq:target-svd} proves boundedness as well as the spectral assertions. It also exhibits the difference between the unstructured tail, which tends to zero, and the positive structured boundary limit in Theorem~\ref{I:B:thm:main}.

\section{A lower bound for all complex competitors}
\label{I:B:sec:complex-lower}

We reduce the boundary obstruction to a finite-dimensional one-letter
extremal problem. All coefficients and comparison realizations in this
section may be complex. The alphabet may contain additional letters besides
$a,b$; their nonnegative energy terms can be discarded in the estimates.

\subsection{A one-letter distance and its passive realization formula}

Identify $\BHtwo$ with $\ell^2(\mathbb N_0)$ by Taylor coefficients, and let
$S^*$ denote its backward shift. For an integer $m\ge0$, define
\[
 \mathcal R_m=\left\{h\in\BHtwo:
 \dim\operatorname{span}\{(S^*)^j h:j\ge0\}\le m\right\},
 \qquad
 E_m=\inf_{h\in\mathcal R_m}\Bnorm{z^m-h}_{H^2}.
\]
Thus $\mathcal R_0=\{0\}$ and $E_0=1$. A finite backward-shift cyclic
span is precisely the scalar stable rational-series condition with its
dimension equal to the one-letter Hankel rank. To see boundedness directly,
restrict $S^*$ to this finite invariant span. It is a contraction with no
unit-circle eigenvector, since the corresponding geometric sequence would
not be square summable. Its spectral radius is therefore below one, and
the generated column sequence defines a bounded finite-rank Hankel operator.

For $m\ge1$, also define
\begin{equation}
 \Gamma_m=
 \max_{\substack{D\in\BC^{m\times m},\ l\in\BC^{1\times m}\\
                 D^*D+l^*l\preceq I_m}}
              \Bnorm{lD^m},
 \qquad \Gamma_0=0.
 \label{I:B:eq:passive-gamma}
\end{equation}
The feasible set is compact and the objective continuous, so the maximum
exists. Smaller state spaces are included by padding with zero coordinates.

\begin{lemma}[One-letter projection identity]
\label{I:B:lem:one-letter-projection}
For every $m\ge0$,
\begin{equation}
 E_m^2=1-\Gamma_m^2.
 \label{I:B:eq:E-gamma}
\end{equation}
The infimum defining $E_m$ is attained. Moreover,
$\kappa_m=\Gamma_m^2=1-E_m^2$ for $m\ge1$, with $\kappa_m$ from Lemma~\ref{I:M:rev13:lem:kappa}.
The projection and attainment argument also holds over $\mathbb R$, with field-specific $E_m^{\mathbb R}$ and $\Gamma_m^{\mathbb R}$.
\end{lemma}

\begin{proof}
The statement is immediate for $m=0$. Let $m\ge1$, and let $K$ range over
the finite-dimensional $S^*$-invariant subspaces of $\BHtwo$ with dimension
at most $m$. Every $h\in\mathcal R_m$ lies in such a space, and every
vector in such a space belongs to $\mathcal R_m$. Consequently
\begin{equation}
 E_m^2=1-\sup_K\Bnorm{P_Kz^m}_{H^2}^2.
 \label{I:B:eq:E-model-projection}
\end{equation}
On a particular $K$, set $D=S^*|_K$ and $lx=x(0)$. The identity
$D^*D+l^*l=I_K$ holds, and
\[
 \Bnorm{P_Kz^m}_{H^2}
   =\sup_{x\in K,\,\Bnorm{x}_{H^2}\le1}|x_m|
   =\Bnorm{lD^m}.
\]
Thus the supremum in \eqref{I:B:eq:E-model-projection} is at most $\Gamma_m$.

Conversely, any passive tuple in \eqref{I:B:eq:passive-gamma} has a contraction
observation map
\[
 O:\BC^m\longrightarrow\BHtwo,\qquad
 (Ox)(z)=\sum_{j\ge0}lD^jx\,z^j.
\]
Indeed, telescoping $D^*D+l^*l\preceq I_m$ bounds every finite output
energy by $\Bnorm{x}^2$. Its range $K=\operatorname{ran}O$ is closed,
finite dimensional and backward-shift invariant, since $S^*O=OD$.
For every unit vector $x$,
\[
 |lD^mx|=|(Ox)_m|
 \le\Bnorm{P_Kz^m}_{H^2}\Bnorm{Ox}_{H^2}
 \le\Bnorm{P_Kz^m}_{H^2}.
\]
This proves the reverse inequality and hence \eqref{I:B:eq:E-gamma}. The spaces used here are exactly the finite backward-shift invariant model spaces in Lemma~\ref{I:M:rev13:lem:kappa}; spaces of smaller dimension may be enlarged, proving $\kappa_m=\Gamma_m^2$. Over the real field the same observation-map argument uses real invariant spaces and real tuples.
Choose a maximizing tuple in \eqref{I:B:eq:passive-gamma}; its associated
$K$ has projection norm exactly $\Gamma_m$. Therefore $P_Kz^m$ attains
$E_m$. Any unobservable unitary part of $D$ disappears in $O$, so no strict
stability assumption on the feasible tuple is needed.
\end{proof}


\begin{lemma}[An attained minimizer uses its full rank budget]
\label{I:B:lem:exact-rank-minimizer}
If $m\ge1$, every minimizer of $E_m$ has backward-shift cyclic dimension
exactly $m$. In particular $E_m>0$.
\end{lemma}

\begin{proof}
Let $h$ be a minimizer and $K_h$ its cyclic span. Projection onto $K_h$
is admissible, so minimality forces $h=P_{K_h}z^m$. The residual
$R=z^m-h$ is nonzero, because $z^m$ itself has cyclic dimension $m+1$.
If $\dim K_h<m$, choose $\lambda\in\mathbb D$ with $R(\lambda)\ne0$.
The reproducing kernel $k_\lambda(z)=(1-\overline\lambda z)^{-1}$ is
a backward-shift eigenvector. The invariant space
$K_h+\operatorname{span}\{k_\lambda\}$ has dimension at most $m$,
and projection onto it strictly improves the approximation, since
$R\perp K_h$ but $\langle R,k_\lambda\rangle\ne0$. This is a
contradiction. Attainment and $h\ne z^m$ also give $E_m>0$.
\end{proof}


\subsection{The exact resolvent defect inequality}

Let $G$ be any bounded word-Hankel operator with rank at most $k$, and
let $\mathcal S_G=\operatorname{ran}G^*$ be its input row space.
For the left creation operators $L_c e_w=e_{cw}$, the Hankel relation
implies that $\mathcal S_G$ is invariant under $L_c^*$. Set
\[
 A=L_a^*|_{\mathcal S_G},\qquad
 B=L_b^*|_{\mathcal S_G},\qquad
 \ell x=x[\Beps].
\]
The partition of nonempty words according to first letter gives
\begin{equation}
 A^*A+B^*B+\ell^*\ell\preceq I_{\mathcal S_G}.
 \label{I:B:eq:boundary-passivity}
\end{equation}
The inequality becomes equality when the alphabet is exactly $\{a,b\}$.

Write $q=1-r^2$. For every contraction $A$,
\begin{equation}
 (I-rA^*)(I-rA)-q(I-A^*A)
       =(A-rI)^*(A-rI)\succeq0.
 \label{I:B:eq:resolvent-defect-exact}
\end{equation}
Since $B^*B\preceq I-A^*A$, congruence by $(I-rA)^{-1}$ yields
\begin{equation}
 \Bnorm{\sqrt q\,B(I-rA)^{-1}}\le1.
 \label{I:B:eq:resolvent-defect-contraction}
\end{equation}
This is a joint passivity consequence; it does not follow merely from
separate spectral-radius bounds on the letter matrices.

\subsection{Compactness at the boundary}

\begin{theorem}[Boundary lower bound over all complex competitors]
\label{I:B:thm:all-k-boundary-lower}
For every fixed integer $k\ge1$,
\begin{equation}
 \liminf_{r\uparrow1}\delta_k(H_{f_{r,k}})
          \ge E_{k-1}=\sqrt{1-\Gamma_{k-1}^2}.
 \label{I:B:eq:all-k-boundary-lower}
\end{equation}
\end{theorem}

\begin{proof}
Let
\[
 v_{r,k}=\sqrt q\sum_{n\ge0}r^n e_{a^n b^k},\qquad
 u_{r,0}=\sqrt q\sum_{n\ge0}r^n e_{a^n}.
\]
The target's singular decomposition gives
$\Bnorm{v_{r,k}}=\Bnorm{u_{r,0}}=1$ and
$u_{r,0}^*H_{f_{r,k}}=v_{r,k}^*$. If $P_G$ projects onto
$\mathcal S_G$, then $G(I-P_G)=0$, and therefore
\begin{equation}
 \Bnorm{H_{f_{r,k}}-G}
   \ge\sqrt{1-\Bnorm{P_Gv_{r,k}}^2}.
 \label{I:B:eq:boundary-kernel-lower}
\end{equation}
The coordinates of $x\in\mathcal S_G$ satisfy
$x[a^n b^k]=\ell B^kA^nx$. Hence
\begin{equation}
 \Bnorm{P_Gv_{r,k}}
 =\Bnorm{\sqrt q\,\ell B^k(I-rA)^{-1}}
 \le\Bnorm{\ell B^{k-1}},
 \label{I:B:eq:boundary-overlap-bound}
\end{equation}
where the final inequality is \eqref{I:B:eq:resolvent-defect-contraction}.

Consider any sequence $r_j\uparrow1$ and arbitrary rank-at-most-$k$
competitors $G_j$. Zero-rank terms have zero overlap. For all remaining
terms, choose orthonormal coordinates on their row spaces and pad by
zero coordinates to dimension $k$. The padded tuples still satisfy
\eqref{I:B:eq:boundary-passivity}. Passing to a subsequence gives
$A_j\to A_0$, $B_j\to B_0$, $\ell_j\to\ell_0$.

If $I-A_0$ is invertible, the resolvent in
\eqref{I:B:eq:boundary-overlap-bound}'s middle expression stays bounded,
so that overlap tends to zero. Otherwise
$\mathcal E=\ker(I-A_0)$ is nonzero and reduces the contraction $A_0$.
Indeed $A_0x=x$ and contraction imply $A_0^*x=x$.
The limiting passive inequality forces $B_0\mathcal E=0$ and
$\ell_0\mathcal E=0$. In the orthogonal decomposition
$\mathcal E\oplus\mathcal F$ one can consequently write
\[
 B_0=\begin{pmatrix}0&C\\0&D\end{pmatrix},\qquad
 \ell_0=(0,l),\qquad \dim\mathcal F\le k-1.
\]
The $\mathcal F$-column part of passivity gives
$D^*D+l^*l\preceq I_{\mathcal F}$. If $k\ge2$,
$\ell_0B_0^{k-1}=(0,lD^{k-1})$ has norm at most $\Gamma_{k-1}$.
If $k=1$, $\mathcal F=\{0\}$ and $\ell_0=0$, giving the same conclusion
with $\Gamma_0=0$.

Thus every subsequential limiting overlap is at most $\Gamma_{k-1}$.
A compactness contradiction makes this uniform over all competitors:
\[
 \limsup_{r\uparrow1}\sup_{\substack{G\text{ bounded word-Hankel}\\
                                  \operatorname{rank}G\le k}}
       \Bnorm{P_Gv_{r,k}}\le\Gamma_{k-1}.
\]
Insert this bound into \eqref{I:B:eq:boundary-kernel-lower} and use
Lemma~\ref{I:B:lem:one-letter-projection}. No existence assumption for a
best word-Hankel competitor is required.
\end{proof}


For $k=2$, the passive problem has one state:
\[
 \Gamma_1=\max_{|d|^2+|l|^2\le1}|ld|=\frac12.
\]
The theorem therefore already gives
$\liminf_{r\uparrow1}\delta_2(H_{f_{r,2}})\ge\sqrt3/2$ without any
Hardy coefficient theorem.

\section{An explicit upper construction}
\label{I:B:sec:lift}

The lower reduction isolates a scalar one-letter approximation problem.
The construction below gives the matching direction without an abstract
extension theorem. It also controls the error before taking the boundary
limit.

\begin{theorem}[A finite-rank lift with an exact orthogonal split]
\label{I:B:thm:upper-lift}
Fix an integer $k\ge1$, put $m=k-1$, and let $0<r<1$,
$q=1-r^2$, and $s=\sqrt q$. Let $h=(h_j)_{j\ge0}$ have a bounded
scalar Hankel operator
\[
 H_h[i,j]=h_{i+j},\qquad i,j\ge0,
 \qquad \operatorname{rank}H_h\le m.
\]
For $m=0$ take $h=0$. Define the word series
\begin{equation}
 g_{r,h}(a^n b^j)=q r^n h_{j-1}
 \quad(n\ge0,\ j\ge1),\qquad
 g_{r,h}(w)=0\quad\text{otherwise}.
 \label{I:B:eq:lift-series}
\end{equation}
Then $H_{g_{r,h}}$ is bounded, has rank at most $m+1=k$, and
\begin{equation}
 \boxed{\displaystyle
 \Bnorm{H_{f_{r,k}}-H_{g_{r,h}}}^2
 \le \Bnorm{z^m-h}_{\BHtwo}^2
       +q\Bnorm{H_{z^m}-H_h}^2.}
 \label{I:B:eq:lift-finite-error}
\end{equation}
Consequently, if
\[
 E_m=\inf_{\substack{H_h\text{ bounded scalar Hankel}\\
                      \operatorname{rank}H_h\le m}}
        \Bnorm{z^m-h}_{\BHtwo},
\]
then
\begin{equation}
 \limsup_{r\uparrow1}\delta_k(H_{f_{r,k}})\le E_{k-1}.
 \label{I:B:eq:lift-limsup}
\end{equation}
All scalar coefficients and competing realizations may be complex.
\end{theorem}

\begin{proof}
If $h\ne0$, choose a scalar realization
$h_j=\alpha^*D^j\beta$ of dimension $n\le m$. The block matrices
\begin{equation}
 A_a=\begin{pmatrix}r&0\\0&0_n\end{pmatrix},\qquad
 A_b=\begin{pmatrix}0&\alpha^*\\0&D\end{pmatrix},\qquad
 \alpha_g^*=q(1,0_n),\qquad
 \beta_g=\begin{pmatrix}0\\\beta\end{pmatrix}
 \label{I:B:eq:lift-realization}
\end{equation}
give a realization of \eqref{I:B:eq:lift-series} in the convention
$g(a_1\cdots a_t)=\alpha_g^*A_{a_1}\cdots A_{a_t}\beta_g$.
Indeed $A_bA_a=0$, so a letter $a$ after a letter $b$ forces the
coefficient to vanish; the empty coefficient also vanishes. Direct block
multiplication gives the stated coefficients on $a^*b^+$. All additional
letter matrices are zero if the alphabet is larger. The case $h=0$ gives
the zero series. Thus the realization dimension is at most $m+1$.

We prove boundedness and the error estimate directly in the infinite
operator norm. Write $\ell^2(\mathbb N_0)$ for the scalar coefficient
space. Define two isometries from that space to $\BHH$ by
\begin{equation}
 V_r e_j=s\sum_{n\ge0}r^n e_{a^n b^{j+1}},\qquad
 W e_j=e_{b^j},\qquad j\ge0,
 \label{I:B:eq:lift-isometries}
\end{equation}
and define the unit vector
\[
 u_r=s\sum_{n\ge0}r^n e_{a^n}.
\]
The geometric normalization gives $V_r^*V_r=W^*W=I$, and
$u_r\perp\operatorname{ran}V_r$. Boundedness of $H_h$ implies
$h=H_he_0\in\ell^2(\mathbb N_0)$. With $\overline h$ denoting
coefficientwise conjugation, exact entry matching yields
\begin{equation}
 H_{g_{r,h}}
   =u_r(V_r\overline h)^*+sV_rH_hW^*.
 \label{I:B:eq:lift-exact-decomposition}
\end{equation}
The conjugation in the first term is necessary. At row $a^n$ and column
$a^t b^j$, that term equals $q r^{n+t}h_{j-1}$, rather than its
conjugate. At row $a^n b^j$ with $j\ge1$ and column $b^\ell$, the
second term equals $q r^n h_{j+\ell-1}$. These account for every
nonzero entry. In particular \eqref{I:B:eq:lift-exact-decomposition} proves
boundedness, so the finite realization gives a bounded word-Hankel
operator of rank at most $k$.

For the scalar polynomial $z^m$, its coefficient vector is $e_m$, and
the same identity reads
\begin{equation}
 H_{f_{r,k}}=u_r(V_re_m)^*+sV_rH_{z^m}W^*.
 \label{I:B:eq:lift-target-decomposition}
\end{equation}
Subtracting \eqref{I:B:eq:lift-exact-decomposition} from
\eqref{I:B:eq:lift-target-decomposition}, the two summands have orthogonal
output ranges. For each input $x\in\BHH$ this gives
\begin{align*}
 \Bnorm{(H_{f_{r,k}}-H_{g_{r,h}})x}^2
 &=\bigl|\langle V_r(e_m-\overline h),x\rangle\bigr|^2
   +q\Bnorm{V_r(H_{z^m}-H_h)W^*x}^2\\
 &\le\left(\Bnorm{e_m-\overline h}_2^2
           +q\Bnorm{H_{z^m}-H_h}^2\right)\Bnorm{x}^2.
\end{align*}
Since $e_m$ has real coefficients,
$\Bnorm{e_m-\overline h}_2=\Bnorm{e_m-h}_2
=\Bnorm{z^m-h}_{\BHtwo}$. This proves
\eqref{I:B:eq:lift-finite-error}.

For every fixed admissible $h$, the second term on the right of
\eqref{I:B:eq:lift-finite-error} tends to zero as $r\uparrow1$. First choose
$h$ within any positive tolerance of $E_m$, next take the boundary limit,
and finally let the tolerance decrease. This proves
\eqref{I:B:eq:lift-limsup} without assuming attainment of $E_m$.
\end{proof}


\subsection{A real two-state competitor and a sharper finite-parameter bound}
\label{I:B:subsec:lift-k2}

For $k=2$, the scalar one-state problem is explicit. Set
\begin{equation}
 y=\frac1{\sqrt2},\qquad c=y(1-y^2)=\frac1{2\sqrt2},\qquad
 h_j=cy^j.
 \label{I:B:eq:lift-k2-parameters}
\end{equation}
The lifted series is
\begin{equation}
 g_r(a^n b^j)=q c r^n y^{j-1}
 \quad(n\ge0,\ j\ge1),\qquad g_r(w)=0\quad\text{otherwise}.
 \label{I:B:eq:lift-k2-series}
\end{equation}
It has the real two-state realization
\begin{equation}
 A_a=\begin{pmatrix}r&0\\0&0\end{pmatrix},\qquad
 A_b=\begin{pmatrix}0&1\\0&y\end{pmatrix},\qquad
 \alpha_g^*=qc(1,0),\qquad\beta_g=(0,1)^{\mathsf T}.
 \label{I:B:eq:lift-k2-realization}
\end{equation}
The following stronger bound shows that this particular competitor
approaches its limiting error at order $q^2$.

\begin{proposition}[A quantitative two-state upper bound]
\label{I:B:prop:lift-k2-sharp}
For $g_r$ in \eqref{I:B:eq:lift-k2-series} and $0<q\le1/8$,
\begin{equation}
 \boxed{\displaystyle
 \frac34+\frac{q^2}{6}
 \le\Bnorm{H_{f_{r,2}}-H_{g_r}}^2
 \le\frac34+\frac{q^2}{6}+\frac43q^3.}
 \label{I:B:eq:lift-k2-sharp-bound}
\end{equation}
In particular,
\begin{equation}
 \delta_2(H_{f_{r,2}})^2
 \le\frac34+\frac{q^2}{6}+\frac43q^3,
 \qquad
 \limsup_{r\uparrow1}\delta_2(H_{f_{r,2}})\le\frac{\sqrt3}{2}.
 \label{I:B:eq:lift-k2-limsup}
\end{equation}
The lower inequality in \eqref{I:B:eq:lift-k2-sharp-bound} concerns this
explicit competitor only, rather than the structured optimum.
\end{proposition}

\begin{proof}
Use the orthonormal vectors
\[
 u=s\sum_{n\ge0}r^n e_{a^n},\qquad
 v_j=s\sum_{n\ge0}r^n e_{a^n b^j}\quad(j\ge1),
\]
and put
\[
 v_y=\sum_{j\ge1}y^{j-1}v_j,\qquad
 j_y=\sum_{\ell\ge0}y^\ell e_{b^\ell},\qquad e=e_\Beps.
\]
Their relevant norms and inner products are
\[
 \Bnorm{v_y}^2=\Bnorm{j_y}^2=(1-y^2)^{-1},\qquad
 \langle v_2,v_y\rangle=y.
\]
The exact target and competitor decompositions become
\begin{align*}
 H_{f_{r,2}}&=uv_2^*+s(v_2e^*+v_1e_b^*),\\
 H_{g_r}&=cu v_y^*+scv_yj_y^*.
\end{align*}
Define
\[
 d=v_2-cv_y,\qquad
 R=v_2e^*+v_1e_b^*-cv_yj_y^*,\qquad A=\Bnorm{d}^2.
\]
For the parameters \eqref{I:B:eq:lift-k2-parameters},
\[
 A=1-2cy+\frac{c^2}{1-y^2}=\frac34.
\]
The output vector $u$ is orthogonal to the range of $R$. Therefore, for
$E=H_{f_{r,2}}-H_{g_r}$,
\begin{equation}
 E=ud^*+sR,\qquad M:=E^*E=dd^*+qR^*R.
 \label{I:B:eq:lift-k2-orthogonal-error}
\end{equation}
Moreover,
\begin{equation}
 \Bnorm R\le1+\frac{c}{1-y^2}=1+y,
 \qquad \Bnorm R^2\le(1+y)^2<3.
 \label{I:B:eq:lift-k2-Rnorm}
\end{equation}
Here the first two summands in $R$ form a partial isometry of norm one.

The leading error vector has the useful cancellations
\begin{align*}
 e^*d&=0,& e_b^*d&=-sc,\\
 j_y^*d&=s\left(y^2-\frac{cy}{1-y^2}\right)=0.
\end{align*}
Consequently
\begin{equation}
 Rd=-scv_1.
 \label{I:B:eq:lift-k2-cancellation}
\end{equation}
Let $z=d/\sqrt A$ and decompose $M$ with respect to
$\operatorname{span}\{z\}\oplus z^\perp$. Its scalar upper-left entry
is
\begin{equation}
 a=\langle z,Mz\rangle
   =A+q\Bnorm{Rz}^2
   =\frac34+\frac{q^2c^2}{A}
   =\frac34+\frac{q^2}{6}.
 \label{I:B:eq:lift-k2-Rayleigh}
\end{equation}
This already proves the lower bound for the explicit competitor.
Let $B$ be the compression of $M$ to $z^\perp$, and let
$v=(I-zz^*)Mz$ be its off-diagonal vector. Since $dd^*$ vanishes on
$z^\perp$, \eqref{I:B:eq:lift-k2-Rnorm} and
\eqref{I:B:eq:lift-k2-cancellation} give
\[
 \Bnorm B\le q\Bnorm R^2<3q,
 \qquad
 \Bnorm v^2\le q^2\Bnorm R^2\Bnorm{Rz}^2
             <\frac{q^3}{2}.
\]
For $q\le1/8$ the spectral gap in this block decomposition obeys
\[
 a-\Bnorm B\ge\frac34-3q\ge\frac38.
\]
The elementary positive-block estimate
\[
 \lambda_{\max}(M)
 \le a+\frac{\Bnorm v^2}{a-\Bnorm B}
\]
therefore yields the upper inequality in
\eqref{I:B:eq:lift-k2-sharp-bound}. For completeness, this block estimate
follows from the Schur complement when the largest eigenvalue is above
$a$: if it is $\lambda$, then
$\lambda-a\le\Bnorm v^2/(\lambda-\Bnorm B)
\le\Bnorm v^2/(a-\Bnorm B)$. If $v=0$, the estimate is immediate.
\end{proof}


Thus the explicit upper competitor has the expansion
\[
 \Bnorm{H_{f_{r,2}}-H_{g_r}}
 =\frac{\sqrt3}{2}+\frac{q^2}{6\sqrt3}+O(q^3).
\]
This expansion is not asserted for the optimal structured error; the
matching boundary value follows by combining
\eqref{I:B:eq:lift-k2-limsup} with the independently proved lower bound.

\section{The exact Hardy coefficient constant}
\label{I:B:sec:hardy}

All scalar functions and approximation classes in this section are complex.
The circle carries normalized measure $d\theta/(2\pi)$.  For $0<p<\infty$,
the Hardy space $H^p$ consists of analytic functions on the unit disc for which
\[
 \Bnorm{F}_{H^p}
 =\sup_{0<\rho<1}\left(\int_0^{2\pi}
       |F(\rho e^{i\theta})|^p\,\frac{d\theta}{2\pi}\right)^{1/p}<\infty.
\]
We write $\BHtwo=H^2$ and identify it isometrically with the square-summable
Taylor coefficient sequences.  For $m\ge0$ define the coefficient-functional
norm
\begin{equation}
 C(m,p)=\sup\bigl\{|[z^m]F|:\Bnorm{F}_{H^p}=1\bigr\}.
 \label{I:B:eq:hardy-C-definition}
\end{equation}
This is the convention of Brevig--Saksman~\cite{brevig2020}; in particular,
it is not a norm of a finite Hankel block.  Point evaluation gives $C(0,p)=1$.

\subsection{The scalar rational problem and finite model spaces}

Let $\mathcal R_m$ be the set of analytic functions $h(z)=\sum_{j\ge0}h_jz^j$
whose coefficient Hankel $(h_{i+j})_{i,j\ge0}$ is bounded and has rank at most
$m$.  Set $\mathcal R_0=\{0\}$.  Equivalently, for $m\ge1$ this is the class
of stable rational functions
\begin{equation}
 h=\frac{N}{D},\qquad \deg N<m,\quad \deg D\le m,\quad D(0)=1,
 \label{I:B:eq:hardy-rational-class}
\end{equation}
after cancellation, with $D$ zero-free on the closed disc.  Polynomial parts
are included: a polynomial of degree $m-1$ can have coefficient-Hankel rank
$m$.  Omitting these parts would change the approximation class.  Define
\begin{equation}
 E_m=\inf_{h\in\mathcal R_m}\Bnorm{z^m-h}_{\BHtwo}.
 \label{I:B:eq:hardy-E-definition}
\end{equation}

For a finite Blaschke product $\Theta$, let
$K_\Theta=\BHtwo\ominus\Theta\BHtwo$.  Its degree is its number of zeros
inside the disc, counted with multiplicity.  A product of degree $l$ can be
written, up to a unimodular constant, as
\[
 \Theta=\frac{D_l^*}{D},\qquad
 D_l^*(z)=z^l\overline{D(1/\bar z)},\quad \deg D\le l,
\]
where $D$ is zero-free on the closed disc.  The reciprocal operation uses
the specified order $l$, even when $\deg D<l$.

\begin{lemma}[Classical finite-model-space formulation]
\label{I:B:lem:hardy-model}
The class $\mathcal R_m$ is the union of the spaces $K_\Theta$ for finite
Blaschke products of degree at most $m$.  If
$\Theta(z)=\sum_{j\ge0}\theta_jz^j$, then
\begin{equation}
 \Bnorm{z^m-P_{K_\Theta}z^m}_{\BHtwo}^2
 =\sum_{j=0}^m|\theta_j|^2.
 \label{I:B:eq:hardy-model-error}
\end{equation}
Consequently,
\begin{equation}
 E_m^2=\inf_{\deg\Theta\le m}\sum_{j=0}^m|\theta_j|^2.
 \label{I:B:eq:hardy-Blaschke-variational}
\end{equation}
\end{lemma}
\begin{proof}
For degree $l$, the standard rational description is
\[
 K_\Theta=\left\{\frac{N}{D}:\deg N<l\right\}.
\]
It follows, for example, from the $l$-element Takenaka--Malmquist orthonormal
basis associated with the zeros of $\Theta$: the common denominator is $D$,
the numerators have degree less than $l$, and dimension is $l$.  This
space is invariant under the scalar backward shift.  Hence the Hankel
columns of each member span a space of dimension at most $l$.  Conversely,
the stable rational realization of a bounded finite-rank scalar Hankel
has the form \eqref{I:B:eq:hardy-rational-class}, at its minimal order; its
denominator supplies such a finite Blaschke product.  This proves the union
assertion, including denominators of degree below the specified order.

Multiplication by $\Theta$ is an isometry on $\BHtwo$, so the orthogonal
projection onto $\Theta\BHtwo$ is
\[
 z^m-P_{K_\Theta}z^m
 =\Theta P_+(\overline\Theta z^m)
 =\Theta\sum_{j=0}^m\overline{\theta_j}z^{m-j}.
\]
Taking the squared norm proves \eqref{I:B:eq:hardy-model-error}, and minimizing
over the union proves \eqref{I:B:eq:hardy-Blaschke-variational}.
\end{proof}


This variational formulation belongs to classical rational approximation;
see also Theorem~1 and the monomial specialization in
Baratchart et al.~\cite{rational2026}.  We give the argument to fix the
rank and polynomial-part conventions precisely.

\subsection{A sharp coefficient-energy identity}

\begin{theorem}[Exact identification with the Hardy coefficient norm]
\label{I:B:thm:hardy-bridge}
For every integer $m\ge0$,
\begin{equation}
 \boxed{\displaystyle E_m=\frac1{C(m,2/3)}.}
 \label{I:B:eq:hardy-exact-bridge}
\end{equation}
The lower bound $E_m\ge1/C(m,2/3)$ is self-contained. Equality for $m\le3$ follows from the published low-degree outer extremals in~\cite{brevig2020}. For general $m$, the matching upper construction uses the outer-extremal theorem of Brevig--S\o vik~\cite{outer2026}.
\end{theorem}
\begin{proof}
The case $m=0$ is immediate, so assume $m\ge1$ and abbreviate
$C=C(m,2/3)$.

\emph{Lower bound, by a direct H\"older argument.}
First take $\Theta$ of degree exactly $m$.  Multiplying by a unimodular
constant does not affect \eqref{I:B:eq:hardy-model-error}, so write
\[
 \Theta=\frac{D_m^*}{D},\qquad
 D_m^*(z)=z^m\overline{D(1/\bar z)}.
\]
Let $b=P_{\le m}\Theta=\sum_{j=0}^m\theta_jz^j$.  Since Taylor truncation
does not affect the coefficient of degree $m$ after multiplication by an
analytic polynomial,
\begin{equation}
 [z^m](bD^2)=[z^m](\Theta D^2)
 =[z^m](D_m^*D)=\Bnorm{D}_{\BHtwo}^2.
 \label{I:B:eq:hardy-key-coefficient}
\end{equation}
The definition of $C$ and H\"older's inequality therefore give
\begin{equation}
 \Bnorm{D}_{\BHtwo}^2
 \le C\Bnorm{bD^2}_{H^{2/3}}
 \le C\Bnorm{b}_{\BHtwo}\Bnorm{D}_{\BHtwo}^2.
 \label{I:B:eq:hardy-holder-bound}
\end{equation}
For clarity, the second step is precisely
\[
 \left(\int |b|^{2/3}|D|^{4/3}\right)^{3/2}
 \le\left(\int |b|^2\right)^{1/2}\int |D|^2,
\]
using H\"older exponents $3$ and $3/2$ and normalized circle measure.
Division by $\Bnorm{D}_{\BHtwo}^2>0$ proves
$\Bnorm{b}_{\BHtwo}\ge1/C$.

If $\deg\Theta=l<m$, append any finite Blaschke factor $\Psi$ of degree
$m-l$.  The inclusion $K_\Theta\subset K_{\Theta\Psi}$ shows that the
projection error for $\Theta$ is at least that for the degree-$m$ product
$\Theta\Psi$.  Thus the same lower bound holds for every degree at most
$m$.  Lemma~\ref{I:B:lem:hardy-model} gives $E_m\ge1/C$.  This part of the
proof uses neither stationarity of a rational minimizer nor existence of
such a minimizer.

\emph{Matching upper bound, from an outer coefficient extremal.}
Coefficient extremals exist by a normal-family argument.  Theorem~1 and
Corollary~3 of Brevig--S\o vik~\cite{outer2026}, combined with the
polynomial structure in Brevig--Saksman~\cite{brevig2020}, imply that an
extremal for \eqref{I:B:eq:hardy-C-definition} at $p=2/3$ is $F=P^3$, where
$P$ is an outer polynomial of exact degree $m$.  Normalize its phase and
norm so that
\begin{equation}
 \Bnorm{P}_{\BHtwo}=1,\qquad [z^m]P^3=C>0.
 \label{I:B:eq:hardy-extremal-normalization}
\end{equation}
For every polynomial $Q$ one has
$\Bnorm{Q^3}_{H^{2/3}}=\Bnorm{Q}_{\BHtwo}^3$.  Hence $P$ maximizes
$\operatorname{Re}[z^m]Q^3$ over the entire unit sphere of the
$(m+1)$-dimensional polynomial Hilbert space; the variations do not need
to remain outer.

Write $P(z)=\sum_{j=0}^mp_jz^j$.  Differentiation with respect to the real
and imaginary parts of all $p_j$, with the unit-sphere constraint, yields
\[
 [z^{m-j}]P^2=\lambda\overline{p_j}\qquad(0\le j\le m).
\]
Multiplying by $p_j$ and summing shows
$\lambda=[z^m]P^3=C$, since $\sum_j|p_j|^2=1$.
Thus, for the reciprocal of order $m$,
\begin{equation}
 P^2=C P_m^*+O(z^{m+1}),\qquad
 P_m^*(z)=z^m\overline{P(1/\bar z)}.
 \label{I:B:eq:hardy-EL}
\end{equation}
The ratio $\Theta=P_m^*/P$ is a finite Blaschke product of degree at most
$m$, after cancelling any common factors on the unit circle.  Outerness
ensures $P(0)\ne0$, so division of \eqref{I:B:eq:hardy-EL} as a Taylor series
at zero gives
\[
 \Theta=P/C+O(z^{m+1}).
\]
Consequently,
\[
 \sum_{j=0}^m|\theta_j|^2
 =\frac{\Bnorm{P}_{\BHtwo}^2}{C^2}=\frac1{C^2}.
\]
The projection $h=P_{K_\Theta}z^m$ belongs to $\mathcal R_m$ and attains
this error.  Therefore $E_m\le1/C$, completing the proof.
\end{proof}


\begin{corollary}[A real optimizer may be chosen]
\label{I:B:cor:hardy-real}
The infimum defining $E_m$ is unchanged if $h$ is required to have real
Taylor coefficients.  In the upper construction one may choose $P$ with
nonnegative real coefficients.
\end{corollary}
\begin{proof}
Take a normalized complex extremal polynomial $P=\sum_jp_jz^j$ as above
and put $Q=\sum_j|p_j|z^j$.  Then $\Bnorm{Q}_{\BHtwo}=1$ and
\[
 [z^m]Q^3\ge\bigl|[z^m]P^3\bigr|=C(m,2/3).
\]
On the other hand, $\Bnorm{Q^3}_{H^{2/3}}=1$, so the coefficient inequality
forces equality.  For $m\le3$, the explicit nonnegative outer extremals in~\cite{brevig2020} give this step directly. For general $m$, the outer-extremal theorem makes $Q$ outer.
The corresponding $\Theta=Q_m^*/Q$ has real coefficients, and the formula
in Lemma~\ref{I:B:lem:hardy-model} makes $P_{K_\Theta}z^m$ real as well.
This real function attains the complex lower bound.
\end{proof}


\section{Completion of the theorem and algebraic construction}
The lower-bound and lift arguments may each be carried out over either scalar field: replace the passive tuples, invariant subspaces and scalar minimizers by real ones for the real problem. They give $\lim\delta_k^{\mathbb F}=E_{k-1}^{\mathbb F}$ independently of the outer-extremal input. In particular, the lower-bound section proves $\liminf\delta_k^{\BC}\ge E_{k-1}$ over the full complex comparison class. The lift section proves $\limsup\delta_k^{\BC}\le E_{k-1}$. The Hardy section identifies $E_{k-1}$ with the reciprocal coefficient constant. This proves the complex assertion in Theorem~\ref{I:B:thm:main}.

The Hardy extremal polynomial may be taken to have nonnegative real coefficients, as proved in the Hardy section. Its finite model-space projection of $z^m$ then has real coefficients, and the lifted word series is real. Therefore the same upper bound holds for $\delta_k^{\BR}$. Since $\delta_k^{\BC}\le\delta_k^{\BR}$, both limiting constants coincide. This proves the full theorem.

\begin{corollary}[Finite algebraic description of the limiting constant]
For each fixed $m\ge0$,
\[
C(m,2/3)=
\max_{\substack{p_0,\ldots,p_m\ge0\\\sum_{j=0}^m p_j^2=1}}
\sum_{i+j+\ell=m}p_i p_j p_\ell.
\]
Consequently $C(m,2/3)$ and its reciprocal are real algebraic numbers. Real quantifier elimination can in principle compute them and an algebraic optimal polynomial. This is exact decidability, not a polynomial-time claim.
\end{corollary}
\begin{proof}
Every polynomial cube in the displayed unit sphere has $H^{2/3}$ norm one, so its coefficient is at most $C(m,2/3)$. The outer-extremal theorem and the nonnegative-coefficient reduction supply a polynomial attaining equality. The feasible set is compact semialgebraic over $\mathbb Q$, and the objective is polynomial. Its maximum and algebraic witnesses can be obtained by real-algebraic decision and witness extraction.
\end{proof}


For a normalized extremal $P(z)=\sum_{j=0}^m p_jz^j$, let
$P^{\sharp}(z)=z^m\overline{P(1/\bar z)}$, with the reversal order fixed at $m$. The Hardy proof gives
\[
\Theta=P^{\sharp}/P,\qquad
h(z)=z^m-\frac{(P^{\sharp}(z))^2}{C(m,2/3)P(z)}.
\]
This $h$ is the projection of $z^m$ onto $K_\Theta$, so it has one-letter Hankel rank at most $m$. Apparent boundary factors cancel. A minimal stable realization of $h$, followed by the explicit word lift, gives a family with at most $m+1=k$ states and limiting error exactly $1/C(m,2/3)$. This construction optimizes the smaller scalar boundary problem; it does not solve the original word-Hankel optimum for each $r$.

For any fixed attained scalar minimizer $h$, the orthogonal lift bound gives
\[
\delta_k^{\BC}(H_{f_{r,k}})^2
\le E_{k-1}^2+q\Bnorm{H_{z^{k-1}}-H_h}^2.
\]
Since $E_{k-1}>0$, this immediately gives a one-sided $E_{k-1}+O(q)$ upper estimate. The next section strengthens the construction to an $O(q^2)$ upper excess for every fixed $k$. A matching finite-$r$ lower rate for the optimum is not proved.

\section{A quadratic approach to the boundary for the lifted competitors}
\label{I:B:sec:curvature}

The orthogonal split in Theorem~\ref{I:B:thm:upper-lift} contains more
information than its direct triangle estimate. For every fixed scalar
competitor, the squared error of its lift has no term of order $q$.
This is a statement about the explicit lifted construction; it does not
give an asymptotic expansion of the globally optimal structured error.

\begin{proposition}[Curvature of a fixed scalar lift]
\label{I:B:prop:lift-curvature}
Let $h$ satisfy the assumptions of Theorem~\ref{I:B:thm:upper-lift}, put
$m=k-1$, and write $d=z^m-h$ for its scalar coefficient residual.
Assume
\begin{equation}
 A=\Bnorm d_{\BHtwo}^2>0,\qquad
 L=\Bnorm{H_d}^2,\qquad
 T=\Bnorm{H_dS_+\overline d}_2^2,
 \label{I:B:eq:curvature-parameters}
\end{equation}
where $S_+e_j=e_{j+1}$ is the scalar forward shift and
$\overline d$ denotes coefficientwise conjugation. Let
$E_r=H_{f_{r,k}}-H_{g_{r,h}}$. If $qL\le A/2$, then
\begin{equation}
 \boxed{\displaystyle
 A+\frac{q^2T}{A}
 \le\Bnorm{E_r}^2
 \le A+\frac{q^2T}{A}+\frac{2q^3LT}{A^2}.}
 \label{I:B:eq:curvature-finite-bound}
\end{equation}
Consequently, for each fixed admissible $h$ with $A>0$,
\begin{align}
 \Bnorm{E_r}^2&=A+\frac{T}{A}q^2+O(q^3),
 \label{I:B:eq:curvature-squared-expansion}\\
 \Bnorm{E_r}&=\sqrt A+\frac{T}{2A^{3/2}}q^2+O(q^3)
 \qquad(q\downarrow0).
 \label{I:B:eq:curvature-norm-expansion}
\end{align}
The constants may depend on the fixed scalar residual through $A,L,T$.
\end{proposition}

\begin{proof}
Use the isometries $V_r,W$ and the unit vector $u_r$ from
\eqref{I:B:eq:lift-isometries}. Define $v_d=V_r\overline d$. The exact
error decomposition is
\[
 E_r=u_rv_d^*+\sqrt q\,V_rH_dW^*,
 \qquad u_r\perp\operatorname{ran}V_r.
\]
In particular,
\begin{equation}
 M:=E_r^*E_r=v_dv_d^*+qWH_d^*H_dW^*.
 \label{I:B:eq:curvature-error-square}
\end{equation}
The overlap between the two embeddings is especially simple:
\begin{equation}
 W^*V_re_j=\sqrt q\,e_{j+1},\qquad
 W^*v_d=\sqrt q\,S_+\overline d.
 \label{I:B:eq:curvature-overlap}
\end{equation}
Indeed the only term in $V_re_j$ supported on a pure power of $b$ is
its $n=0$ term, at the word $b^{j+1}$.

Let $z=v_d/\sqrt A$, so $\Bnorm z=1$, and decompose $M$ relative to
$\operatorname{span}\{z\}\oplus z^\perp$. By
\eqref{I:B:eq:curvature-overlap}, its upper-left scalar entry is
\begin{equation}
 a=\langle z,Mz\rangle
   =A+q\Bnorm{H_dW^*z}^2
   =A+\frac{q^2T}{A}.
 \label{I:B:eq:curvature-rayleigh}
\end{equation}
This is a Rayleigh quotient of $M$ and proves the lower inequality in
\eqref{I:B:eq:curvature-finite-bound} for the explicit error $E_r$.

Let $B_0$ be the compression of $M$ to $z^\perp$, and let
$v=(I-zz^*)Mz$ be the off-diagonal vector. The rank-one term
$v_dv_d^*$ vanishes on $z^\perp$. Since $W$ is an isometry,
\begin{equation}
 \Bnorm{B_0}\le qL.
 \label{I:B:eq:curvature-complement}
\end{equation}
Moreover, using \eqref{I:B:eq:curvature-error-square} and
\eqref{I:B:eq:curvature-overlap},
\begin{align}
 \Bnorm v^2
 &\le q^2\Bnorm{WH_d^*H_dW^*z}^2\\
 &\le q^2L\Bnorm{H_dW^*z}^2
  =\frac{q^3LT}{A}.
 \label{I:B:eq:curvature-offdiagonal}
\end{align}
The assumed inequality $qL\le A/2$ gives
$a-\Bnorm{B_0}\ge A/2$. Applying the positive-block estimate used in
Proposition~\ref{I:B:prop:lift-k2-sharp} yields
\[
 \Bnorm M\le a+\frac{\Bnorm v^2}{a-\Bnorm{B_0}}
 \le a+\frac{2q^3LT}{A^2},
\]
which is the upper inequality in
\eqref{I:B:eq:curvature-finite-bound}. The two expansions follow immediately;
in taking the square root, the next Taylor correction from the $q^2$
term is of order $q^4$ and is absorbed by the displayed $O(q^3)$.
\end{proof}


\subsection{An autocorrelation formula for Hardy extremal lifts}

For the scalar residual constructed from a Hardy coefficient extremal,
the curvature coefficient can be written directly in terms of its outer
polynomial. This makes the cancellation and its quantitative size
explicit.

\begin{corollary}[Curvature from the outer extremal polynomial]
\label{I:B:cor:hardy-curvature}
Let $m=k-1$, let $C=C(m,2/3)$, and choose an outer coefficient-extremal
polynomial
\[
 P(z)=\sum_{j=0}^{m}p_jz^j,\qquad
 \Bnorm P_{\BHtwo}=1,\qquad [z^m]P^3=C>0.
\]
Use reversal order $m$ in
$P^*(z)=z^m\overline{P(1/\overline z)}$ and let $h$ be the scalar
approximant supplied by the Hardy construction, whose residual is
\begin{equation}
 d=z^m-h=\frac{(P^*)^2}{CP}.
 \label{I:B:eq:curvature-hardy-residual}
\end{equation}
Put $J=\Bnorm P_{H^4}^4$ and $L=\Bnorm{H_d}^2$. Then
\begin{equation}
 A=C^{-2},\qquad
 T=\frac{J-1}{2C^4},\qquad
 \frac TA=\frac{J-1}{2C^2}.
 \label{I:B:eq:curvature-hardy-coefficients}
\end{equation}
For $qL\le1/(2C^2)$, the error of the corresponding explicit word lift
satisfies
\begin{equation}
 \boxed{\displaystyle
 \frac1{C^2}+\frac{J-1}{2C^2}q^2
 \le\Bnorm{E_r}^2
 \le\frac1{C^2}+\frac{J-1}{2C^2}q^2
           +L(J-1)q^3.}
 \label{I:B:eq:curvature-hardy-finite}
\end{equation}
In particular,
\begin{equation}
 \Bnorm{E_r}=\frac1C+\frac{J-1}{4C}q^2+O(q^3),
 \qquad
 \delta_k(H_{f_{r,k}})\le E_{k-1}+O(q^2).
 \label{I:B:eq:curvature-hardy-norm}
\end{equation}
The expansion in \eqref{I:B:eq:curvature-hardy-norm} is for this chosen
lift; the final inequality is only an upper bound for the optimum.
\end{corollary}

\begin{proof}
On the unit circle, $|P^*|=|P|$. Hence
\eqref{I:B:eq:curvature-hardy-residual} gives
\begin{equation}
 |d|^2=\frac{|P|^2}{C^2}\quad\text{almost everywhere}.
 \label{I:B:eq:curvature-boundary-modulus}
\end{equation}
Any boundary zeros of $P$ cancel in the rational expression for $d$;
they do not alter this identity or its coefficient consequences.
Integrating \eqref{I:B:eq:curvature-boundary-modulus} proves $A=C^{-2}$.

The $i$th scalar coefficient of $H_dS_+\overline d$ is
\begin{equation}
 (H_dS_+\overline d)_i
   =\sum_{j\ge0}d_{i+j+1}\overline{d_j},\qquad i\ge0.
 \label{I:B:eq:curvature-autocorrelation-d}
\end{equation}
These sums converge absolutely by Cauchy--Schwarz. They are the positive
Fourier coefficients of $|d|^2$. Applying
\eqref{I:B:eq:curvature-boundary-modulus} therefore gives
\[
 (H_dS_+\overline d)_{i-1}
   =\frac1{C^2}\sum_{j=0}^{m-i}p_{j+i}\overline{p_j}
 \quad(1\le i\le m),
\]
and these coefficients vanish for $i>m$. Consequently,
\begin{equation}
 T=\frac1{C^4}\sum_{i=1}^{m}
       \left|\sum_{j=0}^{m-i}p_{j+i}\overline{p_j}\right|^2.
 \label{I:B:eq:curvature-autocorrelation-P}
\end{equation}
Parseval's identity applied to the trigonometric polynomial $|P|^2$
shows that
\[
 \Bnorm P_{H^4}^4
 =\left(\sum_{j=0}^{m}|p_j|^2\right)^2
   +2\sum_{i=1}^{m}
       \left|\sum_{j=0}^{m-i}p_{j+i}\overline{p_j}\right|^2.
\]
The squared norm of $P$ is one, proving
\eqref{I:B:eq:curvature-hardy-coefficients}. Substitute these identities into
Proposition~\ref{I:B:prop:lift-curvature} to obtain
\eqref{I:B:eq:curvature-hardy-finite} and the norm expansion.
Finally, the scalar Hardy construction has
$E_{k-1}=1/C$, and its word lift has at most $k$ states. It is therefore
an admissible competitor for $\delta_k$, proving the upper bound in
\eqref{I:B:eq:curvature-hardy-norm}.
\end{proof}


\section{Scope and principal conclusions}
Let $\Sigma$ be finite, $d=|\Sigma|\ge2$, and let
\[
H_f[u,v]=f(uv),\qquad u,v\in\Sigma^*.
\]
Every norm of a word-Hankel operator below is its operator norm on $\FHH$ with counting measure. For bounded $H_f$ define
\[
\delta_k^{\FC}(H_f)=\inf\{\Fnorm{H_f-H_g}:H_g\text{ bounded word-Hankel},\ \Frank H_g\le k\}.
\]
For real $f$, define $\delta_k^{\FR}$ by requiring real coefficients of $g$. A bare $\delta_k$ means the complex optimum. The standard realization convention is
\[
f(a_1\cdots a_t)=\alpha^*A_{a_1}\cdots A_{a_t}\beta.
\]
Joint mean-square stability means $\rho(\Phi_A)<1$ for $\Phi_A(X)=\sum_a A_aXA_a^*$. Separate letter spectral radii do not replace this hypothesis. On rational input this condition is verifiable by solving $Y-\sum_aA_a^*YA_a=I$ and checking $Y\succ0$. A coefficient transpose $H^{\mathsf T}$ means $\mathcal JH^*\mathcal J$, with $\mathcal J$ coordinatewise conjugation; it has the same norm and singular values as $H$.

The finite-parameter family problems use the binary alphabet $\Sigma=\{a,b\}$. The synthesis theorem later allows any finite alphabet. Let
\[
q=1-r^2\in(0,1),\qquad f_{r,k}\text{ as in \eqref{I:eq:binary-source}}.
\]
Write $D_k^{\mathbb F}(q)=\delta_k^{\mathbb F}(H_{f_{\sqrt{1-q},k}})^2$.
The new finite-parameter and first-correction conclusions are:
\begin{enumerate}[leftmargin=*]
\item \textbf{Exact one-state global optimum.} Put
\[
\mathcal A_q(t)=(1-q)^2t^2+(q^2+2q-1)(t-q).
\]
Then $D_1^{\FC}(q)=D_1^{\FR}(q)$ is the unique selected root
\[
\boxed{\mathcal A_q(t)^2=4q^3(t-q)(1-2q+qt),\quad
\frac{2q}{1+q}<t<1,\quad\mathcal A_q(t)>0.}
\]
There is a single real-analytic optimal value branch throughout $(0,1)$. The complete proof and all minimizers appear below.
\item \textbf{Global two-state first correction.} For each field separately,
\[
\boxed{D_2^{\mathbb F}(q)=\frac34-\frac14q+O(q^2),
\qquad \mathbb F\in\{\FR,\FC\}.}
\]
This is not an assertion that $D_2^{\FR}(q)=D_2^{\FC}(q)$ at finite $q$. It implies that their difference is nonnegative and $O(q^2)$. The remainder constant and a numerical neighborhood are not explicitly evaluated; their existence follows from the analytic proof.
\item \textbf{Quantitative additive global synthesis.} For supplied stable rational input, an explicit finite enumeration produces a stable rational rank-at-most-$k$ competitor and rational $L,U$ with
\[
0\le L\le\delta_k(H_f)\le\Fnorm{H_f-H_g}\le U,\qquad U-L\le\varepsilon.
\]
Its candidate count is explicit and exponential in $2dk^2+4k$. Uniform moment and size bounds on the input give a nontrivial class on which the scheme is polynomial in input dimension, input bit length and $1/\varepsilon$ for fixed $d,k$. It is a theoretical enumeration scheme, not a practical fast minimizer or a polynomial algorithm in all parameters.
\end{enumerate}

The finite-parameter arguments use the boundary theorem as input and establish separate optimality statements. Classical completion, finite nets, rational Stein solves and analytic implicit-function arguments supply the supporting tools.

\section{Canonical facts used in the synthesis argument}
These are background facts. We include short proofs to make the construction's comparison class explicit.

\begin{lemma}[Stable observable canonical coordinates]\label{I:F:lem:canonical-input}
Every bounded word-Hankel $G$ of finite rank $j$ has a minimal $j$-state realization $g(w)=\ell T_w\beta$ such that
\[
\sum_a T_a^*T_a+\ell^*\ell=I,\qquad \Fnorm\ell\le1,
\qquad\Fnorm\beta\le\Fnorm G.
\]
Its completely positive transition map is strictly stable. Padding to a larger dimension preserves the displayed equality as an inequality.
\end{lemma}
\begin{proof}
Let $R_ae_u=e_{ua}$ be the right creation isometries and $\mathcal S=\Fran G$. The word-Hankel identity $R_a^*G=GL_a$, with $L_ae_v=e_{av}$, makes $\mathcal S$ invariant under $R_a^*$. Put $T_a=R_a^*|_{\mathcal S}$, $\beta=Ge_\Feps$ and $\ell x=x[\Feps]$. Then $G e_w=T_w\beta$ and $g(w)=\ell T_w\beta$. Partitioning words by last letter proves the energy identity. The observation map is inclusion: $\ell T_wx=x[w]$.

For a vector in this finite-dimensional space,
\[
\sum_{|w|=t}\Fnorm{T_wx}^2=\sum_{|v|\ge t}|x[v]|^2\longrightarrow0.
\]
Summing over a finite orthonormal basis gives uniform decay of the positive transition powers and hence a strictly smaller than one spectral radius. Reachability follows from the columns of $G$ spanning $\mathcal S$; observability from inclusion. The remaining bounds are immediate.
\end{proof}


The structured optimum is attained: a norm-bounded minimizing sequence has a weak-operator convergent subsequence; the Hankel entry equalities and all rank-$(k+1)$ minors pass to the limit, and the operator norm is lower semicontinuous. This existence fact is used in the radial comparison, not claimed as a new algorithm.

For a stable finite error realization $e(w)=cD_wb$, its rational or algebraic Gramians satisfy
\[
P=bb^*+\sum_aD_aPD_a^*,\qquad
Q=c^*c+\sum_aD_a^*QD_a.
\]
The nonzero eigenvalues of $PQ$ are the squared singular values of $H_e$, including multiplicities. They are real and nonnegative, even if $PQ$ is not Hermitian. This factorization evaluates the infinite operator norm without a word-horizon truncation.

\section{The global one-state optimum at every finite parameter}
\label{I:F:sec:F1}

Let $0<q<1$, $r=\sqrt{1-q}$, and let $f=f_{r,1}$ be the binary
series $f(a^n b)=qr^n$, zero elsewhere. All norms in this section are
ordinary counting-measure word-Hankel operator norms. Write
$D_1^{\FC}(q)=\delta_1^{\FC}(H_f)^2$ and similarly for real competitors.
The result below determines the global optimum, rather than an upper
bound within a prescribed family of amplitudes.

\begin{theorem}[Exact finite-parameter solution of F1]
\label{I:F:thm:F1-exact}
For every $0<q<1$, there is a unique $u=u(q)$ in
$q<u<\sqrt q$ satisfying
\begin{equation}
 q^2(1-u^2)(u-q)^2=(1-q)^3u^2(q-u^2).
 \label{I:F:eq:F1-u-quartic}
\end{equation}
The exact real and complex optima coincide and equal
\begin{equation}
 \boxed{\displaystyle
 D_1^{\FC}(q)=D_1^{\FR}(q)=t(q)
 =\frac{q^2+(1-2q)u(q)^2}{q(1-u(q)^2)}.}
 \label{I:F:eq:F1-value}
\end{equation}
They satisfy $2q/(1+q)<t(q)<1$. The selected functions $u(q)$ and
$t(q)$ are real analytic on $(0,1)$, with no interior change of optimal
branch.

An explicit real optimizer is $g_0(w)=c a^{\#_a(w)}y^{\#_b(w)}$, where
\begin{equation}
 a=\frac{1-u}{r},\quad
 y=\sqrt{\frac{u(1-u)}{1-q}},\quad
 m=\frac{u-q}{1-q},\quad
 \tau=\frac{u}{y(1+u)},\quad c=m\tau.
 \label{I:F:eq:F1-optimizer-parameters}
\end{equation}
These parameters obey $a,y,m,c>0$ and $a^2+y^2=1-m<1$.
All complex optimizers are precisely
\begin{equation}
 g_\theta(w)=c a^{\#_a(w)}y^{\#_b(w)}
             e^{i\theta(\#_b(w)-1)},\qquad
 \theta\in\FR/(2\pi\mathbb Z).
 \label{I:F:eq:F1-all-optimizers}
\end{equation}
The real optimizers are the two cases $\theta=0,\pi$.
\end{theorem}

The equation in $u$ has degree at most four; an exceptional vanishing of
its leading coefficient does not affect the unique interval-selected
analytic branch. A quartic directly in the optimum is given after the
proof.

\subsection{Geometric rank-one competitors and an exact residual lower bound}

Every nonzero bounded rank-one word-Hankel series has a one-state
realization
\[
 g(w)=c\lambda_w,\qquad c\ne0,\qquad
 \lambda_w=\prod_{j=1}^{|w|}\lambda_{w_j},\qquad
 |\lambda_a|^2+|\lambda_b|^2<1.
\]
Indeed minimal rank-one realization matrices are scalars, and boundedness
of the empty column forces the geometric coefficients to be square
summable. Conversely the displayed strict inequality gives a bounded
rank-one Hankel. Put
\[
 x=\lambda_a,\quad z=|\lambda_b|^2,\quad
 m=1-|x|^2-z>0,\quad
 p=\sqrt m\sum_w\lambda_w e_w.
\]
Then $\Fnorm p=1$ and $H_g=(c/m)pp^{\mathsf T}$; its input row space is
$\operatorname{span}\{\overline p\}$, with coefficientwise conjugation.

Write $s=\sqrt q$ and use the real unit vectors
\[
 u_0=s\sum_{n\ge0}r^n e_{a^n},\qquad
 v=s\sum_{n\ge0}r^n e_{a^n b},\qquad e=e_\Feps.
\]
Here $u_0\perp v$, $v\perp e$, and $\langle u_0,e\rangle=s$. The
target has the exact factorization
\begin{equation}
 H_f=u_0v^*+sve^*,\qquad H_f^*H_f=vv^*+qee^*.
 \label{I:F:eq:F1-target}
\end{equation}
Define
\begin{equation}
 A=\frac{qm}{|1-rx|^2},\qquad B=Az.
 \label{I:F:eq:F1-overlaps}
\end{equation}
The squared overlaps of $p$ with $u_0,v,e$ are $A,B,m$; since $v$ is
real, the squared overlap of $\overline p$ with $v$ is also $B$.

Let $P$ project onto $\overline p$. Every amplitude therefore satisfies
\begin{equation}
 \Fnorm{H_f-H_g}^2\ge R^2:=\Fnorm{H_f(I-P)}^2.
 \label{I:F:eq:F1-input-lower}
\end{equation}
The nonzero eigenvalues determining $R^2$ are those of the two-by-two
Hermitian matrix with diagonal entries $1-B,q(1-m)$ and off-diagonal
modulus $\sqrt{qBm}$. Its trace is $1+q-B-qm$ and determinant is
$q(1-B-m)$. For every $q<t<1$, its scalar Schur complement gives
\begin{equation}
 R^2\le t
 \quad\Longleftrightarrow\quad
 B\ge(1-t)\left(1+\frac{qm}{t-q}\right).
 \label{I:F:eq:F1-right-feasibility}
\end{equation}
The complementary diagonal is $t-q+qm>0$, so the equivalence has no
singular-denominator qualification. Rank-one interlacing also gives
$R^2\ge q$.

For fixed $|x|,z,m$, the left side of
\eqref{I:F:eq:F1-right-feasibility} is maximized by taking $x=|x|\ge0$,
because $|1-rx|\ge1-r|x|$. The phase of $\lambda_b$ has no effect.
Thus optimization of this lower bound over \emph{all complex}
competitors reduces exactly to $x=a\ge0$, $z\ge0$ and
$m=1-a^2-z>0$. A matching real completion, rather than an assumed
real-parameter restriction, will establish equality of the original
complex and real optima.

\subsection{A unique scalar branch solves the residual optimization}

For $q<t<1$, set
\[
 \lambda=\frac{1-t}{t-q}>0,\qquad
 d=(1-ra)^2,\qquad c_a=1-a^2.
\]
Condition~\eqref{I:F:eq:F1-right-feasibility} becomes
\begin{equation}
 q(c_a-z)(z-\lambda d)\ge(1-t)d.
 \label{I:F:eq:F1-concave-z}
\end{equation}
Positive feasibility requires $c_a>\lambda d$. The quadratic on the
left is strictly concave in $z$ and uniquely maximized at
\begin{equation}
 z=\frac{c_a+\lambda d}{2},\qquad
 m=\frac{c_a-\lambda d}{2}.
 \label{I:F:eq:F1-z-vertex}
\end{equation}
The maximum is $q(c_a-\lambda d)^2/4$. Therefore feasibility is equivalent
to a positive value satisfying
\begin{equation}
 \frac{\sqrt q}{2}
 \max_a\left\{\frac{1-a^2}{1-ra}-\lambda(1-ra)\right\}
 \ge\sqrt{1-t}.
 \label{I:F:eq:F1-residual-scalar}
\end{equation}

Set $u=1-ra\in(1-r,1]$. The expression to maximize is
\begin{equation}
 F_\lambda(u)
 =\frac{2-(1+(1-q)\lambda)u-q/u}{1-q}.
 \label{I:F:eq:F1-concave-u}
\end{equation}
It is strictly concave, with unconstrained maximum at
\[
 u_\lambda=\sqrt{\frac{q}{1+(1-q)\lambda}}.
\]
If $u_\lambda\le1-r$, its constrained maximum would occur at $1-r$
and equal $-\lambda(1-r)<0$, precluding feasibility. Also
$u_\lambda<\sqrt q<1$. A positive maximal value is equivalent to
$0<\lambda<1/q$; in this case $q<u_\lambda<\sqrt q$, so it lies
strictly inside the allowed interval. The maximum is
\[
 \max F_\lambda=
 \frac{2-2\sqrt{q(1+(1-q)\lambda)}}{1-q}.
\]
Using $t=(1+q\lambda)/(1+\lambda)$, the threshold condition is exactly
\begin{equation}
 \frac{q}{(1-q)^2}
 \left(1-\sqrt{q(1+(1-q)\lambda)}\right)^2
 \ge\frac{(1-q)\lambda}{1+\lambda}.
 \label{I:F:eq:F1-monotone-lambda}
\end{equation}
On $(0,1/q)$ its left side strictly decreases from a positive value to
zero, while its right side strictly increases from zero to a positive
value. There is consequently exactly one equality. Since $t$ strictly
decreases with $\lambda$, that equality is the minimal feasible residual
threshold.

At equality, set $u=u_\lambda$. Then
\begin{equation}
 \lambda=\frac{q-u^2}{(1-q)u^2},\qquad
 z=\frac{u(1-u)}{1-q},\qquad
 m=\frac{u-q}{1-q},
 \label{I:F:eq:F1-stationary-parameters}
\end{equation}
and the threshold is \eqref{I:F:eq:F1-value}. Its complement has the two
equivalent expressions
\begin{equation}
 1-t=\frac{q(u-q)^2}{(1-q)^2u^2}
     =\frac{(1-q)(q-u^2)}{q(1-u^2)}.
 \label{I:F:eq:F1-equality-u}
\end{equation}
Equating them gives \eqref{I:F:eq:F1-u-quartic}, without introducing an
extraneous sign. Conversely any root in $(q,\sqrt q)$ satisfies the
positive equality. The first expression in \eqref{I:F:eq:F1-equality-u}
strictly increases from zero in $u$, while the second strictly decreases
to zero. This proves the asserted root uniqueness directly. Their
difference has positive $u$ derivative, so the implicit function theorem
also proves analyticity on $(0,1)$.

We have thus proved that \emph{every} bounded complex rank-one competitor
has squared error at least the number $t(q)$ in
\eqref{I:F:eq:F1-value}. The zero competitor has error one, and cannot be
optimal since $t(q)<1$.

\subsection{Strict output slack and a genuine Hankel completion}

Take the positive-real geometric parameters in
\eqref{I:F:eq:F1-stationary-parameters}, with $y=\sqrt z$, and let $p$ be
the corresponding real unit vector. The output projection residual
$L^2=\Fnorm{(I-pp^*)H_f}^2$ is the larger eigenvalue of
\[
 \begin{pmatrix}
 1-A&-A\sqrt{qz}\\
 -A\sqrt{qz}&q(1-Az)
 \end{pmatrix}.
\]
For $t>q$, its Schur threshold is
\begin{equation}
 L^2\le t
 \quad\Longleftrightarrow\quad
 A\ge(1-t)\left(1+\frac{qAz}{t-q}\right)
 =(1-t)+\lambda qAz.
 \label{I:F:eq:F1-left-feasibility}
\end{equation}
Equation~\eqref{I:F:eq:F1-equality-u} gives $1-t=Am$. Hence the final
condition is $1-\lambda qz\ge m$. Substituting
\eqref{I:F:eq:F1-stationary-parameters}, its strict version is equivalent to
\[
 qu^2+(1-q)u-q^2>0.
\]
This polynomial is strictly increasing on $u>0$ and at $u=q$ equals
$q(1-q)^2>0$. Thus
\begin{equation}
 \Fnorm{H_f(I-pp^*)}^2=t,\qquad
 \Fnorm{(I-pp^*)H_f}^2<t.
 \label{I:F:eq:F1-two-corners}
\end{equation}
The strict second inequality supplies the required sequentially compatible
completion. In particular, we do not assert that $H_fpp^*$ itself is
Hankel.

\begin{lemma}[Scalar one-corner norm completion]
\label{I:F:lem:F1-completion}
Let $H$ be bounded and let $p$ be a unit vector used as the selected
one-dimensional input and output space. If
\[
 \Fnorm{H(I-pp^*)}\le\gamma,
 \qquad\Fnorm{(I-pp^*)H}<\gamma,
 \quad\gamma>0,
\]
there is an amplitude $\tau$ such that
$\Fnorm{H-\tau pp^*}\le\gamma$. If the first displayed norm equals
$\gamma$, this amplitude is unique. Real data admit a real amplitude.
\end{lemma}

\begin{proof}
This is a scalar contraction-completion result; we provide its elementary
Schur proof rather than relying on an external completion theorem (compare
the classical mechanism discussed in~\cite{parrott2025}). Write
\[
 H=\begin{pmatrix}h&b\\c&D\end{pmatrix}
\]
relative to $\operatorname{span}\{p\}\oplus p^\perp$. Since
$\Fnorm D<\gamma$, the complement block
\[
 \mathcal L=\begin{pmatrix}\gamma I&D\\D^*&\gamma I\end{pmatrix}
\]
is strictly positive. Put $t=\gamma^2$ and $x=h-\tau$.
The norm inequality for $M=\left(\begin{smallmatrix}x&b\\c&D\end{smallmatrix}\right)$
is equivalent to positivity of
$\left(\begin{smallmatrix}\gamma I&M\\M^*&\gamma I\end{smallmatrix}\right)$.
Its Schur complement with respect to $\mathcal L$ is the scalar block
\[
 \begin{pmatrix}
 \gamma\{1-b(tI-D^*D)^{-1}b^*\}
 &x+bD^*(tI-DD^*)^{-1}c\\
 \overline x+c^*(tI-DD^*)^{-1}Db^*
 &\gamma\{1-c^*(tI-DD^*)^{-1}c\}
 \end{pmatrix}.
\]
The given column and row norm bounds make both diagonal entries
nonnegative. Choose $x$ to make the off-diagonal zero. This proves
existence. If the column norm equals $\gamma$, its scalar Schur
complement, the first diagonal entry, is zero; positivity then forces the
off-diagonal to vanish, proving uniqueness. All formulas preserve real
data.
\end{proof}


\subsection{Explicit amplitude and classification of the minimizers}

Apply Lemma~\ref{I:F:lem:F1-completion} with $\gamma=\sqrt t$ and
\eqref{I:F:eq:F1-two-corners}. To identify its unique amplitude, set
\[
 z_0=(tI-H_f^*H_f)^{-1}p.
\]
This inverse exists because $q<t<1$ and
$\operatorname{spec}(H_f^*H_f)=\{0,q,1\}$. The right residual secular
equation gives $p^*z_0=0$; equivalently it follows by multiplying the
characteristic equation associated with
\eqref{I:F:eq:F1-right-feasibility}. Hence
\[
 H_f^*H_fz_0=t z_0-p,
 \qquad\Fnorm{H_fz_0}^2=t\Fnorm{z_0}^2.
\]
Every completion of norm at most $\sqrt t$ preserves this saturated input
direction. Taking its scalar cross term against $p$ gives
\begin{equation}
 \tau=-\frac1{p^*H_fz_0}.
 \label{I:F:eq:F1-amplitude-necessary}
\end{equation}
All quantities are real here. Put
$\alpha=\langle p,u_0\rangle=\sqrt A$ and
$\beta=\langle p,v\rangle=\alpha y$. The target factorization gives
\[
 p^*H_fz_0
 =\frac{\alpha\beta}{t-1}
  +\frac{\sqrt q\,\beta\sqrt m}{t-q}
 =-\frac{\alpha\beta(1-u\lambda)}{1-t}.
\]
The denominator in \eqref{I:F:eq:F1-amplitude-necessary} is nonzero because
\[
 1-u\lambda=\frac{(u-q)(1+u)}{(1-q)u}
             =\frac{m(1+u)}u>0.
\]
Finally $1-t=Am$, so
\[
 \tau=\frac{1-t}{Ay(1-u\lambda)}
       =\frac{u}{y(1+u)}.
\]
The completion lemma proves that this amplitude has error at most
$\sqrt t$, and the universal residual lower bound proves equality.
Since $pp^*=m(\lambda_w)_w(\lambda_w)_w^{\mathsf T}$ for these real
parameters, this is exactly the coefficient series in
\eqref{I:F:eq:F1-optimizer-parameters}. A real competitor has thus attained
the lower bound over all complex competitors, proving
$D_1^{\FR}=D_1^{\FC}=t$.

For completeness, every optimizer must saturate the inequalities in the
residual optimization. The strict phase bound for $|1-rx|$, together with
the strict quadratic maxima, forces $\lambda_a=a>0$ and the unique
modulus $|\lambda_b|=y$. Let
$U_\theta e_w=e^{i\theta\#_b(w)}e_w$. Because every target word has
exactly one $b$,
\[
 U_\theta H_fU_\theta=e^{i\theta}H_f.
\]
Consequently
$H_{g_\theta}=e^{-i\theta}U_\theta H_{g_0}U_\theta$ satisfies
\[
 H_f-H_{g_\theta}
 =e^{-i\theta}U_\theta(H_f-H_{g_0})U_\theta.
\]
Both side factors are unitary, so the errors have the same norm.
Conversely, a phase-$\theta$ competitor is reduced to the real geometric
pair by $G\mapsto e^{i\theta}U_\theta^*GU_\theta^*$ before applying
Lemma~\ref{I:F:lem:F1-completion}. This gives
\eqref{I:F:eq:F1-all-optimizers}. Strict output slack makes
the amplitude unique for each fixed geometric pair, so these exhaust the
optimizers. Since $g_\theta(\Feps)=ce^{-i\theta}\ne0$, real
coefficients force $\theta=0$ or $\pi$; both choices are indeed real.
This completes the proof of Theorem~\ref{I:F:thm:F1-exact}.

\subsection{A quartic directly in the optimum}

\begin{theorem}[Direct polynomial branch selector]
\label{I:F:thm:F1-t-quartic}
For $0<q<1$ put
\[
 B_q=q^2+2q-1,\qquad d_q(t)=1-2q+qt,\qquad
 A_q(t)=(1-q)^2t^2+B_q(t-q).
\]
Then $D_1^{\FC}(q)=D_1^{\FR}(q)$ is the unique real $t$ satisfying
\begin{equation}
 \boxed{\displaystyle A_q(t)^2=4q^3(t-q)d_q(t),\qquad
 \frac{2q}{1+q}<t<1,\quad A_q(t)>0.}
 \label{I:F:eq:F1-direct-quartic}
\end{equation}
The optimizer is reconstructed by
$u=\sqrt{q(t-q)/d_q(t)}$ and
\eqref{I:F:eq:F1-optimizer-parameters}.
\end{theorem}

\begin{proof}
On the stated interval $d_q(t)>0$, and
\[
 u^2=\frac{q(t-q)}{d_q(t)}
\]
lies strictly between $q^2$ and $q$. This is the inverse of
\eqref{I:F:eq:F1-value}. With this substitution,
\eqref{I:F:eq:F1-u-quartic} is equivalently
\[
 (u-q)^2=
 \frac{(1-q)^2(t-q)(1-t)}{d_q(t)}.
\]
Expanding the square and using the expression for $u^2$ gives
\[
 2q\,d_q(t)u=A_q(t).
\]
Its square is \eqref{I:F:eq:F1-direct-quartic}; the strict condition
$A_q(t)>0$ chooses the positive sign and recovers the preceding equality.
Thus the interval and sign conditions make the polynomial formulation
equivalent to the unique $u$ branch already proved. No additional
quantifier-elimination or local-optimization selection is required.
\end{proof}


\subsection{Short certified intervals without global elimination}

The selected scalar branch also provides an elementary certification
procedure. Let $[u_-,u_+]$ be the bracket after $N$ bisection steps from
$[q,\sqrt q]$, using the sign of the polynomial in
\eqref{I:F:eq:F1-u-quartic} moved to the left. Its sign agrees with the
strictly increasing difference in \eqref{I:F:eq:F1-equality-u}; it is
negative at $q$ and positive at $\sqrt q$. Thus bisection is certified
without searching over realizations.

Regarding the expression in \eqref{I:F:eq:F1-value} as $t_q(u)$,
\[
 t_q'(u)=\frac{2u(1-q)^2}{q(1-u^2)^2}\le\frac2{\sqrt q}
 \qquad(q<u<\sqrt q).
\]
Since $u_+-u_-=(\sqrt q-q)2^{-N}$, monotonicity gives certified
squared-error bounds
\[
 t_q(u_-)\le D_1(q)\le t_q(u_+),\qquad
 t_q(u_+)-t_q(u_-)\le2(1-\sqrt q)2^{-N}\le2^{1-N}.
\]
Their positive square roots bound $\delta_1$, with gap at most
$2^{(1-N)/2}$. In particular $N\ge1+2\log_2(1/\epsilon)$ suffices
for a norm interval of width at most $\epsilon$. This is a uniform
count of degree-four sign evaluations, rather than a quantifier-elimination
procedure.

For rational $q$ of binary length $B$, the bisection endpoints lie in the
quadratic field $\mathbb Q(\sqrt q)$. After clearing denominators, each
polynomial sign is the sign of $A+B'\sqrt{pd}$ with integers of bit length
$O(B+N)$ when $q=p/d$ in lowest terms. Signs reduce to integer signs and,
when necessary, comparison of $A^2$ with $(B')^2pd$. Consequently elementary
schoolbook arithmetic gives a polynomial bit bound, for example
$O(N(B+N)^2)$ with a common-denominator representation. The exact optimizer
itself is specified by the degree-at-most-four selected root and the
explicit radicals in \eqref{I:F:eq:F1-optimizer-parameters}; this avoids any
numerical perturbation of the optimal realization.

\section{A two-state branch that improves the first boundary correction}
\label{I:F:sec:f2}

The global two-state first correction is a different problem from the
one-state optimum. The previously available fixed lift has squared error
$3/4+q^2/6+O(q^3)$. The construction below proves that allowing the
competitor to depend on $q$ improves the error already at order $q$.
The matching global lower bound is proved separately by the two-dimensional passive-envelope argument.

\begin{theorem}[An improving real two-state branch]
\label{I:F:thm:f2-improvement}
For $0<q<1/2$, set
\[
 r=\sqrt{1-q},\qquad a=\sqrt{1-q-2q^2},\qquad
 c=\frac{q(1+2q)}2,
\]
and define the standard-order realization
\begin{equation}
 A_a=\begin{pmatrix}a&0\\0&0\end{pmatrix},\qquad
 A_b=\frac1{\sqrt2}\begin{pmatrix}q&1\\-q&1\end{pmatrix},\qquad
 \alpha_g^*=c(1,0),\qquad \beta_g=(-q,1)^{\mathsf T}.
 \label{I:F:eq:f2-branch}
\end{equation}
Let $g_q(w)=\alpha_g^*A_w\beta_g$, with zero transition matrices for any
additional alphabet letters. This is a real, jointly stable realization
of dimension two. For the target $f_{r,2}(a^n b^2)=qr^n$,
\begin{equation}
 \boxed{\displaystyle
 \Fnorm{H_{f_{r,2}}-H_{g_q}}^2
 =\frac34-\frac14q+O(q^2)\qquad(q\downarrow0).}
 \label{I:F:eq:f2-branch-expansion}
\end{equation}
Consequently
\begin{equation}
 D_2^{\FC}(q)\le D_2^{\FR}(q)
 \le\frac34-\frac14q+O(q^2).
 \label{I:F:eq:f2-optimum-upper}
\end{equation}
In particular the fixed lift approaching $3/4$ from above is not
asymptotically optimal at first order.
\end{theorem}

\begin{proof}
We establish joint stability, give exact error Gramians, and differentiate
the isolated leading squared singular value.

\emph{Joint stability.}
Direct multiplication yields
\begin{equation}
 A_a^*A_a+A_b^*A_b
   =\operatorname{diag}(1-q-q^2,1).
 \label{I:F:eq:f2-passivity}
\end{equation}
Let $D_0=\operatorname{diag}(q+q^2,0)$. The matrix
$D_0+A_a^*D_0A_a+A_b^*D_0A_b$ is positive definite. Indeed, its kernel
would be contained in $\ker D_0=\operatorname{span}\{(0,1)^{\mathsf T}\}$, but
$A_b(0,1)^{\mathsf T}=(1,1)^{\mathsf T}/\sqrt2$ has nonzero first coordinate.
Telescoping \eqref{I:F:eq:f2-passivity} therefore gives
\[
 \sum_{|w|=2}A_w^*A_w\prec I.
\]
The column completely positive map is a strict contraction after two
steps, and its adjoint has the same spectral radius. Thus the system is
mean-square stable. Its word-Hankel is bounded of rank at most two.

\emph{Competitor Gramians.}
Let $R_g,Q_g$ be its exact reachability and observation Gramians. Their
Stein equations have the solutions
\begin{equation}
 R_g=\begin{pmatrix}
 \dfrac{1+q^2}{q(1+q)}&\dfrac{1-q}{1+q}\\[4pt]
 \dfrac{1-q}{1+q}&\dfrac{2+q+2q^2+q^3}{1+q}
 \end{pmatrix},\qquad
 Q_g=\frac{q(1+2q)^2}{4(1+q)}I_2.
 \label{I:F:eq:f2-own-gramians}
\end{equation}
These identities can be verified by substitution in
$R_g=\beta_g\beta_g^*+\sum_j A_jR_gA_j^*$ and
$Q_g=\alpha_g\alpha_g^*+\sum_j A_j^*Q_gA_j$.
For example, writing the entries of $R_g$ as $u,v,w$, the off-diagonal
and lower-right equations give
$v=(1-q)/(1+q)$ and $w=2+q^2u-2qv$; the remaining equation gives
$u=(1+q^2)/(q(1+q))$. Joint stability ensures uniqueness.

\emph{Target and cross Gramians.}
Use the target realization
\[
 F_a=\operatorname{diag}(r,0,0),\qquad
 F_b=\begin{pmatrix}0&1&0\\0&0&1\\0&0&0\end{pmatrix},\qquad
 \alpha_f^*=q(1,0,0),\qquad \beta_f=(0,0,1)^{\mathsf T}.
\]
Its Gramians are
$R_f=\operatorname{diag}(q^{-1},1,1)$ and $Q_f=qI_3$.
Put $\Delta=1-ra>0$. Define $X,Y\in\FC^{3\times2}$ by their rows:
\begin{align}
 X_0&=\left(\frac{1+q+q^2-q^3}{2\Delta},
                 \frac{1-q+q^2+q^3}{2}\right),\notag\\
 X_1&=\frac1{\sqrt2}(1-q^2,1+q^2),\qquad X_2=(-q,1),
 \label{I:F:eq:f2-cross-reach}\\
 Y_0&=\frac{qc}{\Delta}(1,0),\qquad
 Y_1=Y_0A_b,\qquad Y_2=Y_0A_b^2.
 \label{I:F:eq:f2-cross-observe}
\end{align}
The triangular form of $F_b$ gives directly
\[
 X=\beta_f\beta_g^*+\sum_jF_jXA_j^*,\qquad
 Y=\alpha_f\alpha_g^*+\sum_jF_j^*YA_j.
\]
For example, the first equation gives successively
$X_2=\beta_g^*$, $X_1=\beta_g^*A_b^*$, and
$X_0=\beta_g^*(A_b^*)^2\operatorname{diag}(\Delta^{-1},1)$,
which is \eqref{I:F:eq:f2-cross-reach}.

For the composite error realization with output
$(\alpha_f^*,-\alpha_g^*)$ and input $(\beta_f,\beta_g)$, the exact
Gramians are therefore
\begin{equation}
 R=\begin{pmatrix}R_f&X\\X^*&R_g\end{pmatrix},\qquad
 Q=\begin{pmatrix}Q_f&-Y\\-Y^*&Q_g\end{pmatrix}.
 \label{I:F:eq:f2-error-gramians}
\end{equation}
Its infinite Hankel squared norm is the largest eigenvalue of
$R^{1/2}QR^{1/2}$, equivalently the largest eigenvalue of $RQ$.
Equations \eqref{I:F:eq:f2-own-gramians}--\eqref{I:F:eq:f2-error-gramians}
are thus a finite exact algebraic description of this chosen error;
there is no word-horizon cutoff.

\emph{Leading eigenvalue perturbation.}
Set $\widehat R=qR$ and $\widehat Q=Q/q$, so
$\widehat R\widehat Q=RQ$. Both extend analytically to $q=0$.
In particular $\Delta=q+q^2+O(q^4)$, removing the apparent
singularities in $qX$ and $Y/q$.
Index the five coordinates by $(f_0,f_1,f_2,g_0,g_1)$ and put
$y=1/\sqrt2$. The expansions are
\[
 \widehat R=R_0+qR_1+O(q^2),\qquad
 \widehat Q=Q_0+qQ_1+O(q^2),
\]
where
\begin{equation}
 R_0=\begin{pmatrix}
 1&0&0&1/2&0\\0&0&0&0&0\\0&0&0&0&0\\
 1/2&0&0&1&0\\0&0&0&0&0
 \end{pmatrix},\quad
 R_1=\begin{pmatrix}
 0&0&0&0&1/2\\0&1&0&y&y\\0&0&1&0&1\\
 0&y&0&-1&1\\1/2&y&1&1&2
 \end{pmatrix},
 \label{I:F:eq:f2-R-expansion}
\end{equation}
\begin{equation}
 Q_0=\begin{pmatrix}
 1&0&0&-1/2&0\\0&1&0&0&-y/2\\0&0&1&0&-1/4\\
 -1/2&0&0&1/4&0\\0&-y/2&-1/4&0&1/4
 \end{pmatrix},\quad
 Q_1=\begin{pmatrix}
 0&0&0&-1/2&0\\0&0&0&-y/2&-y/2\\0&0&0&1/4&-1/2\\
 -1/2&-y/2&1/4&3/4&0\\0&-y/2&-1/2&0&3/4
 \end{pmatrix}.
 \label{I:F:eq:f2-Q-expansion}
\end{equation}
The matrix $Z_0=R_0Q_0$ has only one nonzero row:
$Z_0[0,0]=3/4$, $Z_0[0,3]=-3/8$.
Its sole nonzero eigenvalue $3/4$ is simple. Right and left eigenvectors,
normalized by $w^*v=1$, are
\[
 v=e_0,\qquad w=e_0-\tfrac12 e_3.
\]
The analytic eigenvalue branch consequently has derivative
\begin{equation}
 \lambda'(0)=w^*(R_1Q_0+R_0Q_1)v=-\frac14.
 \label{I:F:eq:f2-eigen-derivative}
\end{equation}
More explicitly, the two terms are $-1/4$ and zero, respectively.
All other eigenvalues tend to zero, so this branch is the largest
squared Hankel singular value for all sufficiently small positive $q$.
Analyticity supplies a bounded $O(q^2)$ remainder and proves
\eqref{I:F:eq:f2-branch-expansion}. Admissibility of the real two-state
construction proves \eqref{I:F:eq:f2-optimum-upper}.
\end{proof}


\begin{remark}[What this does and does not determine]
The theorem refutes first-order optimality of the earlier fixed lift,
and refutes any proposed lower bound of the form
$D_2(q)\ge3/4+o(q)$ near zero. Taken alone it is an upper-bound result. The separate passive-envelope
lower bound, proved in this manuscript, establishes
$D_2^{\FC}(q)\ge3/4-q/4-O(q^2)$ against all complex competitors.
Combining the two proves the common first correction for the real and
complex optima. Local numerical searches played no role in the global
lower-bound certificate; higher coefficients and exact finite-$q$
rank-two optima are not determined here.
\end{remark}

\section{The global complex two-state first correction}
\label{I:F:sec:f2-global-lower}

We now match the real upper branch of
Theorem~\ref{I:F:thm:f2-improvement} by a lower bound against every bounded
complex word-Hankel competitor of rank at most two.  The central step is
an asymptotic maximization over all passive complex two-dimensional
systems.  The maximization is global; no restriction to a chosen real
matrix family is imposed.  The proof does not establish an exact
finite-$q$ classification of the optimum.

\subsection{From arbitrary two-state competitors to a passive matrix problem}

Write $q=1-r^2$, with $0<q<1$, and let
\begin{equation}
 u_r=\sqrt q\sum_{n\ge0}r^n e_{a^n},\qquad
 v_r=\sqrt q\sum_{n\ge0}r^n e_{a^n b^2}.
 \label{I:F:eq:f2-test-vectors}
\end{equation}
Both vectors have norm one in $\FHH$.  For the target
$X_f=H_{f_{r,2}}^{\mathsf T}$, direct substitution gives
$X_fu_r=v_r$.

Let $H_g$ be any bounded complex word-Hankel operator of rank at most
two, and let $\mathcal S=\operatorname{ran}H_g^{\mathsf T}$.
This is a finite-dimensional backward-shift invariant subspace of
$\FHH$.  In orthonormal canonical state coordinates, denote the
restrictions of $L_a^*$ and $L_b^*$ by $A$ and $B$, and denote
empty-word evaluation by the row $\ell$.  If necessary, pad the system
by zero coordinates to dimension two.  The canonical energy identity
implies
\begin{equation}
 A^*A+B^*B+\ell^*\ell\preceq I_2.
 \label{I:F:eq:f2-passive-class}
\end{equation}
Additional alphabet letters only add positive terms to the omitted part
of this identity.  Coordinate evaluation on $a^n b^2$ is the row
$\ell B^2A^n$.  Consequently, if $P_{\mathcal S}$ is orthogonal
projection onto $\mathcal S$,
\begin{equation}
 \Fnorm{P_{\mathcal S}v_r}^2
   =q\Fnorm{\ell B^2(I-rA)^{-1}}^2.
 \label{I:F:eq:f2-projection-passive}
\end{equation}
The resolvent exists because \eqref{I:F:eq:f2-passive-class} implies
$\Fnorm{A}\le1$ and $r<1$.  Since the range of $H_g^{\mathsf T}$ is
contained in $\mathcal S$,
\begin{equation}
 \Fnorm{H_{f_{r,2}}-H_g}^2
 \ge \Fnorm{(I-P_{\mathcal S})v_r}^2
 =1-q\Fnorm{\ell B^2(I-rA)^{-1}}^2.
 \label{I:F:eq:f2-competitor-lower}
\end{equation}
Thus it suffices to understand the universal passive-tuple maximum
\begin{equation}
 \mathcal P(q)=
 \max_{A^*A+B^*B+\ell^*\ell\preceq I_2}
 q\Fnorm{\ell B^2(I-rA)^{-1}}^2,
 \qquad r=\sqrt{1-q}.
 \label{I:F:eq:f2-P-definition}
\end{equation}
The maximum exists: the feasible set is compact and the resolvent is
continuous there for every fixed $r<1$.

\subsection{Exact elimination of the waiting transition}

The variable $A$ can be eliminated without relaxation.  For $0\le t\le1$,
define the compact envelope
\begin{equation}
 h(t)=\max\left\{
  |\ell B^2e_1|^2:
  V=\begin{pmatrix}B\\\ell\end{pmatrix}\in\FC^{3\times2},\quad
  V^*V\preceq I_2,\quad\Fnorm{Ve_1}^2=t
 \right\}.
 \label{I:F:eq:f2-h-definition}
\end{equation}
Here and below $e_1=(1,0)^{\mathsf T}$ is a state coordinate vector.

\begin{lemma}[Exact scalar reduction]
\label{I:F:lem:f2-A-elimination}
For every $0<q<1$,
\begin{equation}
 \mathcal P(q)=
 \max_{0\le t\le1}
 \frac{q\,h(t)}{\bigl(1-\sqrt{1-q}\sqrt{1-t}\bigr)^2}.
 \label{I:F:eq:f2-scalar-reduction}
\end{equation}
Moreover $0\le h(t)\le t$.
\end{lemma}

\begin{proof}
For fixed $A,B,\ell$, change variables in the row operator norm by
$x=(I-rA)y$.  This gives
\[
 \Fnorm{\ell B^2(I-rA)^{-1}}^2
 =\sup_{y\ne0}\frac{|\ell B^2y|^2}{\Fnorm{(I-rA)y}^2}.
\]
Normalize $\Fnorm{y}=1$ and set
$D=I-B^*B-\ell^*\ell\succeq0$ and
$t=\Fnorm{By}^2+|\ell y|^2=1-y^*Dy$.
The constraint $A^*A\preceq D$ yields
\begin{equation}
 \Fnorm{(I-rA)y}\ge1-r\Fnorm{Ay}
                 \ge1-r\sqrt{1-t}.
 \label{I:F:eq:f2-denominator-lower}
\end{equation}
A unitary state change sends $y$ to $e_1$ and preserves all constraints.
Taking suprema gives the upper direction in
\eqref{I:F:eq:f2-scalar-reduction}.

Conversely, fix $B,\ell,y$ satisfying the stated conditions.  If $t<1$,
define the rank-one matrix
\begin{equation}
 A=\frac{y\,y^*D}{\sqrt{1-t}}.
 \label{I:F:eq:f2-A-completion}
\end{equation}
Then $Ay=\sqrt{1-t}\,y$ and
\[
 A^*A=\frac{Dyy^*D}{y^*Dy}\preceq D.
\]
The last inequality is Cauchy--Schwarz in the positive semidefinite
$D$ metric.  Equality holds throughout
\eqref{I:F:eq:f2-denominator-lower}.  If $t=1$, take $A=0$ instead.
This proves attainability of the scalar denominator and hence
\eqref{I:F:eq:f2-scalar-reduction}.
Finally $\Fnorm{B},\Fnorm{\ell}\le1$, so
$|\ell B^2e_1|\le\Fnorm{Be_1}\le\sqrt t$.
\end{proof}


\subsection{The global compact envelope near zero}

\begin{lemma}[Two-term expansion of the complex envelope]
\label{I:F:lem:f2-h-expansion}
The envelope \eqref{I:F:eq:f2-h-definition} satisfies
\begin{equation}
 h(t)=\frac14t+\frac14t^2+O(t^3)
 \qquad(t\downarrow0).
 \label{I:F:eq:f2-h-expansion}
\end{equation}
The remainder applies to the global maximum over all complex
contractions in \eqref{I:F:eq:f2-h-definition}.
\end{lemma}

\begin{proof}
We first parametrize the entire feasible set, then remove the phase
symmetries at its limiting maximizing set.  The resulting maximum is a
nondegenerate analytic critical point in seven real variables.

\emph{An exact compact parametrization.}
Put $t=s^2$, initially allowing either sign of the real parameter $s$.
Write the first column of $V$ as $u=sz$, where
$z\in\FC^3$ and $\Fnorm{z}=1$, and write the second column as $v$.
For $|s|<1$, the Schur complement of $I-V^*V$ gives the exact equivalence
\begin{equation}
 V^*V\preceq I
 \quad\Longleftrightarrow\quad
 \Fnorm{v}^2+\frac{s^2}{1-s^2}|z^*v|^2\le1.
 \label{I:F:eq:f2-column-ellipsoid}
\end{equation}
Thus every feasible second column is uniquely of the form
\begin{equation}
 v=W_s(z)\nu,\qquad \Fnorm{\nu}\le1,\qquad
 W_s(z)=I_3-\bigl(1-\sqrt{1-s^2}\bigr)zz^*.
 \label{I:F:eq:f2-W-parametrization}
\end{equation}
The matrix $W_s(z)$ is analytic and even in $s$ near zero.  In these
coordinates,
\[
 B=\begin{pmatrix}sz_1&v_1\\sz_2&v_2\end{pmatrix},\qquad
 \ell=(sz_3,v_3),
\]
and direct multiplication gives
\begin{align}
 \ell B^2e_1&=s\,\mathcal J(s,z,\nu),\notag\\
 \mathcal J(s,z,\nu)
 &=z_2v_2v_3+s z_2(z_3v_1+z_1v_3)+s^2z_3z_1^2.
 \label{I:F:eq:f2-J-objective}
\end{align}
Consequently
\begin{equation}
 h(s^2)=s^2\mathcal M(s)^2,\qquad
 \mathcal M(s)=\max_{\Fnorm{z}=1,\,\Fnorm{\nu}\le1}
                  |\mathcal J(s,z,\nu)|.
 \label{I:F:eq:f2-M-definition}
\end{equation}

At $s=0$,
\[
 |\mathcal J(0,z,\nu)|=|z_2\nu_2\nu_3|\le\frac12.
\]
Equality is possible exactly on the phase torus
\begin{equation}
 z_1=z_3=\nu_1=0,\qquad
 |z_2|=1,\qquad |\nu_2|=|\nu_3|=2^{-1/2}.
 \label{I:F:eq:f2-maximizing-torus}
\end{equation}
The parameter space is compact, and the objective converges uniformly
as $s\to0$.  Hence every global maximizing pair for sufficiently small
$|s|$ lies arbitrarily close to this torus.

\emph{All complex phase degeneracies can be removed.}
The original optimization is invariant under
\begin{equation}
 B\longmapsto e^{i\theta}U^*BU,\qquad
 \ell\longmapsto e^{i\phi}\ell U,\qquad
 U=\operatorname{diag}(1,e^{i\psi}).
 \label{I:F:eq:f2-phase-action}
\end{equation}
The scalar $\ell B^2e_1$ is multiplied only by
$e^{i(\phi+2\theta)}$.  On the compact parameters, the same action is
\[
 z\longmapsto Dz,\qquad
 \nu\longmapsto e^{i\psi}D\nu,\qquad
 D=\operatorname{diag}\bigl(e^{i\theta},e^{i(\theta-\psi)},e^{i\phi}\bigr).
\]
Indeed \eqref{I:F:eq:f2-W-parametrization} is equivariant under diagonal
unitaries.  The phases of $(z_2,\nu_2,\nu_3)$ change respectively by
$(\theta-\psi,\theta,\phi+\psi)$, an invertible linear transformation
of the three phase parameters.  Since those three coordinates are
nonzero near \eqref{I:F:eq:f2-maximizing-torus}, they may all be fixed to be
positive real in a smooth local gauge.  No transverse complex direction
is excluded by this choice.

A global maximizer also satisfies $\Fnorm{\nu}=1$ for all sufficiently
small $|s|$.  To see this, use radial coordinates
$\nu=\rho\widehat\nu$, $\Fnorm{\widehat\nu}=1$.
At the limiting torus the derivative of
$|\mathcal J|$ in the increasing radial direction is $1$ at $\rho=1$.
It remains positive in a neighborhood of that compact torus.  Global
maximizers approach the torus, so a maximizing point with $\rho<1$
would admit a feasible radial increase of its objective.

\emph{Seven real coordinates and a strictly negative Hessian.}
After fixing the phases and the radial constraint, choose local
coordinates
\begin{equation}
 \begin{gathered}
 z_1,z_3,\nu_1\in\FC,\qquad \xi\in\FR,\qquad
 z_2=\sqrt{1-|z_1|^2-|z_3|^2},\\
 \nu_2=\sqrt{1-|\nu_1|^2}\cos(\pi/4+\xi),\qquad
 \nu_3=\sqrt{1-|\nu_1|^2}\sin(\pi/4+\xi).
 \end{gathered}
 \label{I:F:eq:f2-seven-coordinates}
\end{equation}
Let $x\in\FR^7$ list, in order,
$\Re z_1,\Im z_1,\Re z_3,\Im z_3,\Re\nu_1,\Im\nu_1,\xi$.
The gauge sends the entire limiting maximizing torus to $x=0$.
The objective
$F(s,x)=|\mathcal J(s,z(x),\nu(x))|$ is real analytic near $(0,0)$,
because $\mathcal J(0,0)=1/2\ne0$.
At $s=0$ its expansion is
\begin{equation}
 F(0,x)=\frac12-\frac{|z_1|^2+|z_3|^2}{4}
                  -\frac{|\nu_1|^2}{2}-\xi^2+O(\Fnorm{x}^4).
 \label{I:F:eq:f2-transverse-Hessian}
\end{equation}
Thus
\begin{equation}
 \partial_x^2F(0,0)=
 \operatorname{diag}\!\left(-\frac12,-\frac12,-\frac12,-\frac12,
                           -1,-1,-2\right)\prec0.
 \label{I:F:eq:f2-Hessian-matrix}
\end{equation}
The full expansion through total degree two is
\begin{align}
 F(s,x)=\frac12-\frac14s^2+\frac{s\,\Re z_1}{\sqrt2}
          -\frac{|z_1|^2+|z_3|^2}{4}
          -\frac{|\nu_1|^2}{2}-\xi^2
          +O\bigl((|s|+\Fnorm{x})^3\bigr).
 \label{I:F:eq:f2-full-quadratic}
\end{align}
For completeness, the pure $s^2$ term follows from the exact identity
$\mathcal J(s,0)=\sqrt{1-s^2}/2$.  In
\eqref{I:F:eq:f2-J-objective}, the only mixed term of total degree two
contributing to the modulus is $s\Re z_1/\sqrt2$.

The analytic implicit function theorem applied to
$\partial_xF(s,x)=0$ gives a unique analytic critical branch $x=x(s)$
near zero.  The negative definite Hessian persists in a fixed small
coordinate neighborhood, where the objective is strictly concave.
Every global maximizer lies in that neighborhood after the phase gauge,
by the compact localization already proved.  It must therefore be this
critical branch.  This justifies applying the local analytic expansion
to the \emph{global} maximum $\mathcal M(s)$.
Differentiating the critical equations using
\eqref{I:F:eq:f2-Hessian-matrix}--\eqref{I:F:eq:f2-full-quadratic} gives
\begin{equation}
 x'(0)=(\sqrt2,0,0,0,0,0,0).
 \label{I:F:eq:f2-critical-derivative}
\end{equation}
Substituting into \eqref{I:F:eq:f2-full-quadratic} yields
\begin{equation}
 \mathcal M(s)=\frac12+\frac14s^2+O(s^3).
 \label{I:F:eq:f2-M-preliminary}
\end{equation}

\emph{Physical evenness improves the remainder.}
The feasible physical matrices with first-column squared norm $s^2$
are identical for $s$ and $-s$.  Equivalently, replacing
$(s,z,\nu)$ by $(-s,-z,\nu)$ preserves $v$ and changes
$\mathcal J$ only by its sign.  Therefore
$\mathcal M(-s)=\mathcal M(s)$.  We have proved that the maximizing
value is analytic near zero, so all its odd Taylor coefficients vanish.
Equation \eqref{I:F:eq:f2-M-preliminary} sharpens to
\begin{equation}
 \mathcal M(s)=\frac12+\frac14s^2+O(s^4).
 \label{I:F:eq:f2-M-even}
\end{equation}
Finally \eqref{I:F:eq:f2-M-definition} gives
$h(t)=t(1/2+t/4+O(t^2))^2$, which is
\eqref{I:F:eq:f2-h-expansion}.
\end{proof}


\subsection{Uniform scalar localization and the matching lower bound}

\begin{theorem}[Global passive-tuple asymptotic]
\label{I:F:thm:f2-passive-asymptotic}
The maximum over all complex passive tuples in
\eqref{I:F:eq:f2-P-definition} satisfies
\begin{equation}
 \mathcal P(q)=\frac14+\frac14q+O(q^2)
 \qquad(q\downarrow0).
 \label{I:F:eq:f2-P-expansion}
\end{equation}
\end{theorem}

\begin{proof}
At the trial point $t=q$, the denominator in
\eqref{I:F:eq:f2-scalar-reduction} equals $q^2$ exactly.  Lemma~
\ref{I:F:lem:f2-h-expansion} therefore gives the lower bound
\begin{equation}
 \mathcal P(q)\ge\frac{h(q)}q
              =\frac14+\frac14q+O(q^2).
 \label{I:F:eq:f2-P-trial}
\end{equation}
We next localize every global scalar maximizer.  For $0\le q,t\le1$,
\[
 1-\sqrt{(1-q)(1-t)}
 \ge\frac{q+t-qt}{2}\ge\frac{q+t}{4}.
\]
Together with $h(t)\le t$, this gives, for $x=t/q$,
\begin{equation}
 \frac{q h(t)}{(1-\sqrt{(1-q)(1-t)})^2}
 \le\frac{16qt}{(q+t)^2}
 =\frac{16x}{(1+x)^2}.
 \label{I:F:eq:f2-scalar-coarse-localization}
\end{equation}
For $x\notin I=[1/100,100]$ the last expression is less than $1/5$.
The trial value \eqref{I:F:eq:f2-P-trial} exceeds $1/5$ for all sufficiently
small $q$.  Thus every global maximizing point has $x\in I$.
In particular $t=qx\to0$ uniformly on all possible maximizers, which
permits the global envelope expansion to be used uniformly.

On the fixed compact interval $I$, Taylor expansion gives
\begin{align}
 1-\sqrt{(1-q)(1-qx)}
 &=\frac q2(1+x)+\frac{q^2}{8}(x-1)^2+O_I(q^3),\notag\\
 \frac{q h(qx)}{(1-\sqrt{(1-q)(1-qx)})^2}
 &=f_0(x)+qf_1(x)+O_I(q^2),
 \label{I:F:eq:f2-uniform-scalar-expansion}
\end{align}
where
\begin{equation}
 f_0(x)=\frac{x}{(1+x)^2},\qquad
 f_1(x)=\frac{x}{(1+x)^2}
       \left(x-\frac{(x-1)^2}{2(1+x)}\right).
 \label{I:F:eq:f2-f0-f1}
\end{equation}
The remainder is uniform in $x\in I$: the numerator expansion comes
from \eqref{I:F:eq:f2-h-expansion}, and the leading denominator is bounded
away from zero after division by $q^2$.

The exact identity
\[
 \frac14-f_0(x)=\frac{(x-1)^2}{4(1+x)^2}
\]
gives a quadratic maximum at $x=1$.  On $I$ there is $c_I>0$ such that
$f_0(x)\le1/4-c_I(x-1)^2$.  Also $f_1(1)=1/4$ and
$|f_1(x)-f_1(1)|\le L_I|x-1|$ for a finite Lipschitz constant $L_I$.
Completing the square yields uniformly on $I$
\begin{align*}
 f_0(x)+qf_1(x)
 &\le\frac14+\frac q4-c_I(x-1)^2+qL_I|x-1|\\
 &\le\frac14+\frac q4+\frac{L_I^2}{4c_I}q^2.
\end{align*}
Combining this with \eqref{I:F:eq:f2-uniform-scalar-expansion} gives the
matching upper bound.  Equation \eqref{I:F:eq:f2-P-trial} then proves
\eqref{I:F:eq:f2-P-expansion}.
\end{proof}


\begin{theorem}[The two-state optimum to first boundary order]
\label{I:F:thm:f2-global-first-correction}
For each coefficient field $\mathbb F\in\{\FC,\FR\}$, the squared
structured two-state optimum for $f_{\sqrt{1-q},2}$ satisfies
\begin{equation}
 \boxed{\displaystyle
 D_2^{\mathbb F}(q)=\frac34-\frac14q+O(q^2)
 \qquad(q\downarrow0).}
 \label{I:F:eq:f2-global-first-correction}
\end{equation}
In particular
$0\le D_2^{\FR}(q)-D_2^{\FC}(q)=O(q^2)$.
This asserts equality of the displayed asymptotic coefficients, not
finite-$q$ equality of the real and complex optima.
\end{theorem}

\begin{proof}
Equations \eqref{I:F:eq:f2-competitor-lower} and
\eqref{I:F:eq:f2-P-expansion} apply to every bounded complex rank-at-most-two
word-Hankel competitor.  Taking the infimum yields
\[
 D_2^{\FC}(q)\ge\frac34-\frac14q-O(q^2).
\]
Since the complex competitor class contains the real class,
$D_2^{\FC}(q)\le D_2^{\FR}(q)$.  The independent real construction of
Theorem~\ref{I:F:thm:f2-improvement}, specifically
\eqref{I:F:eq:f2-optimum-upper}, gives
$D_2^{\FR}(q)\le3/4-q/4+O(q^2)$.  These inequalities squeeze each field's
optimum to \eqref{I:F:eq:f2-global-first-correction} and give the stated
$O(q^2)$ bound on their possible gap.
\end{proof}


The analytic proof supplies a neighborhood and bounded remainder
constants, but no numerical threshold $q_0$ or explicit numerical
remainder constant is provided here.  In particular,
\eqref{I:F:eq:f2-P-expansion} does not prove the stronger proposed exact
inequality $4r^2q\Fnorm{\ell B^2(I-rA)^{-1}}^2\le1$ for every sufficiently
large $r<1$.  Nor does \eqref{I:F:eq:f2-global-first-correction} determine the
full function $D_2^{\mathbb F}(q)$ away from the boundary, its second
correction coefficient, or exact real--complex equality at positive
$q$.

\section{A quantitative global synthesis scheme by radial damping}
\label{I:F:sec:g1-damping}

This section gives an explicit solution of the additive-certification task G1
for supplied stable realizations over the complex rationals. It uses a finite
net, rather than quantifier elimination. Its cost is exponential in the
reduced parameter count, and its certificate ledger may have the same size.
Neither a new short optimality certificate nor publication priority is claimed.

Let
\[
 f(w)=cA_wb,\qquad \Phi_A(X)=\sum_{a\in\Sigma}A_aXA_a^*,
 \qquad \rho(\Phi_A)<1,
\]
be an $n$-state input realization with complex rational entries. The norm
and comparison class are the ordinary counting-measure word-Hankel norm
and all bounded complex word-Hankel competitors of rank at most $k$.
Define the input moment
\begin{equation}
 J_2(f)^2=\sum_{w\in\Sigma^*}(|w|+1)|w|^2|f(w)|^2.
 \label{I:F:eq:g1-input-moment}
\end{equation}
This is finite under the supplied joint stability assumption.

\begin{theorem}[Certified additive synthesis from complete rational input]
\label{I:F:thm:g1-net}
Fix $1\le k<n$ and rational $\varepsilon>0$. From the input one can compute
rational bounds $M\ge\max\{1,\Fnorm{H_f}\}$ and
$J\ge\max\{1,J_2(f)\}$. Put
\begin{align}
 e&=\min\{\varepsilon,1\},& B&=2M,&
 \gamma&=1-\frac{e}{3J},& \mu&=1-\gamma^2,\nonumber\\
 \eta&=\frac{e\mu^2}{15B},&
 v&=\left\lceil\sqrt{2dk^2}\right\rceil,&
 p_0&=2dk^2+4k.
 \label{I:F:eq:g1-parameters}
\end{align}
An explicit enumeration of at most
\begin{equation}
 N\le\left(3+\frac{2v}{\eta}\right)^{p_0}
   \le\left(3+\frac{240vMJ^2}{e^3}\right)^{p_0}
 \label{I:F:eq:g1-net-size}
\end{equation}
stable rational $k$-state candidates returns a competitor $g$ and rational
$L,U$ satisfying
\begin{equation}
 0\le L\le\delta_k(H_f)\le\Fnorm{H_f-H_g}\le U,
 \qquad U-L\le\varepsilon.
 \label{I:F:eq:g1-global-gap}
\end{equation}
Every candidate has a joint column-contraction margin $\mu$. Let
\[
 Z=M+B/\mu,\qquad
 p_{\rm tr}=\left\lceil
 \frac{18Z\max\{1,\lceil\log_2(n+k)\rceil\}}e\right\rceil.
\]
Construction and independent verification cost
$N\,\operatorname{poly}(n+k,d,H,p_{\rm tr},\log(Z/e))$ bit operations,
where $H$ bounds the rational input and grid coefficient heights after a
common denominator is chosen. A complete interval ledger has the same
size bound. Thus for fixed $d,k$ and polynomially bounded numerical $M,J$,
this is a polynomial-time additive approximation scheme in input length,
$n$, and $1/e$. No polynomial dependence on $d,k$ or $\log(1/e)$ is asserted.
\end{theorem}

The parameter restriction applies to a known nearoptimal damped competitor,
not to the original optimum. The proof below quantifies its loss using the
input alone, thereby avoiding a restricted-comparison-class gap.

\subsection{Input moments and radial comparison}

Write $R=(I-\Phi_A)^{-1}$ and $S_0=bb^*$. Four successive rational Stein
solves give $M_j=R^{j+1}S_0$, $0\le j\le3$:
\[
 M_0=S_0+\Phi_A(M_0),\qquad
 M_j=M_{j-1}+\Phi_A(M_j)\quad(1\le j\le3).
\]
All are positive semidefinite. Counting word splittings and using the scalar
generating functions gives
\begin{align}
 \Fnorm{H_f}_{\rm HS}^2&=cM_1c^*,\label{I:F:eq:g1-HS-input}\\
 J_2(f)^2&=c\bigl(2\Phi_A(M_2)+6\Phi_A^2(M_3)\bigr)c^*.
 \label{I:F:eq:g1-moment-Stein}
\end{align}
Indeed the coefficient energy at length $t$ is
$c\Phi_A^t(S_0)c^*$, each coefficient occurs $t+1$ times in its Hankel,
and
\[
 \sum_{t\ge0}t^2(t+1)x^t
 =\frac{2x}{(1-x)^3}+\frac{6x^2}{(1-x)^4}.
\]
Integer ceilings of the square roots in
\eqref{I:F:eq:g1-HS-input}--\eqref{I:F:eq:g1-moment-Stein}, bounded below by one,
supply rational $M,J$. These calculations use no approximation of
$\delta_k$.

For $0<\gamma<1$, let $D_\gamma e_w=\gamma^{|w|}e_w$ and
$h_\gamma(w)=\gamma^{|w|}h(w)$. Entrywise,
\begin{equation}
 H_{h_\gamma}=D_\gamma H_hD_\gamma.
 \label{I:F:eq:g1-radial-Hankel}
\end{equation}
This preserves the rank upper bound, and its realization is obtained by
replacing every letter matrix by $\gamma A_a$. For the input,
\begin{equation}
 \Fnorm{H_f-H_{f_\gamma}}
 \le\Fnorm{H_f-H_{f_\gamma}}_{\rm HS}
 \le(1-\gamma)J_2(f),
 \label{I:F:eq:g1-radial-bias}
\end{equation}
because $1-\gamma^t\le t(1-\gamma)$.

Let $G=H_h$ attain $\delta_k(H_f)$. Since zero is admissible,
$\Fnorm G\le2\Fnorm{H_f}\le2M$. The canonical observable realization on
$\operatorname{ran}G$, using right backward shifts, has parameters
$T_a,\ell,\beta$ satisfying
\[
 \sum_aT_a^*T_a+\ell^*\ell=I,
 \qquad\Fnorm\ell\le1,
 \qquad\Fnorm\beta=\Fnorm{Ge_\varepsilon}\le B.
\]
Padding to $k$ states changes the equality to an inequality if needed.
Therefore its damped realization belongs to
\begin{equation}
 \sum_a\widehat T_a^*\widehat T_a\preceq\gamma^2I,
 \qquad\Fnorm\ell\le1,
 \qquad\Fnorm\beta\le B.
 \label{I:F:eq:g1-compact-class}
\end{equation}
Using that $D_\gamma$ is a contraction gives
\begin{align}
 \Fnorm{H_f-H_{h_\gamma}}
 &\le\Fnorm{H_f-H_{f_\gamma}}
       +\Fnorm{D_\gamma(H_f-G)D_\gamma}\nonumber\\
 &\le\delta_k(H_f)+(1-\gamma)J_2(f)
 \le\delta_k(H_f)+e/3.
 \label{I:F:eq:g1-nearoptimal-damped}
\end{align}
The original optimizer may have an arbitrarily poor stability margin; no
such margin enters this argument.

\subsection{An infinite-Hankel Lipschitz modulus}

Consider two parameter triples $(A,c,b)$ and $(A',c',b')$ satisfying
\eqref{I:F:eq:g1-compact-class}. Suppose
\[
 \Fnorm{\begin{bmatrix}A_a-A'_a\end{bmatrix}^{\rm col}}\le\eta,
 \qquad\Fnorm{c-c'}\le\eta,
 \qquad\Fnorm{b-b'}\le B\eta.
\]
Joint column contractivity gives
$\sum_{|w|=t}\Fnorm{A_wx}^2\le\gamma^{2t}\Fnorm x^2$.
Telescoping the product $A_w-A'_w$ and applying this energy inequality on
both sides of each differing letter gives, at each word length $t$,
\[
 \left(\sum_{|w|=t}|cA_wb-c'A'_wb'|^2\right)^{1/2}
 \le2B\eta\gamma^t+Bt\eta\gamma^{t-1}.
\]
For $t=0$ the last term is zero. The $t$ product-replacement terms are
summed by Minkowski in the length-$t$ coefficient space. A second
Minkowski inequality with weights $t+1$ yields
\begin{align}
 \Fnorm{H_h-H_{h'}}
 &\le\Fnorm{H_h-H_{h'}}_{\rm HS}\nonumber\\
 &\le\frac{2B\eta}{\mu}
       +\frac{B\eta\sqrt{2+4\gamma^2}}{\mu^2}
 <\frac{5B\eta}{\mu^2}=e/3.
 \label{I:F:eq:g1-Lipschitz}
\end{align}
Here
\[
 \sum_{t\ge0}(t+1)\gamma^{2t}=\mu^{-2},\qquad
 \sum_{t\ge1}(t+1)t^2\gamma^{2t-2}
 =(2+4\gamma^2)\mu^{-4}.
\]
This supplies the summability needed for the infinite counting-measure
operator norm, rather than transferring a worst-word bound without a tail
argument.

\subsection{A rational stable net}

There are $N_A=2dk^2$ real coordinates in the complex vertical letter stack
and $2k$ in each boundary vector. Set
$v_A=\lceil\sqrt{N_A}\rceil$ and $v_c=\lceil\sqrt{2k}\rceil$.
Use the following finite lattices, including ceiling endpoints.
\begin{enumerate}[leftmargin=*]
\item For each real letter coordinate in $[-\gamma,\gamma]$, use step
$\eta/v_A$. Retain a raw stack only if
$\Fnorm{A^{\rm raw}}\le\gamma+\eta/2$, checked by the rational Hermitian
inequality $(\gamma+\eta/2)^2I-(A^{\rm raw})^*A^{\rm raw}\succeq0$.
Scale it by $\gamma/(\gamma+\eta/2)$.
\item For each real readout coordinate in $[-1,1]$, use step $\eta/v_c$.
Retain $\Fnorm{c^{\rm raw}}\le1+\eta/2$ and divide by $1+\eta/2$.
\item For each real initial-state coordinate in $[-B,B]$, use step
$B\eta/v_c$. Retain $\Fnorm{b^{\rm raw}}\le B(1+\eta/2)$ and divide by
$1+\eta/2$.
\end{enumerate}
Every retained output is rational and satisfies
\eqref{I:F:eq:g1-compact-class}. Every target triple in that class is covered:
rounding contributes at most $\eta/2$, $\eta/2$, $B\eta/2$ respectively,
and final scaling contributes at most the same amount. Thus the total
errors meet \eqref{I:F:eq:g1-Lipschitz}. The number of choices per real
coordinate is at most $3+2v_A/\eta$, proving the first inequality in
\eqref{I:F:eq:g1-net-size}. Since $e\le1$, $J\ge1$,
\[
 \mu\ge e/(2J),\qquad \eta\ge e^3/(120MJ^2),
\]
which proves the second count. The known column bound makes every candidate
mean-square stable, with CP spectral radius at most $\gamma^2$.

\subsection{Finite rational norm certificates and the global lower bound}

For each candidate form the stable block realization of $f-g_i$, of
dimension $m_0=n+k$, and solve its rational Stein equations
\[
 P_i=b_ib_i^*+\sum_aD_{i,a}P_iD_{i,a}^*,\qquad
 Q_i=c_i^*c_i+\sum_aD_{i,a}^*Q_iD_{i,a}.
\]
Its squared Hankel norm is the largest eigenvalue of $P_iQ_i$; all these
eigenvalues are nonnegative. No matrix square root is needed for the
following interval computation. Let $p=p_{\rm tr}$ and
$a_i=\operatorname{tr}((P_iQ_i)^p)$. If $a_i=0$, the error norm is zero.
Otherwise
\begin{equation}
 (a_i/m_0)^{1/(2p)}\le\Fnorm{H_f-H_{g_i}}
                         \le a_i^{1/(2p)}.
 \label{I:F:eq:g1-trace-interval}
\end{equation}
The uniform norm bound $Z=M+B/\mu$ follows from
$\Fnorm{H_{g_i}}\le\Fnorm{H_{g_i}}_{\rm HS}\le B/\mu$.
The width of the exact endpoints in \eqref{I:F:eq:g1-trace-interval} is at most
\[
 Z\bigl(m_0^{1/(2p)}-1\bigr)
 \le\frac{Z\log m_0}{p}\le e/18,
\]
using $p\ge\log m_0$ and $e^x-1\le2x$ for $0\le x\le1/2$.
Enclose the lower root from below and the upper root from above, each
within $e/18$, by rational bisection on $[0,2Z]$. Its comparisons are
inequalities between rational $2p$th powers. The resulting interval
$[a_i^-,u_i]$ contains the norm and has width at most $e/6$.

Choose a candidate minimizing $u_i$, let $U$ be that endpoint, and set
$L=\max\{0,U-5e/6\}$. The existence of the damped competitor
\eqref{I:F:eq:g1-nearoptimal-damped}, its net approximation
\eqref{I:F:eq:g1-Lipschitz}, and the norm-interval width imply
\[
 U\le\delta_k(H_f)+e/3+e/3+e/6
   =\delta_k(H_f)+5e/6.
\]
The chosen candidate is admissible and has actual error at most $U$.
This proves \eqref{I:F:eq:g1-global-gap}.

The verifier may rerun the finite net or inspect a full ledger containing
the exact Stein solutions, rational trace powers and root brackets. It
checks the interval widths and the minimum-upper-endpoint choice. The
ledger and its verification may have $N$ entries. A short certificate is
not asserted.

\subsection{Complexity qualifications}

For rational input, vectorized Stein systems have order $(n+k)^2$ and
rational coefficients. Fraction-free exact linear algebra, rational matrix
powers to exponent $p_{\rm tr}$, and scalar rational-power bisection have
polynomial bit cost in the parameters stated in Theorem~\ref{I:F:thm:g1-net}.
All input/grid rational denominators can be combined explicitly; their
heights are polynomial in input bit length, $\log M$, $\log J$,
$\log(1/e)$, $d$ and $k$. Also
$p_{\rm tr}=O(MJ\log(n+k)/e^2)$. Thus the full precision and verification
costs are included in the finite bound.

A checkable positive class fixes $d,k$ and bounds the input moments $M,J$
by constants or by specified polynomials in input dimension. Membership is
checked through the input Stein computations; it is not defined by an easy
optimum. On this class the construction is polynomial in $n$, rational
input length and $1/\varepsilon$. This does not establish a strict
complexity improvement over every optimized QE formulation or earlier
algorithm; that comparison and scientific priority need a separate audit.

For individually encoded algebraic entries, the compositum field degree
can be exponential in the total encoding length. A polynomial bit claim
therefore needs a supplied common number field and its degree, height and
embedding costs as parameters. The rational subclass avoids this issue.
Finally, this scheme bypasses finite-block completion: it supplies no
quantitative convergence modulus for the original $\beta_{k,L}$ hierarchy.

\section{The resulting theorem and its scope}
Let $\Sigma=\{a,b\}$, let $H_f[u,v]=f(uv)$ on $\ell^2(\Sigma^*)$, and put
\[
 q=1-r^2\in(0,1),\qquad
 f_{r,d}\text{ denotes \eqref{I:eq:binary-source} with terminal index }d.
\]
For a field $\mathbb F\in\{\RR,\RC\}$ define
\[
 D_{n,d}^{\mathbb F}(q)=
 \inf_{\substack{H_g\text{ bounded word-Hankel over }\mathbb F\\
                 \Rrank H_g\le n}}
       \Rnorm{H_{f_{r,d}}-H_g}^{2}.
\]
Write $D_k^{\mathbb F}=D_{k,k}^{\mathbb F}$. Ordinary matrix rank is
used throughout; the comparison class contains every bounded structured
Hankel of that rank. The letter matrices are required to be jointly
mean-square stable only for the explicit constructions, not as an
extra uniform-margin restriction on the optimum.

For integers $p\ge j\ge1$, let
\[
 \gamma_{j,p}^{\mathbb F}
 =\max_{D^*D+l^*l\preceq I_j}\Rnorm{lD^p},\qquad
 (E_{j,p}^{\mathbb F})^2=1-(\gamma_{j,p}^{\mathbb F})^2.
\]
Equivalently, $E_{j,p}^{\mathbb F}$ is the distance in scalar $H^2$
from $z^p$ to stable rational coefficient sequences of Hankel rank at
most $j$, over the indicated field. Let $\mathcal M_{j,p}^{\mathbb F}$
contain the lossless triples $(D,l,x)$ with
\[
 D^*D+l^*l=I_j,\quad \Rnorm x=1,\quad
 lD^p x=\gamma_{j,p}^{\mathbb F}>0.
\]
Every such $D$ is strictly stable and $x$ is cyclic. For one such triple,
write $\gamma=\gamma_{j,p}^{\mathbb F}$ and define
\begin{align*}
 a_i&=lD^i x\quad(0\le i\le p),\\
 S&=\sum_{i=0}^{p-1}a_i\,x^*D^{p-i}x,
 &v&=\sum_{i=0}^{p-2}a_iD^{p-2-i}x,\\
 K&=\RHerm\!\left(\sum_{i=0}^{p}D^{p-i}x\,lD^i\right),
 &\RHerm T&=(T+T^*)/2.
\end{align*}
The empty sum defining $v$ is zero. The response matrix $K$ is positive
definite. Define the finite variational coefficient
\begin{equation}
 b_{n,d}^{\mathbb F}=
 \max_{(D,l,x)\in\mathcal M_{n-1,d-1}^{\mathbb F}}
 \left\{|a_{d-2}|^2-\gamma\Re S+
                     \gamma v^*K^{-1}v\right\}.
 \label{I:R:eq:main-response-coefficient}
\end{equation}
This maximization is part of the answer: the general theorem does not
assert a closed-form evaluation or an efficient optimizer for every
budget.

\begin{theorem}[First-order structured approximation at every fixed budget]
\label{I:R:thm:main-general}
For every fixed pair of integers $2\le n\le d$, over each field separately,
\begin{equation}
 \boxed{\displaystyle
 D_{n,d}^{\mathbb F}(q)=
 (E_{n-1,d-1}^{\mathbb F})^2-b_{n,d}^{\mathbb F}q+o(q),
 \qquad q\downarrow0.}
 \label{I:R:eq:main-general}
\end{equation}
The maximum in \eqref{I:R:eq:main-response-coefficient} exists. No uniqueness
or nondegenerate Hessian hypothesis is imposed on its scalar extremizers.
There is a stable $n$-state construction with squared error
$(E_{n-1,d-1}^{\mathbb F})^2-b_{n,d}^{\mathbb F}q+O(q^2)$.
The global lower remainder is $o(q)$ in general.
On the diagonal $n=d=k$, the leading constants and the response
coefficients agree over the two fields. Consequently
$D_k^{\RR}(q)-D_k^{\RC}(q)=o(q)$.
The general diagonal field comparison uses the external scalar-extremal
classification identified in Section~\ref{I:R:subsec:real-diagonal-orbits};
the explicit cases through $k=4$ use the published low-degree extremals.
\end{theorem}

The proof has two distinct parts. Section~\ref{I:R:sec:all-budget-response}
proves the unrestricted complex passive-envelope response and its
real analogue. Section~\ref{I:R:sec:v4-matching} constructs a genuine stable
word-Hankel attaining that response to first order. A scalar compactness
argument then gives \eqref{I:R:eq:main-general}. The independent-budget
version changes the fast state dimension to $j=n-1$ and the coefficient
index to $p=d-1$ in both proofs.

\begin{corollary}[Explicit consequences]
\label{I:R:cor:main-explicit}
For each field separately,
\[
 \boxed{D_3^{\mathbb F}(q)=\frac9{16}-\frac17q+O(q^2).}
\]
At four states,
\[
 \boxed{D_4^{\mathbb F}(q)=
 \frac{1103-33\sqrt{33}}{2048}
 -\frac{855+103\sqrt{33}}{15872}q+o(q).}
\]
More generally, for every fixed $d\ge2$, put
$a_d=(d-1)^{d-1}/d^d$. Then
\[
 \boxed{D_{2,d}^{\mathbb F}(q)=
       1-a_d-\frac{a_d}{(d-1)^2}q+O(q^2).}
\]
These common expansions do not assert finite-$q$ equality between the
real and complex optima.
\end{corollary}

On the diagonal $n=d=k$, the earlier Hardy identification is
\[
 E_{k-1,k-1}=\frac1{C(k-1,2/3)},\qquad
 C(m,p)=\sup_{\Rnorm F_{H^p}=1}|[z^m]F|.
\]
The Hardy constants and outer-extremal structure belong to the cited
sources~\cite{brevig2020,outer2026}. Their use here is as scalar input to the
passive response and structured-completion arguments. The value $9/16$
at three states follows from $C(2,2/3)=4/3$; the new calculation here is the finite-parameter coefficient and its global matching proof.

\subsection{How the rank-one benchmark constrains the extension}
The exact one-state optimum is the unique $\tau\in(q,1)$ satisfying
\[
 \boxed{1-\tau=
 \frac{q}{\left(1+\sqrt{q\tau/(\tau-q)}\right)^2}.}
\]
Section~\ref{I:R:sec:full-channel} derives this monotone radical formulation
from the complete target singular decomposition and relates it to the
selected quartic. The dominant-vector bound alone would instead give
$1-q/(1+\sqrt q)^2$. It misses the true error at order $q^{5/2}$.
The higher-state first-order result therefore needs the separate
structured upper construction; sharpness of a projection envelope is
insufficient as a general finite-$q$ principle.

\subsection{Canonical row spaces and the global projection bound}
Let $G=H_g$ be any bounded rank-at-most-$n$ word-Hankel, and regard
its coefficient rows as columns. Their span is
$S=\operatorname{ran}G^{\mathsf T}$, where coefficient transpose
preserves the operator norm. If $L_c e_w=e_{cw}$, the identity
$L_c^*(G^{\mathsf T}e_u)=G^{\mathsf T}e_{uc}$ makes $S$ invariant
under $L_a^*,L_b^*$. In orthonormal coordinates on $S$, let
$A=L_a^*|_S$, $B=L_b^*|_S$, and $\ell x=x[\varepsilon]$.
Partitioning words by their first letter gives
\[
 A^*A+B^*B+\ell^*\ell=I_S.
\]
The observation map $Ox=(\ell T_{\operatorname{rev}(w)}x)_w$
is the isometric inclusion of $S$. Moreover
\[
 \sum_{|w|=h}\Rnorm{T_{\operatorname{rev}(w)}x}^2
       =\sum_{|v|\ge h}|x[v]|^2\longrightarrow0.
\]
Finite dimensionality makes this decay uniform on the unit sphere,
so the tuple is jointly mean-square stable. Padding smaller spaces
preserves the energy equality as an inequality.

For the target, put
\[
 u_r=\sqrt q\sum_{j\ge0}r^je_{a^j},\qquad
 v_r=\sqrt q\sum_{j\ge0}r^je_{a^jb^d}.
\]
Both vectors are unit, and $H_{f_{r,d}}^{\mathsf T}u_r=v_r$.
Since $G^{\mathsf T}u_r\in S$,
\[
 \Rnorm{H_{f_{r,d}}-G}^2\ge1-\Rnorm{P_Sv_r}^2,
 \qquad
 \Rnorm{P_Sv_r}^2
   =q\Rnorm{\ell B^d(I-rA)^{-1}}^2.
\]
Thus maximizing the right overlap over all passive $n$-state tuples
is a valid lower bound against the entire comparison class. This
argument imposes no prescribed competitor stability margin.


% Begin all_budget_response.tex

\section{An all-budget first variation of the passive envelope}
\label{I:R:sec:all-budget-response}

Let $k\ge2$, put $m=k-1$, and let $q=1-r^2\in(0,1)$.  All state
spaces in this section are complex Euclidean spaces.  Define
\begin{equation}
 \mathcal P_k(q)=
 \max_{A^*A+B^*B+\ell^*\ell\preceq I_k}
 q\Rnorm{\ell B^k(I-rA)^{-1}}^2
 \label{I:R:eq:Pk-definition}
\end{equation}
and
\begin{equation}
 h_k(t)=\max\left\{|\ell B^ke_1|^2:
 V=\begin{pmatrix}B\\ \ell\end{pmatrix},\quad
 V^*V\preceq I_k,\quad \Rnorm{Ve_1}^2=t\right\},
 \qquad 0\le t\le1.
 \label{I:R:eq:hk-definition}
\end{equation}
Both maxima are over compact sets; the inverse in
\eqref{I:R:eq:Pk-definition} exists because $\Rnorm A\le1$ and $r<1$.
The intrinsic scalar boundary constant is
\begin{equation}
 \Gamma_m=\max_{D^*D+l^*l\preceq I_m}\Rnorm{lD^m},
 \qquad 0<\Gamma_m<1.
 \label{I:R:eq:Gamma-intrinsic}
\end{equation}
Its classical observation-space interpretation gives
$E_m^2=1-\Gamma_m^2$, where $E_m$ is the distance from $z^m$ to scalar
stable rational series of Hankel rank at most $m$.  The separate
identification $E_m=1/C(m,2/3)$ uses the polynomial coefficient-extremal
structure of Brevig--Saksman~\cite{brevig2020} and the outer-extremal theorem
of Brevig--S\o vik~\cite{outer2026}.  The argument below uses only
\eqref{I:R:eq:Gamma-intrinsic} and classical finite model spaces.

This section generalizes the passive projection mechanism.
Section~\ref{I:R:sec:v4-matching} then constructs a word-Hankel
approximant attaining its first-order response.

\subsection{Exact elimination of the slow transition}

\begin{proposition}[Exact scalar reduction in every dimension]
\label{I:R:prop:allk-eliminate-A}
For every $k\ge1$ and $0<q<1$,
\begin{equation}
 \boxed{\displaystyle
 \mathcal P_k(q)=\max_{0\le t\le1}
 \frac{q h_k(t)}{(1-r\sqrt{1-t})^2}.}
 \label{I:R:eq:allk-elimination}
\end{equation}
\end{proposition}
\begin{proof}
For fixed $B,\ell$, put $S=I-B^*B-\ell^*\ell\succeq0$.  The row
operator norm can be written as the supremum of
$|\ell B^ky|/\Rnorm{(I-rA)y}$ over unit vectors $y$.  Set
$t=\Rnorm{By}^2+|\ell y|^2$, so $y^*Sy=1-t$.  Every feasible $A$ satisfies
\[
 \Rnorm{(I-rA)y}\ge1-r\Rnorm{Ay}\ge1-r\sqrt{1-t}.
\]
If $t<1$, choose
\[
 A=\frac{y(y^*S)}{\sqrt{1-t}}.
\]
Then $Ay=\sqrt{1-t}\,y$, and the positive-semidefinite
Cauchy--Schwarz inequality gives
\[
 A^*A=\frac{Syy^*S}{y^*Sy}\preceq S.
\]
Thus the denominator bound is attained.  If $t=1$, choose $A=0$.
Conjugating the state coordinates to send $y$ to $e_1$ proves both
inequalities in \eqref{I:R:eq:allk-elimination}.
\end{proof}


\subsection{Cyclicity and quantitative control of the passivity defect}

We first prove the structural fact responsible for normal coercivity.
For integers $n\ge j\ge0$, write
\[
 \gamma_{j,n}=\max\{\Rnorm{lD^n}:D^*D+l^*l\preceq I_j\},
 \qquad \gamma_{0,n}=0.
\]

\begin{lemma}[Strict improvement with the scalar state dimension]
\label{I:R:lem:strict-dimension-improvement}
Fix $n\ge1$.  Then
$\gamma_{0,n}<\gamma_{1,n}<\cdots<\gamma_{n,n}$.
\end{lemma}
\begin{proof}
Identify $H^2$ with its square-summable coefficient sequences and let
$S^*$ denote the backward shift.  The observation-map formulation is
\begin{equation}
 \gamma_{j,n}=
 \sup_{\substack{K\subset H^2,\ S^*K\subset K\\ \dim K\le j}}
 \Rnorm{P_K z^n}.
 \label{I:R:eq:gamma-model-space}
\end{equation}
Indeed, a passive pair $(D,l)$ has a contractive observation map
$Ox=(lD^ax)_{a\ge0}$, with finite-dimensional invariant range $K$.
Consequently
$|lD^nx|\le\Rnorm{P_Kz^n}\Rnorm{Ox}\le\Rnorm{P_Kz^n}\Rnorm x$.
Conversely, on a finite-dimensional invariant $K$, take
$D=S^*|_K$ and $lx=x[0]$.  Then $D^*D+l^*l=I_K$ and
$\Rnorm{lD^n}=\Rnorm{P_Kz^n}$.  This proves
\eqref{I:R:eq:gamma-model-space}.  Compactness of passive parameter sets
also shows that its supremum is attained.

The classical finite-model-space description writes
$K=K_\Theta=H^2\ominus\Theta H^2$ for a finite Blaschke product
$\Theta$.  If $\dim K\le j\le n$, then $z^n\notin K$, because the
backward-shift cyclic span of $z^n$ has dimension $n+1$.  Thus
$f=z^n-P_Kz^n$ is a nonzero member of $\Theta H^2$.  Reproducing kernels
span $H^2$, so there is a point $a$ in the disk with
$\langle f,\Theta k_a\rangle\ne0$.  If $b_a$ is the elementary
Blaschke factor at $a$, then
\[
 K_{\Theta b_a}=K_\Theta\mathbin{\oplus}\Theta\operatorname{span}\{k_a\}
\]
has one additional dimension and strictly increases the projection
norm of $z^n$.  Apply this construction to an attained $j$-state maximum
for each $j<n$.
\end{proof}


\begin{lemma}[Every scalar boundary maximizer is cyclic and lossless]
\label{I:R:lem:cyclic-lossless-maximizer}
Suppose $D^*D+l^*l\preceq I_m$, $\Rnorm x=1$, and
$|lD^mx|=\Gamma_m$.  Then $x$ is cyclic for $D$,
$D^*D+l^*l=I_m$, and $\rho(D)<1$.
\end{lemma}
\begin{proof}
The observation map $O$ is contractive and its range is backward-shift
invariant of dimension at most $m$.  Hence
\[
 \Gamma_m=|lD^mx|\le\Gamma_m\Rnorm{Ox}\le\Gamma_m,
\]
which forces $\Rnorm{Ox}=1$.  If the cyclic span of $x$ under $D$ had
dimension below $m$, restricting to that span would contradict
Lemma~\ref{I:R:lem:strict-dimension-improvement}, with $n=m$.
Put $T=I-D^*D-l^*l\succeq0$.  The finite-horizon energy identity is
\[
 1=\sum_{j=0}^{N-1}|lD^jx|^2
   +\sum_{j=0}^{N-1}\Rnorm{T^{1/2}D^jx}^2+\Rnorm{D^Nx}^2.
\]
Since the first sum tends to $\Rnorm{Ox}^2=1$, all the defect terms
vanish and $D^Nx\to0$.  Cyclicity gives $T=0$ and $D^N\to0$ on the
whole state space.  The latter is equivalent to $\rho(D)<1$ in finite
dimension.
\end{proof}


Define the compact normalized boundary set
\begin{equation}
 \mathcal M_m=\{(D,l,x):D^*D+l^*l=I_m,\quad \Rnorm x=1,
                    \quad lD^mx=\Gamma_m>0\}.
 \label{I:R:eq:boundary-maximizer-set}
\end{equation}
It is nonempty by attainment and Lemma~\ref{I:R:lem:cyclic-lossless-maximizer}.
Every member has an invertible cyclic matrix
$[x,Dx,\ldots,D^{m-1}x]$.  Compactness therefore bounds its least
singular value uniformly away from zero.  This also proves
$0<\Gamma_m<1$: positivity follows from a scalar stable realization,
and equality to one would put $z^m$ in an invariant space of dimension
at most $m$.

\begin{lemma}[Quantitative passivity-defect control]
\label{I:R:lem:quantitative-defect-control}
There are $C_m,\varepsilon_m>0$ such that every passive $(D,l)$ with
$\Gamma_m^2-\Rnorm{lD^m}^2\le\varepsilon_m$ satisfies
\begin{equation}
 \operatorname{tr}(I-D^*D-l^*l)
 \le C_m\bigl(\Gamma_m^2-\Rnorm{lD^m}^2\bigr).
 \label{I:R:eq:defect-gap-control}
\end{equation}
\end{lemma}
\begin{proof}
Put $a=lD^m$, $\mu=\Rnorm a$, and $x=a^*/\mu$.
By compactness and Lemma~\ref{I:R:lem:cyclic-lossless-maximizer}, when
$\mu$ is sufficiently close to $\Gamma_m$ the matrix
$K_x=[x,Dx,\ldots,D^{m-1}x]$ has least singular value at least a fixed
$\kappa_m>0$.  With $T=I-D^*D-l^*l$, the observation-space comparison
and telescoping give
\begin{align*}
 \mu^2&\le\Gamma_m^2\Rnorm{Ox}^2,\\
 \sum_{j=0}^{m-1}\Rnorm{T^{1/2}D^jx}^2
 &\le1-\Rnorm{Ox}^2
 \le\frac{\Gamma_m^2-\mu^2}{\Gamma_m^2}.
\end{align*}
The sum on the left is
$\Rnorm{T^{1/2}K_x}_{\mathrm F}^2\ge\kappa_m^2\operatorname{tr}T$.
This proves \eqref{I:R:eq:defect-gap-control}.
\end{proof}


\subsection{An explicit positive normal-response matrix}

For $(D,l,x)\in\mathcal M_m$, abbreviate $\Gamma=\Gamma_m$ and set
\begin{align}
 a_j&=lD^jx\quad(0\le j\le m),\qquad a_m=\Gamma,
 \label{I:R:eq:response-a}\\
 S&=\sum_{j=0}^{m-1}a_j\,x^*D^{m-j}x,
 &v&=\sum_{j=0}^{m-2}a_jD^{m-2-j}x,
 \label{I:R:eq:response-S-v}\\
 K&=\operatorname{Herm}\left(\sum_{j=0}^{m}D^{m-j}xlD^j\right),
 &\operatorname{Herm}(Z)&=\frac{Z+Z^*}{2}.
 \label{I:R:eq:response-K}
\end{align}
An empty sum is zero, so $v=0$ when $m=1$.

\begin{lemma}[Strict positivity of the response matrix]
\label{I:R:lem:response-K-positive}
For every $(D,l,x)\in\mathcal M_m$,
\begin{equation}
 \boxed{\displaystyle
 K\succeq\Gamma\sum_{j=0}^{m-1}D^jxx^*D^{*j}\succ0.}
 \label{I:R:eq:response-K-coercivity}
\end{equation}
Consequently $K^{-1}$ is uniformly bounded on $\mathcal M_m$.
\end{lemma}
\begin{proof}
Fix $Y\succeq0$ and perturb the lossless isometry
$U=\begin{pmatrix}D\\l\end{pmatrix}$ to
$U_\epsilon=U(I-\epsilon Y)^{1/2}$.  Its components satisfy
\[
 l_\epsilon D_\epsilon^mx
 =\Gamma-\frac\epsilon2
   \sum_{j=0}^{m}lD^jYD^{m-j}x+O(\epsilon^2).
\]
Its passivity defect is exactly $\epsilon Y$.  The observation-space
inequality and the first $m$ terms of the energy identity give
\[
 |l_\epsilon D_\epsilon^mx|^2
 \le\Gamma^2\left(1-\epsilon
   \sum_{j=0}^{m-1}x^*D^{*j}YD^jx+O(\epsilon^2)\right).
\]
Comparison of first derivatives yields
\[
 \operatorname{tr}(KY)\ge\Gamma\operatorname{tr}\left(
     \sum_{j=0}^{m-1}D^jxx^*D^{*j}Y\right).
\]
This holds for every $Y\succeq0$, proving the first inequality in
\eqref{I:R:eq:response-K-coercivity}.  Its right side is positive definite
by cyclicity.  Uniform boundedness of the inverse follows from
compactness.
\end{proof}


\begin{theorem}[All-budget first variation of the passive envelope]
\label{I:R:thm:all-budget-response}
For every fixed $k\ge2$, put $m=k-1$ and define
\begin{equation}
 b_k=\max_{(D,l,x)\in\mathcal M_m}
 \left\{|a_{m-1}|^2-\Gamma_m\operatorname{Re}S
                  +\Gamma_m v^*K^{-1}v\right\}.
 \label{I:R:eq:bk-response-formula}
\end{equation}
Then
\begin{equation}
 \boxed{\displaystyle
 h_k(t)=\Gamma_m^2t+b_kt^2+o(t^2),\qquad
 \mathcal P_k(q)=\Gamma_m^2+b_kq+o(q).}
 \label{I:R:eq:all-budget-envelope-expansion}
\end{equation}
The maximum in \eqref{I:R:eq:bk-response-formula} exists.  No uniqueness or
nondegenerate Hessian assumption on the scalar extremizers is needed.
In particular, for the binary target family, the projection lower bound
implies
\begin{equation}
 D_k(q)\ge E_{k-1}^2-b_kq+o(q).
 \label{I:R:eq:all-budget-D-lower}
\end{equation}
\end{theorem}

\begin{proof}
\emph{Exact first-column coordinates.}
Put $t=s^2$, with $s$ real near zero.  Every contraction $V$ in
\eqref{I:R:eq:hk-definition} can be written
\begin{equation}
 V=[sz,W_s(z)N],\quad \Rnorm z=1,\quad \Rnorm N\le1,\qquad
 W_s(z)=I-(1-\sqrt{1-s^2})zz^*.
 \label{I:R:eq:first-column-chart}
\end{equation}
Indeed, the Schur complement of $I-V^*V$ says that the remaining columns
$V_F$ satisfy
\[
 V_F^*\left(I+\frac{s^2}{1-s^2}zz^*\right)V_F\preceq I;
\]
its inverse positive square root is $W_s(z)$.  Partition
\[
 z=\begin{pmatrix}z_0\\z_F\\z_o\end{pmatrix},\qquad
 N=\begin{pmatrix}c\\D\\l\end{pmatrix},
 \qquad z_F\in\RC^m,
\]
where $c$ is a row.  Define $J(s,z,N)=\ell B^{m+1}e_1/s$ by its
continuous polynomial extension at $s=0$, and
$M(s)=\max_{z,N}|J(s,z,N)|$.  Then
\begin{equation}
 h_k(s^2)=s^2M(s)^2,\qquad J(0,z,N)=lD^mz_F,\qquad M(0)=\Gamma.
 \label{I:R:eq:M-and-J}
\end{equation}
At every leading maximizer, Lemma~\ref{I:R:lem:cyclic-lossless-maximizer}
forces $z_0=z_o=0$, $\Rnorm{z_F}=1$, $D^*D+l^*l=I_m$, and $c=0$.
Thus the leading maximizers are precisely the embedded scalar boundary
set, up to an output phase.

\emph{Uniform localization of the normal directions.}
Expansion of the matrix product gives
\begin{align}
 J(s,z,N)&=lD^mz_F+sJ_1(z,N)+O(s^2),
 \label{I:R:eq:J-first-order}\\
 J_1(z,N)&=z_ocD^{m-1}z_F+z_0lD^{m-1}z_F
       +\sum_{j=0}^{m-2}(lD^jz_F)(cD^{m-2-j}z_F).
 \label{I:R:eq:J1-explicit}
\end{align}
The remainder is uniform on the compact parameter domain.  Put
$\delta=\Gamma^2-|lD^mz_F|^2$.  Near the leading maximizers,
\[
 |z_0|^2+|z_o|^2\le\frac\delta{\Gamma^2},\qquad
 \Rnorm c^2\le\operatorname{tr}(I-D^*D-l^*l)\le C\delta
\]
by Lemma~\ref{I:R:lem:quantitative-defect-control}.  Therefore
$|J_1|\le C\sqrt\delta$ there.  Away from a fixed neighborhood of the
leading maximizers, $\delta$ is bounded below, so the same estimate holds
globally after enlarging $C$.  It follows that
\begin{equation}
 |J(s,z,N)|\le\Gamma-\frac\delta{2\Gamma}
                         +C|s|\sqrt\delta+Cs^2.
 \label{I:R:eq:normal-gap-upper}
\end{equation}
A fixed leading maximizer gives $M(s)\ge\Gamma-Cs^2$.
Completing the square in \eqref{I:R:eq:normal-gap-upper} proves
\begin{equation}
 M(s)=\Gamma+O(s^2),\qquad
 \delta=O(s^2)\quad\hbox{at every global maximizer.}
 \label{I:R:eq:normal-localization}
\end{equation}
In particular, all maximizing $z_0,z_o,c$ are $O(s)$ and their remaining
passivity defect is $O(s^2)$.

\emph{Polar coordinates and the complex phase.}
At any maximizing point write
\[
 z_0=s\lambda,\quad z_o=s\zeta,\quad c=sw,\quad
 z_F=\sqrt{1-s^2(|\lambda|^2+|\zeta|^2)}\,x,\qquad \Rnorm x=1.
\]
Let $Y_{\rm raw}=\begin{pmatrix}D_{\rm raw}\\l_{\rm raw}\end{pmatrix}$
denote the lower block of $N$, and put
$T=I-Y_{\rm raw}^*Y_{\rm raw}$.  Then
$T\succeq c^*c$, $\operatorname{tr}T=O(s^2)$, and the polar
factorization is
\begin{equation}
 Y_{\rm raw}=U(I-T)^{1/2},\quad U^*U=I_m,\quad
 T=s^2(w^*w+F),\qquad F\succeq0.
 \label{I:R:eq:normal-polar}
\end{equation}
The variables $\lambda,\zeta,w,F$ are uniformly bounded at maximizing
points.  Write $U=\begin{pmatrix}D\\l\end{pmatrix}$.  The slow-coordinate
unitary $\operatorname{diag}(e^{i\psi},I_m)$ preserves feasibility and
rotates both $J$ and $lD^mx$ by $e^{i\psi}$.  It therefore permits the
normalization $a=lD^mx>0$.  More explicitly, it changes
$x$ to $e^{i\psi}x$, $c$ to $e^{-i\psi}c$, $z_o$ to
$e^{i\psi}z_o$, and leaves $z_0,D,l$ unchanged.  Thus this normalization
discards no complex direction.  The base triples $(D,l,x)$ approach
$\mathcal M_m$, and always satisfy $a\le\Gamma$.

\emph{The normal quadratic form.}
Expand $(I-T)^{1/2}$, $W_s(z)$, and the finite product $\ell B^{m+1}e_1/s$.
Uniformly over the localized parameters,
\begin{equation}
 |J(s)|=a+s^2Q+O(s^3).
 \label{I:R:eq:J-normal-expansion}
\end{equation}
At a boundary triple in $\mathcal M_m$, the quadratic coefficient is
\begin{align}
 Q={}&-\frac12\operatorname{Re}S
   +\operatorname{Re}(\lambda a_{m-1})-\frac\Gamma2|\lambda|^2
   +\operatorname{Re}(wv)-\frac12wKw^*\nonumber\\
   &\hspace{18mm}-\frac\Gamma2|\zeta|^2
                       -\frac12\operatorname{tr}(KF).
 \label{I:R:eq:normal-quadratic-Q}
\end{align}
Off the boundary set the same formula uses the continuously varying
phase-normalized coefficients and replaces $\Gamma$ in the two
normalization costs by $a$.

For completeness, each contribution in
\eqref{I:R:eq:normal-quadratic-Q} can be read directly from the block
product.  Normalizing $z_F$ contributes
$-\Gamma(|\lambda|^2+|\zeta|^2)/2$.  The slow diagonal entry
$s^2\lambda$ contributes $\lambda a_{m-1}$.  A single slow-to-fast return
through the row $sw$ contributes
\[
 \sum_{j=0}^{m-2}(lD^jx)(wD^{m-2-j}x)=wv.
\]
Right contraction by $(I-s^2(w^*w+F))^{1/2}$ contributes
\[
 -\frac12\operatorname{Re}\sum_{j=0}^{m}
 lD^j(w^*w+F)D^{m-j}x
 =-\frac12wKw^*-\frac12\operatorname{tr}(KF).
\]
The first-column Schur factor contributes the perturbation
$D\mapsto D-s^2xx^*D/2$, hence $-S/2$.  Finally the slow readout is
$s^2\zeta$; using it requires a return to the slow state and first affects
$J$ at order $s^3$.  These observations prove the displayed expansion,
including its empty-sum case $m=1$.

Lemma~\ref{I:R:lem:response-K-positive} allows the quadratic form to be
maximized explicitly:
\begin{equation}
 \begin{gathered}
 F=0,\quad\zeta=0,\quad
 \lambda=\frac{\overline{a_{m-1}}}{\Gamma},\quad w=v^*K^{-1},\\
 Q_{\max}=-\frac12\operatorname{Re}S
       +\frac{|a_{m-1}|^2}{2\Gamma}+\frac12v^*K^{-1}v.
 \end{gathered}
 \label{I:R:eq:normal-quadratic-max}
\end{equation}
For the global upper bound, every sequence of maximizers has a
subsequence whose base triples approach $\mathcal M_m$.  Their base
contribution $a-\Gamma$ is nonpositive, all quadratic coefficients
converge continuously, and $K$ remains uniformly positive.  Therefore
\[
 \limsup_{s\to0}\frac{M(s)-\Gamma}{s^2}
 \le\max_{\mathcal M_m}Q_{\max}.
\]
For the reverse bound choose a boundary triple maximizing
$Q_{\max}$, take the parameters in
\eqref{I:R:eq:normal-quadratic-max}, and use the exact square roots in
\eqref{I:R:eq:first-column-chart} and \eqref{I:R:eq:normal-polar}.
They form a feasible family for all sufficiently small $|s|$ and give
the reverse liminf.  Consequently
\[
 M(s)=\Gamma+s^2\max_{\mathcal M_m}Q_{\max}+o(s^2).
\]
Equation~\eqref{I:R:eq:M-and-J} proves the first expansion in
\eqref{I:R:eq:all-budget-envelope-expansion}, with the coefficient
\eqref{I:R:eq:bk-response-formula}.

\emph{The final scalar maximization.}
The elementary estimate $h_k(t)\le t$ and
\[
 1-\sqrt{(1-q)(1-t)}\ge\frac{q+t}{4}
\]
show that the objective in \eqref{I:R:eq:allk-elimination} is at most
$16qt/(q+t)^2$.  The trial point $t=q$ has value tending to
$\Gamma^2>0$, so all maximizers have $x=t/q$ in a fixed compact
subinterval of $(0,\infty)$ when $q$ is small.  Uniformly on such
intervals,
\begin{align*}
 \frac{q h_k(qx)}{(1-\sqrt{(1-q)(1-qx)})^2}
 ={}&\frac{4\Gamma^2x}{(1+x)^2}\\
 &+q\left(\frac{4b_kx^2}{(1+x)^2}
        -\frac{2\Gamma^2x(x-1)^2}{(1+x)^3}\right)+o(q).
\end{align*}
The leading term has its unique maximum $\Gamma^2$ at $x=1$.
Thus all maximizing $x$ tend to one.  The leading term is never above
$\Gamma^2$ and the bracket tends to $b_k$ along every maximizing
sequence.  Taking $x=1$ gives the matching lower bound, proving the
second expansion.  The established projection inequality
$D_k(q)\ge1-\mathcal P_k(q)$ proves
\eqref{I:R:eq:all-budget-D-lower}.
\end{proof}


\subsection{Independent state budget and terminal delay}

The state budget and the number of terminal letters need not coincide.
This gives a larger class of targets without another nondegeneracy
assumption.

\begin{theorem}[Independent budget and delay]
\label{I:R:thm:independent-budget-delay-response}
Let $n\ge2$ and $d\ge n$.  Put $j=n-1$ and $p=d-1$, and define
\begin{align*}
 \mathcal P_{n,d}(q)&=
 \max_{A^*A+B^*B+\ell^*\ell\preceq I_n}
 q\Rnorm{\ell B^d(I-rA)^{-1}}^2,\\
 h_{n,d}(t)&=\max\{|\ell B^de_1|^2:
                  V=\begin{pmatrix}B\\\ell\end{pmatrix}
                  \text{ is a contraction},\ \Rnorm{Ve_1}^2=t\},\\
 \Gamma_{j,p}&=\max_{D^*D+l^*l\preceq I_j}\Rnorm{lD^p}.
\end{align*}
Let
\[
 \mathcal M_{j,p}=\{(D,l,x):D^*D+l^*l=I_j,\ \Rnorm x=1,
                                     \ lD^px=\Gamma_{j,p}>0\}.
\]
For a triple in this set put $\Gamma=\Gamma_{j,p}$ and
\begin{align*}
 a_i&=lD^ix\quad(0\le i\le p),\\
 S_p&=\sum_{i=0}^{p-1}a_i x^*D^{p-i}x,
 &v_p&=\sum_{i=0}^{p-2}a_iD^{p-2-i}x,\\
 K_p&=\operatorname{Herm}\sum_{i=0}^{p}D^{p-i}xlD^i,
 &b_{n,d}&=\max_{\mathcal M_{j,p}}
    \bigl\{|a_{p-1}|^2-\Gamma\operatorname{Re}S_p
                           +\Gamma v_p^*K_p^{-1}v_p\bigr\}.
\end{align*}
Then $\mathcal M_{j,p}$ is compact and nonempty,
\begin{equation}
 K_p\succeq\Gamma\sum_{i=0}^{j-1}D^ixx^*D^{*i}\succ0,
 \label{I:R:eq:independent-delay-K-positive}
\end{equation}
and
\begin{equation}
 \boxed{\displaystyle
 h_{n,d}(t)=\Gamma_{j,p}^2t+b_{n,d}t^2+o(t^2),\qquad
 \mathcal P_{n,d}(q)=\Gamma_{j,p}^2+b_{n,d}q+o(q).}
 \label{I:R:eq:independent-delay-expansion}
\end{equation}
For the target supported on $a^*b^d$, its squared rank-$n$ error therefore
satisfies
\[
 D_{n,d}(q)\ge1-\Gamma_{j,p}^2-b_{n,d}q+o(q).
\]
\end{theorem}
\begin{proof}
Proposition~\ref{I:R:prop:allk-eliminate-A} concerns neither the exponent
of $B$ nor its state dimension, so it gives exactly
\[
 \mathcal P_{n,d}(q)=\max_t
       \frac{q h_{n,d}(t)}{(1-r\sqrt{1-t})^2}.
\]
Since $p\ge j$, Lemma~\ref{I:R:lem:strict-dimension-improvement} applies to
the coefficient $z^p$ at every dimension up to $j$.  Its observation-map
proof implies that a maximizing unit $x$ for $lD^px$ is cyclic in the
$j$-dimensional state space.  The energy identity then forces
$D^*D+l^*l=I_j$ and strict stability exactly as in
Lemma~\ref{I:R:lem:cyclic-lossless-maximizer}.  Quantitative defect control
uses the invertible matrix $[x,Dx,\ldots,D^{j-1}x]$, so
\eqref{I:R:eq:defect-gap-control} holds with $lD^p$ in place of $lD^m$.

The exact first-column chart has fast state dimension $j$, and the
normalized product is now $J=\ell B^{p+1}e_1/s$.  Its leading term is
$lD^pz_F$; its first derivative is
\[
 z_ocD^{p-1}z_F+z_0lD^{p-1}z_F
   +\sum_{i=0}^{p-2}(lD^iz_F)(cD^{p-2-i}z_F).
\]
The same quantitative defect estimate controls this derivative by the
square root of the leading gap.  Consequently every maximizing normal
coordinate is $O(s)$ and the remaining defect is $O(s^2)$.
The polar coordinates \eqref{I:R:eq:normal-polar} therefore apply with
$j$ fast states.  Expanding the product gives the quadratic form
\eqref{I:R:eq:normal-quadratic-Q} with $a_{m-1},S,v,K$ replaced by
$a_{p-1},S_p,v_p,K_p$.

To justify the inverse, perturb the lossless fast isometry by
$(I-\epsilon Y)^{1/2}$.  Differentiating its coefficient $lD^px$ gives
$K_p$, while telescoping just the first $j$ energy terms gives the right
side of \eqref{I:R:eq:independent-delay-K-positive}.  Cyclicity makes that
matrix positive definite.  Thus the normal quadratic form can again be
maximized by a positive matrix inverse.  Compactness over
$\mathcal M_{j,p}$ and the nonpositive tangential base gap give the
matching upper and lower limits, proving the expansion of $h_{n,d}$.
The scalar maximization in $t/q$ is identical to that in
Theorem~\ref{I:R:thm:all-budget-response}, with leading maximum at $t/q=1$.
This proves the expansion of $\mathcal P_{n,d}$.  The target's projection
inequality gives the final lower bound.
\end{proof}


\subsection{Low-order response checks}

\begin{proposition}[The two-state response]
\label{I:R:prop:two-state-response-check}
For $k=2$, formula \eqref{I:R:eq:bk-response-formula} gives $b_2=1/4$.
\end{proposition}
\begin{proof}
Here $m=1$, and all boundary maximizers are equivalent by scalar phases
to $D=l=1/\sqrt2$, $x=1$, $\Gamma=1/2$.  Then
$a_0=1/\sqrt2$, $S=1/2$, $v=0$, and $K=1$.
Thus $b_2=1/2-1/4=1/4$.
\end{proof}


\begin{proposition}[Response of the degree-two scalar extremizer]
\label{I:R:prop:degree-two-response-check}
The scalar extremizer with denominator $1+z+z^2/2$ has
$\Gamma_2^2=7/16$ and normal-response coefficient $1/7$.
If all scalar degree-two boundary maximizers are equivalent to this
extremizer under the scalar symmetries, then $b_3=1/7$.
\end{proposition}
\begin{proof}
Use the observable companion coordinates
\[
 T=\begin{pmatrix}0&1\\-1/2&-1\end{pmatrix},\quad
 l_0=(1,0),\quad
 Q=\begin{pmatrix}8/5&4/5\\4/5&12/5\end{pmatrix},
 \qquad Q-T^*QT=l_0^*l_0.
\]
The canonical Euclidean realization is
$D=Q^{1/2}TQ^{-1/2}$, $l=l_0Q^{-1/2}$.  Equivalently, compute in the
$Q$ inner product, in which
\[
 \Gamma=\frac{\sqrt7}{4},\qquad
 x_{\rm raw}=\frac{(-1,-3)^{\mathsf T}}{2\sqrt7},\quad
 a_0=-\frac1{2\sqrt7},\quad a_1=-\frac3{2\sqrt7},\quad a_2=\Gamma.
\]
Direct rational algebra gives
\[
 S=\frac{23\sqrt7}{196},\qquad
 v_{\rm raw}=\frac{(1,3)^{\mathsf T}}{28},\qquad
 K_{\rm raw}=\begin{pmatrix}
 2\sqrt7/7&-\sqrt7/7\\-\sqrt7/28&13\sqrt7/28
 \end{pmatrix}.
\]
The last matrix is $Q$-self-adjoint and represents the same positive
response operator.  In particular,
\[
 v_{\rm raw}^*QK_{\rm raw}^{-1}v_{\rm raw}=\frac{3\sqrt7}{196}.
\]
The three terms in \eqref{I:R:eq:bk-response-formula} are therefore
\[
 \frac9{28}-\frac{23}{112}+\frac3{112}=\frac17.
\]
The value for this extremizer follows without any uniqueness assertion.
The last statement follows by taking the maximum in
\eqref{I:R:eq:bk-response-formula}.
\end{proof}


\subsection{The exact rank-one benchmark and the role of matching}

The case $k=1$ is singular because its limiting passive-overlap constant is zero (the limiting squared approximation error is one).
For a single-column contraction $V=(B,\ell)^{\mathsf T}$,
\[
 h_1(t)=\max_{|B|^2+|\ell|^2=t}|\ell B|^2=\frac{t^2}{4}.
\]
Proposition~\ref{I:R:prop:allk-eliminate-A} then gives the exact envelope
\begin{equation}
 \mathcal P_1(q)=\frac{q}{(1+\sqrt q)^2},\qquad
 t_{\max}=\frac{2\sqrt q}{1+\sqrt q}.
 \label{I:R:eq:rankone-envelope-benchmark}
\end{equation}
Thus
\[
 1-\mathcal P_1(q)
 =1-q+2q^{3/2}-3q^2+4q^{5/2}+O(q^3).
\]
The exact rank-one Hankel optimum instead has the expansion
\[
 D_1(q)=1-q+2q^{3/2}-3q^2+5q^{5/2}+O(q^3).
\]
The envelope is sharp through $q^2$ in this benchmark and has a strict
gap at order $q^{5/2}$.  Accordingly, the all-budget response theorem
supplies the global lower mechanism. The matching, $q$-dependent
sequential realization needed to identify the first correction of
$D_k$ is supplied in Section~\ref{I:R:sec:v4-matching}.

All arguments above also hold over the real field, restricting
$\mathcal M_m$ and the normal variables accordingly.  In
Lemma~\ref{I:R:lem:strict-dimension-improvement}, a real point $a\in(-1,1)$
can be chosen because a nonzero real analytic residual cannot vanish
on that entire interval.  The scalar Hardy extremal can be chosen real,
so the leading $\Gamma_m$ agrees in the two fields.
Section~\ref{I:R:subsec:real-diagonal-orbits} proves equality of the
diagonal response maxima by classifying the phases of all complex
scalar extremizers. For independent budget and delay, the two
field-specific response maxima remain distinct definitions.

% End all_budget_response.tex


% Begin matching.tex

\section{A structured realization matching the envelope to first order}
\label{I:R:sec:v4-matching}

The projection estimate becomes useful for optimal approximation only
when a genuine Hankel realization attains it to the required order.
The following construction supplies this step in every fixed dimension.
It applies over the complex field; real input data give a real
competitor.  It does not assert exact finite-parameter completion.

\begin{theorem}[Structured first-order matching]\label{I:R:thm:v4-matching}
Fix integers $n\ge2$ and $d\ge n$, and put $m=n-1$.  Let
$D\in\mathbb C^{m\times m}$ and $l\in\mathbb C^{1\times m}$ satisfy
\[
 D^*D+l^*l=I_m,\qquad \rho(D)<1,\qquad
 \Gamma=\|lD^{d-1}\|\in(0,1).
\]
Write $x=(lD^{d-1})^*/\Gamma$.  Fix any
$\lambda\in\mathbb C$ and $w\in\mathbb C^{1\times m}$.
For sufficiently small positive $s$, set $q=s^2$, $r=\sqrt{1-q}$ and
\begin{align}
 F_s&=(I_m-s^2w^*w)^{1/2},&
 N_s&=\begin{pmatrix}sw\\DF_s\\lF_s\end{pmatrix},&
 z_s&=\begin{pmatrix}s\lambda\\
       \sqrt{1-s^2|\lambda|^2}\,x\\0\end{pmatrix},\label{I:R:eq:v4-match-trial}\\
 W_s&=I_{n+1}-(1-r)z_sz_s^*,&
 V_s&=[\,sz_s,\ W_sN_s\,]
       =\begin{pmatrix}B_s\\l_s\end{pmatrix},&
 v_s&=\begin{pmatrix}r\\-sN_s^*z_s\end{pmatrix},\qquad
 A_s=e_0v_s^*.\label{I:R:eq:v4-match-tuple}
\end{align}
Define the row and the boundary vector
\begin{equation}
 \omega_s=s\,l_sB_s^d(I-rA_s)^{-1},\qquad
 \beta_s=s^{-1}(I-rA_s)\omega_s^*.
 \label{I:R:eq:v4-match-beta}
\end{equation}
For a word $u$, let $T_u$ denote the standard product of the letters
$T_a=A_s,T_b=B_s$, and set
\begin{equation}
 g_s(u)=l_sT_{\operatorname{rev}(u)}\beta_s.
 \label{I:R:eq:v4-match-series}
\end{equation}
This is an $n$-state word series: in the standard, unreversed convention
its readout is $\beta_s^{\mathsf T}$, its letters are
$A_s^{\mathsf T},B_s^{\mathsf T}$, and its boundary is
$l_s^{\mathsf T}$.  For all sufficiently small $s>0$ it is jointly
mean-square stable, and
\begin{equation}
 \boxed{\displaystyle
 \|H_{f_{r,d}}-H_{g_s}\|^2
     =1-\|\omega_s\|^2+O(q^2).}
 \label{I:R:eq:v4-match-conclusion}
\end{equation}
In particular, if a trial direction has
$\|\omega_s\|^2=\Gamma^2+cq+O(q^2)$, it gives a genuine structured
upper bound $1-\Gamma^2-cq+O(q^2)$.  Constants and neighborhoods may
depend on the fixed data.  The same conclusion holds for any other
analytic branch having the exact energy identity and the block orders
used in the proof below.
\end{theorem}

\begin{proof}
We keep ordinary coefficient transpose ${}^{\mathsf T}$ distinct from
Hermitian adjoint ${}^*$.  This distinction matters for complex data.

\emph{Passivity and the cancellation.}
Direct calculation gives $N_s^*N_s=I_m$, $\|z_s\|=1$ and
\[
 I_n-V_s^*V_s=v_sv_s^*,\qquad
 A_s^*A_s+B_s^*B_s+l_s^*l_s=I_n.
\]
The zero in the last component of $z_s$ is exact.  On
$\mathbb Ce_0\oplus\mathbb C^m$, the construction therefore has
\begin{align}
 B_s&=\begin{pmatrix}O(s^2)&O(s)\\
                sx+O(s^3)&D+O(s^2)\end{pmatrix},&
 l_s&=(0,\ l+O(s^2)),\label{I:R:eq:v4-match-blocks}\\
 A_s&=\begin{pmatrix}r&-sx^*D+O(s^3)\\0&0\end{pmatrix},&
 (I-rA_s)^{-1}
 &=\begin{pmatrix}q^{-1}&-rsq^{-1}x^*D+O(s)\\0&I_m\end{pmatrix}.
 \label{I:R:eq:v4-match-resolvent}
\end{align}
For fixed $d$, multiplication of the blocks gives
\[
 l_sB_s^de_0=s\,lD^{d-1}x+O(s^3)=s\Gamma+O(s^3),\qquad
 l_sB_s^d|_{\mathbb C^m}=lD^d+O(s^2).
\]
The identity $\Gamma x^*=lD^{d-1}$ cancels the order-$s$ fast
component of the overlap:
\begin{equation}
 \omega_s=(\Gamma+O(s^2),\
 s\{lD^d-\Gamma x^*D\}+O(s^3))
          =(\Gamma+O(q),O(s^3)).
 \label{I:R:eq:v4-match-cancellation}
\end{equation}
Consequently $\beta_s=(s\Gamma+O(s^3),O(s^2))$.

\emph{Joint stability by finite observability.}
Let $\Phi_s(X)=A_sXA_s^*+B_sXB_s^*$. The exact energy identity gives
$\Phi_s^*(I)=I-l_s^*l_s$. Since $\rho(D)<1$ and
$D^*D+l^*l=I_m$, the pair $(l,D)$ is observable: if all $lD^jx$
vanish, telescoping gives $\|D^t x\|=\|x\|$ for every $t$, hence $x=0$.
Thus the $m$ rows $lD^j$, $0\le j<m$, have full column rank.
Furthermore,
\[
 l_sB_s^dA_s=(l_sB_s^de_0)v_s^*
           =(s\Gamma+O(s^3))v_s^*.
\]
For small $s>0$ its scalar coefficient is nonzero. The $m+1$ rows
$v_s^*,l_s,l_sB_s,\ldots,l_sB_s^{m-1}$ converge to
$e_0^*,(0,l),(0,lD),\ldots,(0,lD^{m-1})$, a full-rank system.
Consequently a finite observation Gramian
\[
 Q_{s,L}=\sum_{|u|\le L}T_u^*l_s^*l_sT_u\succ0,
 \qquad L=d+1,
\]
is positive definite. Telescoping the energy identity yields
$(\Phi_s^*)^{L+1}(I)=I-Q_{s,L}\preceq(1-\eta_s)I$ with
$\eta_s=\lambda_{\min}(Q_{s,L})>0$. Positivity now implies
$\rho(\Phi_s)^{L+1}\le1-\eta_s<1$.
The transposed realization has the same stability spectral radius.

\emph{A reachability Gram estimate.}
Let
\[
 P_s=\sum_u T_{\operatorname{rev}(u)}\beta_s\beta_s^*
                    T_{\operatorname{rev}(u)}^*
      =\begin{pmatrix}p&z^*\\z&X\end{pmatrix}.
\]
It exists by joint stability and solves
$P_s=\beta_s\beta_s^*+A_sP_sA_s^*+B_sP_sB_s^*$.
The blocks above yield, with uniform implicit constants,
\begin{align*}
 qp&=q\Gamma^2+O(q^2)+O(s\|z\|)+O(q\|X\|)+O(q^2p),\\
 \|z\|&=O(s^3(1+p))+O(q\|z\|)+O(s\|X\|),\\
 \|X\|&=O(qp)+O(s\|z\|)+O(q^2).
\end{align*}
The last line follows by inverting the uniformly stable fast Stein
map.  Substitution in the middle line and absorption give
$\|z\|=O(s^3(1+p))$ and $\|X\|=O(qp+q^2)$.
The scalar equation becomes
$p=\Gamma^2+O(q(1+p))$.  Absorbing the last small multiple of $p$
proves
\begin{equation}
 (P_s)_{00}=\Gamma^2+O(q),\qquad
 (P_s)_{0F}=O(s^3),\qquad (P_s)_{FF}=O(q),\qquad
 \|P_s-\omega_s^*\omega_s\|=O(q).
 \label{I:R:eq:v4-match-gram}
\end{equation}

\emph{Exact interpolation and a fixed spectral gap.}
Define the reversed observation isometry and reachability operator by
\[
 (\mathcal O_s c)[u]=l_sT_{\operatorname{rev}(u)}c,\qquad
 \mathcal R_s e_u=T_{\operatorname{rev}(u)}\beta_s.
\]
The energy identity and joint stability imply
$\mathcal O_s^*\mathcal O_s=I_n$.
Exact coefficient multiplication gives
$H_{g_s}^{\mathsf T}=\mathcal O_s\mathcal R_s$ and
$\mathcal R_s\mathcal R_s^*=P_s$.
Let $S_s=\operatorname{ran}\mathcal O_s$, and introduce the real
orthonormal target vectors
\[
 u_r=s\sum_{j\ge0}r^j e_{a^j},\qquad
 u_h=s\sum_{j\ge0}r^j e_{a^jb^h}\quad(1\le h\le d),
 \qquad v_r=u_d.
\]
The target decomposes exactly as
\[
 H_{f_{r,d}}=u_rv_r^*+
       s\sum_{j=0}^{d-1}u_{d-j}e_{b^j}^*.
\]
By \eqref{I:R:eq:v4-match-beta},
\[
 \mathcal R_su_r=s(I-rA_s)^{-1}\beta_s=\omega_s^*,\qquad
 \mathcal O_s^*v_r=\omega_s^*,\qquad
 H_{g_s}^{\mathsf T}u_r=P_{S_s}v_r.
\]
Set $T=H_{f_{r,d}}^{\mathsf T}-H_{g_s}^{\mathsf T}$.
Then $Tu_r=(I-P_{S_s})v_r$ and
\begin{align*}
 T^*Tu_r
 &=H_{f_{r,d}}(I-P_{S_s})v_r\\
 &=a_su_r-s\sum_{j=0}^{d-1}u_{d-j}
                  [l_sB_s^j\omega_s^*],\qquad
 a_s=1-\|\omega_s\|^2.
\end{align*}
The first equality uses that the residual is orthogonal to the range
of $H_{g_s}^{\mathsf T}$.  The target is real, so its Hermitian adjoint
in this expression is its coefficient transpose.  For $j\ge1$,
$l_sB_s^je_0=O(s)$, while $l_se_0=0$.
Together with \eqref{I:R:eq:v4-match-cancellation}, this shows that the
component of $T^*Tu_r$ perpendicular to $u_r$ has norm $O(q)$.

Finally the isometry and \eqref{I:R:eq:v4-match-gram} give exactly
\[
 \|H_{g_s}^{\mathsf T}P_{u_r^\perp}\|^2
 =\|\mathcal R_sP_{u_r^\perp}\mathcal R_s^*\|
 =\|P_s-\omega_s^*\omega_s\|=O(q).
\]
Also $\|H_{f_{r,d}}^{\mathsf T}P_{u_r^\perp}\|=s$.
Hence the compression of $T^*T$ to $u_r^\perp$ has norm $O(q)$.
In this orthogonal decomposition,
\[
 T^*T=\begin{pmatrix}a_s&O(q)\\O(q)&O(q)\end{pmatrix},\qquad
 a_s\longrightarrow1-\Gamma^2>0.
\]
The scalar Schur-complement estimate across this fixed gap gives
$\|T\|^2=a_s+O(q^2)$.  Ordinary coefficient transpose preserves
operator norm, proving \eqref{I:R:eq:v4-match-conclusion}.
\end{proof}


\begin{corollary}[Matching the three-state correction]
\label{I:R:cor:v4-three-state-upper}
For the three-state, three-delay problem there is an explicit real
stable competitor with squared error
\[
 \frac9{16}-\frac17q+O(q^2).
\]
Combined with the independently established complex global envelope
$\mathcal P_{3,3}(q)=7/16+q/7+O(q^2)$, this gives, for each field
separately,
\[
 D_3^{\mathbb F}(q)=\frac9{16}-\frac17q+O(q^2),
 \qquad\mathbb F\in\{\mathbb R,\mathbb C\}.
\]
\end{corollary}

\begin{proof}
Take
\[
 D=\begin{pmatrix}-4/7&\sqrt5/7\\-5\sqrt5/14&-3/7\end{pmatrix},
 \qquad l=(-\sqrt7/14,\sqrt{35}/7).
\]
One checks $D^*D+l^*l=I_2$, $\sigma(D)=\{(-1+i)/2,(-1-i)/2\}$,
and $lD^2=(\sqrt7/4,0)$, so $\Gamma^2=7/16$ and $x=e_1$.
The explicit trial from the envelope computation takes
$N_s=\exp(sK)N_0$, $N_0=\begin{pmatrix}0_{1\times2}\\D\\l\end{pmatrix}$,
with the only nonzero upper-triangular entries
$K_{01}=2/7$, $K_{02}=\sqrt5/35$ and $K^{\mathsf T}=-K$, and
\[
 z_s=\frac{(-6s/7,1,0,0)^{\mathsf T}}
                   {\sqrt{1+36s^2/49}}.
\]
It has all the exact identities and parity orders in the theorem's
proof.  Its direct envelope expansion is
$\|\omega_s\|^2=7/16+q/7+O(q^2)$.
Apply \eqref{I:R:eq:v4-match-conclusion}, then sandwich the real upper
bound with the lower bound over all complex competitors.
This does not assert real--complex equality at finite $q$.
\end{proof}


% End matching.tex


\begin{proof}[Proof of Theorem~\ref{I:R:thm:main-general}]
Use the independent-budget/delay envelope theorem with fast dimension
$j=n-1$ and coefficient index $p=d-1$. It gives the global lower bound
$(E_{j,p}^{\mathbb F})^2-b_{n,d}^{\mathbb F}q+o(q)$.
Choose a boundary triple attaining the response maximum in
\eqref{I:R:eq:main-response-coefficient}. Cauchy--Schwarz equality implies
$x=(lD^p)^*/\gamma$. In Theorem~\ref{I:R:thm:v4-matching}, take
\[
 \lambda=\overline{a_{p-1}}/\gamma,\qquad w=v^*K^{-1}.
\]
These are precisely the maximizing normal directions of the envelope.
The evenness needed for the remainder follows from the exact trial.
Put $E=\operatorname{diag}(-1,I_m)$ and
$\widetilde E=\operatorname{diag}(-1,I_m,1)$. The definitions give
$N_{-s}=\widetilde E N_s$, $z_{-s}=\widetilde E z_s$,
$A_{-s}=EA_sE$, $B_{-s}=EB_sE$, and $l_{-s}=l_sE$.
It follows that $\omega_{-s}=-\omega_sE$: its slow coordinate
$J_s=\omega_se_0$ is even, and its fast coordinates are odd.
The apparent division by $s$ in $J_s$ is removable by
\eqref{I:R:eq:v4-match-cancellation}. Hence $J_s$ and
$\|\omega_s\|^2$ are real analytic and even. Their third-order
remainder is therefore $O(s^4)=O(q^2)$.
The fixed analytic trial has
$\Rnorm{\omega_s}^2=\gamma^2+b_{n,d}^{\mathbb F}q+O(q^2)$:
its slow component is the divided coefficient $J_s$, and the fast
component is $O(q^{3/2})$. The matching theorem therefore gives the
upper bound $(E_{j,p}^{\mathbb F})^2-b_{n,d}^{\mathbb F}q+O(q^2)$.
The two estimates prove the stated expansion in each field separately.
\end{proof}



% Begin three_state.tex

\section{The three-state envelope and its sharp first correction}
\label{I:R:sec:three-state}

The next state budget admits an unconditional complex envelope expansion.
The global Hessian calculation below improves the general first-response
remainder from $o(q)$ to $O(q^2)$.  This section establishes the envelope
and the resulting Hankel lower bound. Combined with
Theorem~\ref{I:R:thm:v4-matching}, it proves the three-state equality in
Corollary~\ref{I:R:cor:main-explicit}.

\begin{theorem}[Three-state passive envelope]\label{I:R:thm:three-envelope}
Let
\[
 h_3(t)=\max\{|\ell B^3e_1|^2:
 V=[B;\ell]\in\RC^{4\times3},\ V^*V\preceq I_3,
 \ \|Ve_1\|^2=t\}.
\]
As $t\downarrow0$ and $q\downarrow0$, respectively,
\begin{align}
 h_3(t)&=\frac7{16}t+\frac17t^2+O(t^3),
 \label{I:R:eq:three-h}\\
 \mathcal P_3(q)&:=\sup_{A^*A+B^*B+\ell^*\ell\preceq I_3}
 q\|\ell B^3(I-\sqrt{1-q}A)^{-1}\|^2
 =\frac7{16}+\frac17q+O(q^2).
 \label{I:R:eq:three-P}
\end{align}
Both expansions hold over the full complex comparison class.  Consequently,
for the binary target $f_{r,3}(a^jb^3)=(1-r^2)r^j$ and $q=1-r^2$,
\[
 D_3^{\RC}(q)\ge\frac9{16}-\frac17q-Cq^2
\]
for some $C<\infty$ and all sufficiently small positive $q$.
The same lower bound holds for real competitors.
\end{theorem}

\subsection{The entire boundary maximizing set}

Put $t=s^2$.  The exact contraction chart is
\[
 V=[sz,W_s(z)N],\quad \|z\|=1,\quad
 W_s(z)=I-(1-\sqrt{1-s^2})zz^*,\quad N^*N\preceq I_2.
\]
For fixed $s,z$, the amplitude $J_s=\ell B^2z_B$ is a holomorphic
polynomial in $N$, and $\ell B^3e_1=sJ_s$.  Applying the maximum-modulus
principle to the two singular-value variables of $N$ shows that its maximum
can be attained on $N^*N=I_2$.  Thus the exact maximization takes place on
\[
 X=S^7\times\operatorname{St}_{\RC}(4,2),\qquad
 h_3(s^2)=s^2\max_X|J_s|^2.
\]
This compact real-analytic manifold has dimension $7+12=19$.
The reduction asserts equality of maximum values; it does not assert that
every maximizer in the original contraction ball is isometric.

At $s=0$, write $z=(z_0,z_F,z_o)$ and $N=[c;D;l]$, with
$z_F\in\RC^2$. Then
\[
 J_0=lD^2z_F,\qquad
 \max|J_0|^2=\Gamma_2^2=1-E_2^2=\frac7{16}.
\]
Here $E_2=3/4$ is the established scalar rational-approximation constant.
The cyclic-lossless boundary lemma gives, at every maximum,
\[
 z_0=z_o=0,\quad \|z_F\|=1,\quad c=0,\quad D^*D+l^*l=I_2.
\]
For clarity, its rigidity has a short direct proof in this dimension.
A nonobservable pair reduces to one state, where
$\max |ld^2|^2=\max_{0\le x\le1}(1-x)x^2=4/27<7/16$.
For an observable pair let $G\succ0$ be its observation Gramian, so
$G\preceq I$.  Whitening cannot decrease $\|lD^2\|$; equality at the
maximum forces $G w^*=w^*$ for $w=lD^2$.
With $Q=I-G$, passivity gives $D^*QD\preceq Q$, making
$\ker Q$ invariant under $D$.
The maximizing vector $w^*$ is cyclic for the whitened pair, and the
whitened and original dynamics agree on $\ker Q$.
It follows that $Q=0$.  Losslessness then forces $c=0$.

There is also no unclassified scalar extremizer at this budget.
Write a minimizing degree-two model-space denominator as $P$ and put
$b=P_{\le2}(P^*/P)$.  Equality in the Hardy coefficient and H\"older
bounds makes $bP^2$ a coefficient extremal in $H^{2/3}$.
The Brevig--Saksman equality case identifies it, up to phases, with the
cube of
\[
 P_\theta(z)=\frac23\left(1+e^{i\theta}z+
                 \frac12 e^{2i\theta}z^2\right).
\]
Thus $bP^2=\lambda P_\theta^3$.
The right side has degree six, forcing $\deg b=\deg P=2$.
Its two roots are distinct; since $P^2$ divides $P_\theta^3$, necessarily
$P$ is proportional to $P_\theta$.
Lower-dimensional model spaces have strictly smaller monomial projection,
so no lower-degree case is omitted.
The canonical pair is consequently unique up to unitary state coordinates,
circle rotation, and readout phase.

A convenient representative is
\begin{equation}
 z_*=\begin{pmatrix}0\\1\\0\\0\end{pmatrix},\qquad
 N_*=
 \begin{pmatrix}
 0&0\\
 -4/7&\sqrt5/7\\
 -5\sqrt5/14&-3/7\\
 -\sqrt7/14&\sqrt{35}/7
 \end{pmatrix}.
 \label{I:R:eq:three-base}
\end{equation}
Its middle block $D$ and last row $l$ satisfy
$N_*^*N_*=I_2$ and $lD^2=(\sqrt7/4,0)$.
The nonzero lower-left entry of $D$ verifies cyclicity directly.
All boundary maxima form the orbit of \eqref{I:R:eq:three-base} under
\[
 B\mapsto e^{i\theta}U^*BU,\qquad
 \ell\mapsto e^{i\phi}\ell U,\qquad
 U=\operatorname{diag}(1,U_2),\quad U_2\in U(2).
\]
This orbit has dimension six.  Its continuous stabilizer is trivial:
fixing $z_*$ makes $U_2$ diagonal; the nonzero diagonal and off-diagonal
entries of $D$ force the letter and basis phases to vanish; the nonzero
entries of $l$ then force the readout phase to vanish.

\subsection{Exact response in rational coordinates}

The all-budget normal-response formula provides a compact exact computation
of the coefficient.  In a lossless canonical system, with
$x=(lD^2)^*/\Gamma$, put
\[
 a_j=lD^jx,\quad
 S=a_0x^*D^2x+a_1x^*Dx,\quad v=a_0x,
\quad
 K=\operatorname{Herm}\!\left(
 D^2xl+DxlD+xlD^2\right).
\]
The coefficient is
\begin{equation}
 b_3=|a_1|^2-\Gamma\operatorname{Re}S+
                  \Gamma v^*K^{-1}v.
 \label{I:R:eq:three-response}
\end{equation}
It suffices to use rational observable coordinates with their Gram metric:
\[
 T=\begin{pmatrix}0&1\\-1/2&-1\end{pmatrix},\quad
 l_0=(1,0),\quad
 Q=\begin{pmatrix}8/5&4/5\\4/5&12/5\end{pmatrix},\quad
 Q-T^*QT=l_0^*l_0.
\]
The corresponding canonical quantities are represented in the $Q$ inner
product by
\[
 \Gamma=\frac{\sqrt7}{4},\quad
 x=\frac{(-1,-3)^T}{2\sqrt7},\quad
 a_0=-\frac1{2\sqrt7},\quad a_1=-\frac3{2\sqrt7},
\]
\[
 S=\frac{23\sqrt7}{196},\qquad
 v=\frac{(1,3)^T}{28},\qquad
 K_Q=\begin{pmatrix}
 2\sqrt7/7&-\sqrt7/7\\
 -\sqrt7/28&13\sqrt7/28
 \end{pmatrix}.
\]
Here $K_Q$ represents the operator $K$ and is self-adjoint for the metric
$Q$. Direct exact algebra gives
\[
 v^*QK_Q^{-1}v=\frac{3\sqrt7}{196},\qquad
 b_3=\frac9{28}-\frac{23}{112}+\frac3{112}=\frac17.
\]
Uniqueness of the scalar boundary orbit makes this the global response,
rather than the response of only one chosen scalar extremizer.

\subsection{Complex transverse curvature and the stronger remainder}

The following exact certificate verifies every complex tangent direction.
Use seven real sphere coordinates
\[
 v=(x_1+iy_1,\ iy_2,\ x_3+iy_3,\ x_4+iy_4)^T,
 \qquad z(s)=\frac{z_*+sv}{\sqrt{1+s^2\|v\|^2}},
\]
and sixteen real coordinates for $N(s)=e^{isH}N_*$, $H=H^*$.
Let $h_0,\ldots,h_3$ be its diagonal entries and list its upper off-diagonal
entries lexicographically as
\[
 h_4+ih_5,\ h_6+ih_7,\ h_8+ih_9,\
 h_{10}+ih_{11},\ h_{12}+ih_{13},\ h_{14}+ih_{15}.
\]
For the resulting 23-vector $\xi$, direct multiplication gives
\[
 |J_s(z(s),N(s))|^2=\frac7{16}+
 s^2\left(a+g^T\xi+\frac12\xi^TH\xi\right)+O(s^3),
\]
where
\[
 a=-\frac{23}{112},\qquad
 g^T\xi=-\frac34x_1-\frac17h_5-
             \frac{5\sqrt5}{56}h_7-\frac{\sqrt7}{56}h_9.
\]
The real symmetric matrix $H$ has characteristic polynomial
\begin{align*}
 \det(\lambda I-H)
 ={}&\frac{\lambda^{10}(8\lambda+7)^4
                 (16\lambda^2+42\lambda+25)^2}{368293445632}\\
 &\quad\cdot(10976\lambda^2+56620\lambda+53525)
       (32\lambda^3+308\lambda^2+655\lambda+350).
\end{align*}
Every nonzero eigenvalue is negative: all roots are real by symmetry,
and each displayed factor has positive coefficients and nonzero constant.
Thus $\operatorname{rank}H=13$.
The 23-coordinate chart has four redundant directions, since the map
$H\mapsto e^{iH}N_*$ has the four-dimensional stabilizer $U(2)$ on the
orthogonal complement of $\operatorname{ran}N_*$.
The physical Hessian therefore has a six-dimensional kernel on the
19-dimensional manifold.  These are precisely the six symmetry directions
already exhibited; the transverse Hessian is strictly negative.

For an additional direct check, a solution of $H\xi=-g$ has only
\[
 x_1=-6/7,\qquad h_5=-2/7,\qquad h_7=-\sqrt5/35
\]
nonzero, and its optimized quadratic value is
\[
 a+\tfrac12g^T\xi=-23/112+39/112=1/7.
\]
The displayed chart specifies $H$ explicitly; its characteristic polynomial,
the response equation and the optimized value can be checked with exact
algebraic arithmetic.

Compactness forces every maximizing sequence for $s\to0$ toward the
single boundary orbit. On a local analytic slice transverse to that orbit,
the implicit-function theorem supplies a unique analytic maximizing branch.
The exact symmetry makes its value independent of the chosen orbit point,
so it is the global maximum for all sufficiently small $|s|$.
The physical class depends only on $s^2$; hence this analytic value is even.
Its expansion is therefore
\[
 \max_X|J_s|^2=\frac7{16}+\frac17s^2+O(s^4),
\]
which proves \eqref{I:R:eq:three-h}.

An exactly feasible real branch attaining this response is
\[
 z(s)=\frac{e_2-(6s/7)e_1}{\sqrt{1+36s^2/49}},\qquad
 N(s)=e^{sL}N_*,
\]
where $L$ is real skew-symmetric with
$L_{12}=2/7$, $L_{13}=\sqrt5/35$ and all other upper entries zero.
Both $\|z(s)\|=1$ and $N(s)^TN(s)=I_2$ hold exactly.
This branch certifies the real lower side of the envelope expansion;
the full complex Hessian certifies its global upper side.

Finally, exact elimination of $A$ gives
\[
 \mathcal P_3(q)=\max_{0\le t\le1}
 \frac{q h_3(t)}{(1-\sqrt{(1-q)(1-t)})^2}.
\]
The bound $h_3(t)\le t$ localizes all maximizing $t/q$ to a compact
subinterval of $(0,\infty)$.  The leading rescaled objective is
$7x/[4(1+x)^2]$, with a unique nondegenerate maximum at $x=1$.
Substitution of \eqref{I:R:eq:three-h} proves \eqref{I:R:eq:three-P}.
The analytic scalar critical branch additionally satisfies
$t_*(q)=q+(32/49)q^2+O(q^3)$.

% End three_state.tex


% Begin four_state.tex

\section{An exact four-state first correction}
\label{I:R:sec:four-state-correction}

The all-state response formula can be evaluated exactly at the
degree-three scalar extremizer.  Its classification follows from
equality in the Hardy-space bridge, rather than from a new scalar
Hessian calculation.  The scalar extremal and its uniqueness are the
external results of Brevig--Saksman~\cite[Theorem 2]{brevig2020}; the response computation below is
an additional finite algebraic calculation.

\subsection{Classification through equality in the Hardy bridge}

\begin{lemma}[Transfer of scalar-extremizer classification]
\label{I:R:lem:Hardy-classification-transfer}
Let $m\ge1$.  Suppose every normalized extremal for the $m$th
coefficient functional on $H^{2/3}$ is an outer polynomial cube, and
let its collection of cube roots be $\mathcal P_m$.  The minimizing
degree-$m$ finite Blaschke products in the rational approximation of
$z^m$ are precisely
\[
 \Theta=P^*/P,\qquad P\in\mathcal P_m,
\]
up to unimodular constants and the normalization of $P$.
Consequently, uniqueness of the coefficient extremal up to rotation
gives uniqueness of the boundary passive extremizer up to unitary
state coordinates and the corresponding scalar phases.
\end{lemma}
\begin{proof}
Write $C=C(m,2/3)$.  The established Hardy bridge gives
$E_m=1/C$.  Every minimizing model space has exact dimension $m$:
otherwise adjoining a reproducing-kernel direction with nonzero
overlap with the residual strictly improves its fit.  Write its
Blaschke product as $\Theta=D^*/D$, with $D$ an outer polynomial of
degree at most $m$, and set $b=P_{\le m}\Theta$.
The coefficient identity and H\"older inequality used in the bridge
are
\begin{equation}
 \Rnorm D_2^2=[bD^2]_m
 \le C\Rnorm{bD^2}_{2/3}
 \le C\Rnorm b_2\Rnorm D_2^2.
 \label{I:R:eq:classification-Holder-chain}
\end{equation}
At a minimizing space $\Rnorm b_2=E_m=1/C$, so equality holds
throughout.  In particular normalized $bD^2$ is a coefficient
extremal.  It is outer by the stated hypothesis, hence $b$ is outer
as well.  Equality in H\"older gives
$|b|/\Rnorm b_2=|D|/\Rnorm D_2$ almost everywhere on the circle.
Uniqueness of an outer function with its boundary modulus therefore
gives $b=cD$ for a nonzero scalar $c$.  It follows that the normalized
cube of $D$ is a coefficient extremal: $D$ is a scalar multiple of
some $P\in\mathcal P_m$.  Conversely the stationarity identity
$P^2=C P^*+O(z^{m+1})$ in the upper half of the Hardy bridge makes
$P^*/P$ a minimizing Blaschke product.  This proves the asserted
classification.

For a lossless observable boundary tuple $(D,l,x)$, its observation
map is an isometry onto the minimizing model space and
$O^*e_m=(lD^m)^*=\Gamma_m x$.  Thus the model space fixes the tuple
up to a unitary state change and fixes the unit cyclic vector by this
display.  A rotation of the scalar variable only produces the
corresponding phase changes in the tuple.
\end{proof}


The hypothesis is supplied in general by the inspected outer-extremal
theorem and polynomial structure.  For $m=2,3$, the complete
Brevig--Saksman classifications also state directly that the
extremizers are the displayed outer cubes, unique up to rotation.
The response
\begin{equation}
 b(D,l,x)=|a_{m-1}|^2-\Gamma_m\Re S
                         +\Gamma_m v^*K^{-1}v,
 \label{I:R:eq:four-state-response-recalled}
\end{equation}
with the notation of the all-state theorem, is invariant under these
state unitaries and phases.  For example, under
$D\mapsto e^{i\theta}D$, $x\mapsto e^{-im\theta}x$, one has
$a_j\mapsto e^{i(j-m)\theta}a_j$, $S\mapsto S$,
$v\mapsto e^{-i(m+2)\theta}v$, and $K\mapsto K$.
Thus evaluating one real representative determines the global
response for $m=3$.

\subsection{An exact algebraic certificate}

Set
\begin{equation}
 \alpha=\frac{\sqrt{15-\sqrt{33}}}{8},\qquad
 \beta=6\alpha-\frac{64}{3}\alpha^3,
 \qquad\gamma=2-8\alpha^2.
 \label{I:R:eq:four-state-alpha}
\end{equation}
Then $64\alpha^4-30\alpha^2+3=0$, $0<\alpha^2<1/4$, and
$P(z)=1+\beta z+\gamma z^2+\alpha z^3$ is the
Brevig--Saksman extremal denominator.  Equivalently
$\beta=\sqrt{3+\sqrt{33}/3}/2$ and
$\gamma=(1+\sqrt{33})/8$.

Use the observable companion coordinates
\[
 T=\begin{pmatrix}0&1&0\\0&0&1\\-\alpha&-\gamma&-\beta\end{pmatrix},
 \qquad l_0=(1,0,0).
\]
An exact solution of $Q-T^{\mathsf T}QT=l_0^{\mathsf T}l_0$ is
\begin{equation}
 Q=\frac1{333}
 \begin{pmatrix}
 342+768\alpha^2&1116\alpha-4224\alpha^3&528-2304\alpha^2\\
 1116\alpha-4224\alpha^3&798-576\alpha^2&1980\alpha-6272\alpha^3\\
 528-2304\alpha^2&1980\alpha-6272\alpha^3&858-192\alpha^2
 \end{pmatrix}.
 \label{I:R:eq:four-state-Q}
\end{equation}
Its leading principal minors are
\[
 \frac{38}{37}+\frac{256}{111}\alpha^2,\quad
 \frac{244}{111}+\frac{896}{333}\alpha^2,\quad
 \frac{968}{333}+\frac{11008}{999}\alpha^2,
\]
all positive.  Since $(T,l_0)$ is observable, the positive Stein
identity gives strict stability.  Whitening yields the Euclidean
lossless pair $D=Q^{1/2}TQ^{-1/2}$, $l=l_0Q^{-1/2}$.

The following calculation stays in companion coordinates and the
$Q$ inner product.  Put
\[
 x_0=Q^{-1}(l_0T^3)^{\mathsf T}
 =\begin{pmatrix}
 -\alpha^3\\21\alpha^2/8-9/16\\27\alpha^3/4-15\alpha/8
 \end{pmatrix},\qquad
 G=l_0T^3x_0=\frac{45}{64}-\frac{33}{32}\alpha^2.
\]
Then $\Gamma_3=\sqrt G$, while the normalized cyclic vector has raw
coordinates $x=x_0/\sqrt G$.
To evaluate \eqref{I:R:eq:four-state-response-recalled} without any
additional square roots, define
\begin{align*}
 c_j&=l_0T^jx_0,\\
 S_0&=\sum_{j=0}^{2}c_jx_0^{\mathsf T}QT^{3-j}x_0,\\
 v_0&=c_0Tx_0+c_1x_0,\\
 M_0&=\sum_{j=0}^{3}T^{3-j}x_0l_0T^j,\\
 K_0&=\tfrac12(M_0+Q^{-1}M_0^{\mathsf T}Q).
\end{align*}
Here $K=K_0/\sqrt G$, $v=v_0/G$, and $S=S_0/G^{3/2}$ in raw
coordinates.  The symmetric matrix $H_0=QK_0$ simplifies to
\begin{equation}
 H_0=\begin{pmatrix}
 21/32+7\alpha^2/16&\alpha/2-3\alpha^3&5/8-9\alpha^2/4\\
 \alpha/2-3\alpha^3&7/4-5\alpha^2/2&3\alpha-20\alpha^3/3\\
 5/8-9\alpha^2/4&3\alpha-20\alpha^3/3&3-19\alpha^2/3
 \end{pmatrix}.
 \label{I:R:eq:four-state-H0}
\end{equation}
Its leading principal minors are
\[
 \frac{21}{32}+\frac7{16}\alpha^2,\qquad
 \frac{1287}{1024}-\frac{915}{512}\alpha^2,\qquad
 \frac{153}{64}-\frac{189}{32}\alpha^2.
\]
They are positive using $0<\alpha^2<1/4$.  Thus the response matrix
is positive definite directly in this example.

Reduction of the following scalar products modulo
$64\alpha^4-30\alpha^2+3$ gives
\begin{align}
 c_2^2&=\frac{6075}{32768}-\frac{7767}{16384}\alpha^2,\notag\\
 S_0&=\frac{6669}{65536}-\frac{1809}{32768}\alpha^2,\notag\\
 v_0^{\mathsf T}QK_0^{-1}v_0
 &=\frac{171}{65536}+\frac{5625}{32768}\alpha^2.
 \label{I:R:eq:four-state-scalar-certificate}
\end{align}
Therefore the exact response is
\begin{equation}
 b_4=\frac{c_2^2-S_0+v_0^{\mathsf T}QK_0^{-1}v_0}{G}
    =\frac{75}{496}-\frac{103}{248}\alpha^2
    =\boxed{\frac{855+103\sqrt{33}}{15872}}.
 \label{I:R:eq:four-state-b4}
\end{equation}
The Stein equation, response matrix and coefficient admit direct exact
arithmetic checks in $\mathbb Q(\alpha)$ using the displayed identities.
Numerically $b_4=0.09114730044061346723567\ldots$; this decimal is
not used to obtain the coefficient.

\begin{theorem}[Four-state first boundary correction]
\label{I:R:thm:four-state-first-correction}
For $f_{r,4}(a^j b^4)=qr^j$, zero elsewhere, the squared optimum for
each field satisfies
\begin{equation}
 \boxed{\displaystyle
 D_4^{\mathbb F}(q)
  =\frac{1103-33\sqrt{33}}{2048}
    -\frac{855+103\sqrt{33}}{15872}\,q+o(q),
 \qquad\mathbb F\in\{\mathbb R,\RC\}.}
 \label{I:R:eq:four-state-optimum}
\end{equation}
\end{theorem}
\begin{proof}
Lemma~\ref{I:R:lem:Hardy-classification-transfer} and the external
degree-three classification identify every complex boundary
extremizer with the real representative used above, modulo symmetries
preserving the response.  The all-state normal-response lower bound
therefore has global coefficient \eqref{I:R:eq:four-state-b4}.
The real structured first-order matching construction uses this real
lossless tuple and its maximizing real normal directions.  It gives
the same coefficient in a rank-at-most-four competitor, with
$O(q^2)$ error relative to its analytic overlap branch.  Since
$\Gamma_3^2=G=(945+33\sqrt{33})/2048$, the complex lower and real
upper bounds give \eqref{I:R:eq:four-state-optimum} for each field.
The global lower remainder established here is $o(q)$; a global
$O(q^2)$ remainder is not inferred without a further tangential
nondegeneracy argument.
\end{proof}


% End four_state.tex


% Begin real_diagonal_lemma.tex

\section{Real representatives of every diagonal boundary extremizer}
\label{I:R:subsec:real-diagonal-orbits}

The diagonal case admits a further field comparison.  This argument
uses the exact scalar Hardy bridge and the strict increase of the
coefficient constants supplied by the October 2026 outer-extremal
result~\cite{outer2026}; it is not asserted for independent budget and delay.

\begin{lemma}[Positive coefficients and affine phases]
\label{I:R:lem:positive-Hardy-extremizers}
Fix $m\ge1$ and let $F=P^3$ be any normalized extremal for the
$m$-th coefficient on $H^{2/3}$, with $\Rnorm P_2=1$ and $P$ outer of
exact degree $m$.  There exist a polynomial
$Q(z)=\sum_{j=0}^m q_jz^j$ with every $q_j>0$ and real phases
$\alpha,\theta$ such that
\[
 P(z)=e^{i\alpha}Q(e^{i\theta}z).
\]
\end{lemma}
\begin{proof}
Write $P(z)=\sum_{j=0}^m p_jz^j$ and put
$Q(z)=\sum_{j=0}^m|p_j|z^j$.  Then $\Rnorm Q_2=1$ and
\[
 C(m,2/3)=|[P^3]_m|\le[Q^3]_m\le C(m,2/3).
\]
Thus $Q^3$ is also a coefficient extremal.  Its support
$S=\{j:q_j>0\}$ contains $0,m$: $P$ is outer of exact degree $m$.
If $q_j=0$, varying that real coefficient and renormalizing the
Euclidean norm shows $[Q^2]_{m-j}=0$; otherwise a positive first-order
variation would increase $[Q^3]_m$.  Consequently
\[
 i,h\in S,\quad i+h\le m\quad\Longrightarrow\quad m-i-h\in S.
\]
Taking $h=0$ shows $m-S=S$.  Applying this reflection once more gives
closure under addition when the sum is at most $m$.  Applying the rule
to $h$ and $m-i$ gives $i-h\in S$ whenever $i\ge h$ belong to $S$.
Let $g$ be the smallest positive element of $S$.  Repeated subtraction
shows every element of $S$ is divisible by $g$, and addition shows all
multiples up to $m$ occur.  Hence $g\mid m$ and
$Q(z)=R(z^g)$ with $\deg R=m/g$.  If $g>1$, then
\[
 C(m,2/3)=[Q^3]_m=[R^3]_{m/g}
 \le C(m/g,2/3)<C(m,2/3),
\]
a contradiction to the strict increase of coefficient constants.
Therefore $g=1$ and every $q_j$ is positive.

Equality in the first triangle inequality now says that every monomial
$p_ip_hp_k$ with $i+h+k=m$ has the same argument.  Put
$p_j=q_je^{i\alpha_j}$.  Comparing the nonzero terms with index triples
$(0,j,m-j)$ and $(1,j-1,m-j)$ gives
\[
 \alpha_j-\alpha_{j-1}=\alpha_1-\alpha_0\pmod{2\pi}
 \qquad(1\le j\le m).
\]
Thus $\alpha_j=\alpha_0+j\theta$ modulo $2\pi$, proving the claim.
\end{proof}


\begin{proposition}[All diagonal boundary models have real representatives]
\label{I:R:prop:real-diagonal-boundary-models}
Every maximizing complex passive triple in $\mathcal M_m$ is equivalent
to a real triple through state unitaries and phases of the scalar
transition and output.  In particular,
\[
 b_{k,k}^{\RR}=b_{k,k}^{\RC}\qquad(k=m+1\ge2).
\]
Consequently the real and complex structured optima have the same
first-order diagonal expansion, without an assertion of equality at
finite $q$.
\end{proposition}
\begin{proof}
A maximizing passive pair is lossless, stable and observable by
Lemma~\ref{I:R:lem:cyclic-lossless-maximizer}.  Its observation map is an
isometry onto an $m$-dimensional model space $K_\Theta$ minimizing the
error of approximation to $z^m$.  Write
$\Theta=e^{i\phi}D_0^*/D_0$, with $D_0$ outer of degree at most $m$
and reversal of order $m$, and set $b=P_{\le m}\Theta$.
The Hardy bridge gives $\Rnorm b_2=1/C(m,2/3)$, while its exact
coefficient identity gives
\[
 \Rnorm{D_0}_2^2=|[bD_0^2]_m|
 \le C(m,2/3)\Rnorm{bD_0^2}_{2/3}
 \le C(m,2/3)\Rnorm b_2\Rnorm{D_0}_2^2
 =\Rnorm{D_0}_2^2.
\]
Therefore $bD_0^2$, after normalization and a phase, is a Hardy
coefficient extremal and hence outer.  Since $D_0$ is outer, $b$ is
outer as well.  Equality in H\"older gives
$|b|^2=c|D_0|^2$ almost everywhere on the circle.  Uniqueness of outer
functions with prescribed boundary modulus then gives $b=\mu D_0$.
It follows that the normalized $D_0^3$ is a coefficient extremal.
Lemma~\ref{I:R:lem:positive-Hardy-extremizers} shows that $D_0$ has real
positive coefficients after a scalar phase and circle rotation.
Thus $K_\Theta$ is a circle rotation of a model space invariant under
coefficient conjugation.  Choosing a real orthonormal basis of that
space gives real canonical matrices and readout.  Every original
maximizing passive pair is unitarily equivalent to its canonical
model, and its maximizing $x$ equals $(lD^m)^*/\Gamma_m$.
This proves the stated orbit classification.

For clarity, the response formula is invariant under
\[
 (D,l,x)\mapsto(U^*DU,lU,U^*x),\quad
 (D,l,x)\mapsto(D,e^{i\varphi}l,e^{-i\varphi}x),\quad
 (D,l,x)\mapsto(e^{i\theta}D,l,e^{-im\theta}x).
\]
The first operation conjugates $K$ and transforms $v$ by $U^*$.
The second fixes $a_j,S,K$ and multiplies $v$ by a phase.
For the third, $a_j$ acquires $e^{i(j-m)\theta}$,
$S,K$ are unchanged, and $v$ acquires $e^{-i(m+2)\theta}$.
Every term in the response coefficient is unchanged.  Since every
complex boundary triple has a real representative, maximizing over
complex triples gives the same value as maximizing over real triples.
The structured first-order theorem supplies the final conclusion.
\end{proof}


% End real_diagonal_lemma.tex


% Begin full_channel.tex

\section{All-state, all-delay projection bounds retaining every singular channel}
\label{I:R:sec:full-channel}

The dominant-vector argument controls only one singular direction of the target.
We now retain all of them and eliminate the waiting transition exactly. This
gives a finite-dimensional characterization of an invariant-subspace relaxation
for every state budget. Structured Hankel completion remains an additional step.

Fix integers $n,m\ge1$, $0<q<1$, and $r=\sqrt{1-q}$. On the binary alphabet
$\{a,b\}$ let
\[
 f_{r,m}(a^j b^m)=qr^j\quad(j\ge0),\qquad
 f_{r,m}(w)=0\text{ otherwise},\qquad
 D_{n,m}(q)=\delta_n^{\RC}(H_{f_{r,m}})^2.
\]
All norms here use counting measure on words. With $e_w$ the coordinate vectors,
put
\[
 v=\sqrt q\sum_{j\ge0}r^j e_{a^j b^m},\qquad
 e_j=e_{b^j}\quad(0\le j<m),\qquad
 \mathcal U=(v,e_0,\ldots,e_{m-1}).
\]
These columns are orthonormal, and direct grouping of the target's rows gives
\begin{equation}
 H_{f_{r,m}}^*H_{f_{r,m}}
   =vv^*+q\sum_{j=0}^{m-1}e_je_j^*
   =\mathcal U C\mathcal U^*,\qquad C=\operatorname{diag}(1,qI_m).
 \label{I:R:eq:fc-singular}
\end{equation}
Indeed the pure-$a$ rows give the unit singular channel; rows $a^i b^h$,
$1\le h\le m$, give the other $m$ orthogonal channels.

Let $L_a,L_b$ be the left creation shifts and define
\begin{equation}
 R_{n,m}(q)=\inf_{\substack{\dim S\le n\\L_a^*S,L_b^*S\subseteq S}}
          \Rnorm{H_{f_{r,m}}(I-P_S)}^2.
 \label{I:R:eq:fc-relaxation}
\end{equation}
For every bounded rank-at-most-$n$ word-Hankel $G$, the row space
$S=\operatorname{ran}G^*$ is invariant under these backward shifts and
$G(I-P_S)=0$. Consequently
\begin{equation}
 D_{n,m}(q)\ge R_{n,m}(q).
 \label{I:R:eq:fc-relaxation-bound}
\end{equation}
The definition allows invariant spaces that need not be cyclic row spaces of a
word-Hankel operator. Equality in \eqref{I:R:eq:fc-relaxation-bound} therefore
requires a separate argument.

\subsection{An exact residual matrix for the relaxation}

Use standard matrix products $T_{a_1\cdots a_j}=T_{a_1}\cdots T_{a_j}$, with
$T_a=A,T_b=B$. For a passive tuple
\begin{equation}
 A^*A+B^*B+\ell^*\ell\preceq I_n,
 \label{I:R:eq:fc-passive}
\end{equation}
the appropriate observation convention is
\[
 (\mathcal Ox)[w]=\ell T_{\operatorname{rev}(w)}x.
\]
The reversal is essential for left-backward-shift coordinates: the coordinate
at $a^j b^m$ is $\ell B^m A^j x$. Telescoping
\eqref{I:R:eq:fc-passive} over complete word levels proves
$\Rnorm{\mathcal O}\le1$, and
$L_c^*\mathcal O=\mathcal OT_c$ for $c=a,b$.
Define
\begin{equation}
 w_0=\sqrt q\,\ell B^m(I-rA)^{-1},\qquad
 Z=\begin{pmatrix}\ell\\\ell B\\\vdots\\\ell B^{m-1}\end{pmatrix},
 \qquad W=\begin{pmatrix}w_0\\Z\end{pmatrix},\qquad
 E=C-C^{1/2}WW^*C^{1/2}.
 \label{I:R:eq:fc-E}
\end{equation}
Here $W=\mathcal U^*\mathcal O$. The inverse exists uniformly over the passive
class because $r<1$ and $\Rnorm{A}\le1$.

\begin{theorem}[Exact passive representation of the projection relaxation]
\label{I:R:thm:fc-residual}
For every $n,m\ge1$ and $0<q<1$,
\begin{equation}
 R_{n,m}(q)=\min_{A^*A+B^*B+\ell^*\ell\preceq I_n}
                  \lambda_{\max}(E).
 \label{I:R:eq:fc-passive-min}
\end{equation}
\end{theorem}
\begin{proof}
For an invariant $S$ in \eqref{I:R:eq:fc-relaxation}, take orthonormal state
coordinates, set $A=L_a^*|_S$, $B=L_b^*|_S$, and let $\ell$ evaluate the empty
coordinate. The observation map is the isometric inclusion of $S$ and the
energy identity gives \eqref{I:R:eq:fc-passive}, with equality. Formula
\eqref{I:R:eq:fc-singular} then yields
\[
 \Rnorm{H_{f_{r,m}}(I-P_S)}^2
 =\lambda_{\max}\!\left(C^{1/2}\mathcal U^*(I-P_S)\mathcal U C^{1/2}\right)
 =\lambda_{\max}(E).
\]
Conversely, for any passive tuple, $K=\operatorname{ran}\mathcal O$ is a
backward-invariant space of dimension at most $n$, and
$\mathcal O\mathcal O^*\preceq P_K$. Thus its matrix $E$ dominates the exact
residual matrix for $K$ in the positive-semidefinite order. Its largest
eigenvalue is at least $R_{n,m}(q)$. The two comparisons prove equality of the
infima. The passive class is compact and its objective continuous, so the
minimum exists. Smaller state spaces are included by zero padding.
\end{proof}


\subsection{Threshold reduction and exact elimination of the waiting matrix}

For a threshold $\tau>q$, set
\begin{equation}
 Q_\tau=I_n+\frac{q}{\tau-q}Z^*Z\succ0.
 \label{I:R:eq:fc-Q}
\end{equation}
The lower-right block of $\tau I-E$ is
$(\tau-q)I_m+qZZ^*\succ0$. Its Schur complement, followed by the Woodbury
identity, gives the exact equivalence
\begin{equation}
 E\preceq\tau I_{m+1}
 \quad\Longleftrightarrow\quad
 w_0Q_\tau^{-1}w_0^*\ge1-\tau.
 \label{I:R:eq:fc-threshold}
\end{equation}
This includes every subdominant target channel through $Z^*Z$.

\begin{lemma}[A weighted trust-region elimination]
\label{I:R:lem:fc-trust}
Fix $B,\ell$ with $D=I-B^*B-\ell^*\ell\succeq0$, and let $Q\succ0$.
For a unit vector $y$, put $d=y^*Dy$ and
\begin{equation}
 \mathcal T_Q(y,d)=\min_{\Rnorm{z}^2\le d}(y-rz)^*Q(y-rz).
 \label{I:R:eq:fc-trust-value}
\end{equation}
Then
\begin{align}
 &\max_{A^*A\preceq D}
 q\ell B^m(I-rA)^{-1}Q^{-1}(I-rA^*)^{-1}(B^*)^m\ell^*
 \nonumber\\[-2pt]
 &\hspace{30mm}=\max_{\Rnorm{y}=1}
       \frac{q|\ell B^m y|^2}{\mathcal T_Q(y,y^*Dy)}.
 \label{I:R:eq:fc-eliminate-A}
\end{align}
If $d=0$, then $\mathcal T_Q(y,0)=y^*Qy$. If $d>0$, its unique minimizer is
\begin{equation}
 z_\mu=r(r^2Q+\mu I)^{-1}Qy,
 \qquad \Rnorm{z_\mu}^2=d,
 \label{I:R:eq:fc-secular}
\end{equation}
where the last equation has a unique positive solution $\mu$.
\end{lemma}
\begin{proof}
As $A$ ranges over $A^*A\preceq D$, its value $Ay$ ranges over exactly the ball
$\Rnorm{z}\le\sqrt d$. Necessity follows from $\Rnorm{Ay}^2\le y^*Dy$. For
$d>0$, the rank-one choice
\[
 A=zy^*D/d
\]
has $Ay=z$ and $A^*A\preceq D$, since
$Dyy^*D\preceq dD$ and $\Rnorm{z}^2\le d$. The case $d=0$ is immediate.

For a fixed $A$, the weighted row norm on the left of
\eqref{I:R:eq:fc-eliminate-A} is
\[
 \max_{x\ne0}\frac{q|\ell B^m(I-rA)^{-1}x|^2}{x^*Qx}
 =\max_{y\ne0}\frac{q|\ell B^m y|^2}
                          {((I-rA)y)^*Q((I-rA)y)}.
\]
Homogeneity permits $\Rnorm{y}=1$. Maximizing over $A,y$ and using the exact
ball of attainable $Ay$ gives \eqref{I:R:eq:fc-eliminate-A}. The denominator is
strictly positive because $r<1$ and $d\le1$.

The trust-region objective is strictly convex. Its unconstrained minimizer is
$y/r$, whose norm $1/r$ exceeds the feasible radius $\sqrt d\le1$.
The constraint is active, and the multiplier equations give
\eqref{I:R:eq:fc-secular}. Diagonalizing $Q$ shows that
$\Rnorm{z_\mu}^2$ decreases strictly from $1/r^2$ to zero as $\mu$ increases
from zero to infinity. This proves existence, uniqueness and sufficiency.
\end{proof}


Combining the preceding statements gives an unconditional all-state threshold
formula. Define
\begin{equation}
 \Psi_{n,m}(q,\tau)=
 \max_{\substack{B^*B+\ell^*\ell\preceq I_n\\\Rnorm{y}=1}}
 \frac{q|\ell B^m y|^2}{\mathcal T_{Q_\tau}(y,y^*Dy)}.
 \label{I:R:eq:fc-Psi}
\end{equation}
Here $Z,Q_\tau,D$ depend only on $B,\ell$. Then, for every $\tau>q$,
\begin{equation}
 \boxed{R_{n,m}(q)\le\tau
       \quad\Longleftrightarrow\quad
       \Psi_{n,m}(q,\tau)\ge1-\tau.}
 \label{I:R:eq:fc-general-threshold}
\end{equation}
The remaining optimization is compact. Each objective evaluation requires one
positive scalar secular root, in place of an optimization over $A$.
Equation \eqref{I:R:eq:fc-general-threshold} characterizes the relaxation at
thresholds above $q$; it does not assert structured attainability, or determine
values below $q$ when the state budget is large.

\subsection{The exact rank-one benchmark is recovered}

For $n=m=1$, write $\kappa=q/(\tau-q)$ and $K=1+\kappa$. Phases may be aligned
so that $a=|A|\ge0$ minimizes $|1-rA|$. With $x=|B|^2$ and $z=|\ell|^2$,
the overlap in \eqref{I:R:eq:fc-threshold} is
\[
 \frac{qzx}{(1-ra)^2(1+\kappa z)},\qquad a^2+x+z\le1.
\]
At fixed $a$, the available budget $t=1-a^2$ is fully used. Direct
differentiation gives
\[
 z_* =\frac{t}{1+\sqrt{1+\kappa t}},\qquad
 \max_{0\le z\le t}\frac{z(t-z)}{1+\kappa z}
       =\frac{t^2}{(1+\sqrt{1+\kappa t})^2}.
\]
The remaining maximization is elementary. Equivalently maximize
\[
 F(a)=\frac{\sqrt{1+\kappa(1-a^2)}-1}{1-ra},\qquad0\le a\le1.
\]
Its interior stationary equation is
\[
 \sqrt{1+\kappa(1-a^2)}+\frac{\kappa a}{r}=1+\kappa.
\]
Solving with this unsquared sign condition gives its unique maximum at
$a_*=r\sqrt K/(\sqrt K+\sqrt q)$, where
$F(a_*)=\kappa/(1+\sqrt{qK})$. Consequently
\begin{equation}
 \boxed{\Psi_{1,1}(q,\tau)
      =\frac{q}{\left(1+\sqrt{q\tau/(\tau-q)}\right)^2}.}
 \label{I:R:eq:fc-scalar-Psi}
\end{equation}

\begin{theorem}[Monotone radical representation of the exact F1 value]
\label{I:R:thm:fc-F1-radical}
The exact rank-one value $D_1(q)=D_{1,1}(q)$ is the unique root $\tau\in(q,1)$
of
\begin{equation}
 \boxed{1-\tau
       =\frac{q}{\left(1+\sqrt{q\tau/(\tau-q)}\right)^2}.}
 \label{I:R:eq:fc-F1-root}
\end{equation}
It equals the real optimum as well.
\end{theorem}
\begin{proof}
The right side of \eqref{I:R:eq:fc-F1-root} increases continuously from zero as
$\tau$ increases from $q$, while $1-\tau$ decreases from $1-q$ to zero.
There is exactly one crossing in $(q,1)$. Equations
\eqref{I:R:eq:fc-general-threshold}--\eqref{I:R:eq:fc-scalar-Psi} identify it as the
projection optimum $R_{1,1}(q)$. The previously proved rank-one output
completion theorem supplies a real rank-one Hankel attaining this relaxation,
so its value is the exact real and complex F1 optimum.

For an algebraic comparison with the selected quartic, put
\[
 \mathcal L_q(\tau)=(1+q)(\tau-q)-(1-q)\tau^2,\qquad
 \mathcal A_q(\tau)=(1-q)^2\tau^2+(q^2+2q-1)(\tau-q).
\]
The radical equation implies
$\mathcal L_q(\tau)=2(\tau-q)\sqrt{q(1-\tau)}>0$, and direct expansion gives
\begin{align*}
 &\mathcal A_q(\tau)^2-4q^3(\tau-q)(1-2q+q\tau)\\
 &\qquad=(1-q)^2\bigl[\mathcal L_q(\tau)^2
                        -4q(1-\tau)(\tau-q)^2\bigr].
\end{align*}
Thus \eqref{I:R:eq:fc-F1-root} provides a monotone, unambiguous representation of
the already selected quartic branch.
\end{proof}


This benchmark also detects the limitation of keeping only the dominant
singular direction. In scalar dimension its local envelope is
$h_{1,1}(t)=t^2/4$, and exact maximization gives
\[
 P_{1,1}(q)=\frac{q}{(1+\sqrt q)^2},\qquad
 t_* =\frac{2\sqrt q}{1+\sqrt q}.
\]
The corresponding lower bound $L_{\mathrm{dom}}=1-P_{1,1}$ agrees with the
exact F1 value through $q^2$, but
\begin{align*}
 D_1(q)&=1-q+2q^{3/2}-3q^2+5q^{5/2}-8q^3+O(q^{7/2}),\\
 L_{\mathrm{dom}}(q)&=1-q+2q^{3/2}-3q^2+4q^{5/2}-5q^3+O(q^{7/2}),\\
 D_1(q)-L_{\mathrm{dom}}(q)&=q^{5/2}-3q^3+O(q^{7/2}).
\end{align*}
The full-channel threshold therefore passes an exact benchmark that the
dominant-vector envelope fails at higher order. Its equality with the Hankel
optimum has been established here only for F1; higher-state matching
constructions remain separate problems.

% End full_channel.tex


% Begin terminal_languages.tex

\section{Arbitrary terminal languages: an exact boundary reduction}
\label{I:R:sec:terminal-languages}

The waiting-letter mechanism extends beyond a monomial suffix.  The
appropriate scalar approximation problem is replaced by simultaneous
approximation of all terminal residuals in one common invariant space.
This section proves that extension, without a uniqueness or a
nondegeneracy assumption on its terminal extremizers.

Let $\Delta$ be a finite alphabet, let $a\notin\Delta$, and write
$\Sigma=\{a\}\sqcup\Delta$.  Let
$\tau:\Delta^*\to\RC$ have bounded counting-measure word-Hankel operator
$H_\tau$.  In particular $\tau\in\ell^2(\Delta^*)$; a polynomial
terminal language always satisfies this hypothesis.  For
$q=1-r^2\in(0,1)$ define
\begin{equation}
 f_{r,\tau}(a^j w)=qr^j\tau(w)
 \quad(j\ge0,\ w\in\Delta^*),
 \qquad f_{r,\tau}=0\text{ on the remaining words}.
 \label{I:R:eq:terminal-target}
\end{equation}
The parsing in \eqref{I:R:eq:terminal-target} is unique.  Fix a state budget
$n\ge1$ and put $\tau_j(v)=\tau(jv)$ for $j\in\Delta$.
Let $\mathfrak K_{n-1}$ be the collection of subspaces of
$\ell^2(\Delta^*)$ of dimension at most $n-1$ invariant under every
left backward shift.  Define
\begin{equation}
 E_{n-1}(\tau)^2=
 \inf_{K\in\mathfrak K_{n-1}}
       \sum_{j\in\Delta}\Rnorm{\tau_j-P_K\tau_j}_2^2.
 \label{I:R:eq:terminal-E}
\end{equation}
The same $K$ is used for every $j$; separate residual fits do not give
the stated state budget.  When $n=1$, use $K=\{0\}$.

\subsection{Exact elimination and dimension loss}

We use the canonical backward-shift convention
$B_w=B_{w_s}\cdots B_{w_1}$ for $w=w_1\cdots w_s$, and
$B_\varepsilon=I$.  Thus coordinate evaluation is $\ell B_wx$.
For a passive tuple
\begin{equation}
 A^*A+\sum_{j\in\Delta}B_j^*B_j+\ell^*\ell\preceq I_n,
 \qquad V=\begin{pmatrix}(B_j)_{j\in\Delta}\\\ell\end{pmatrix},
 \label{I:R:eq:terminal-passive}
\end{equation}
set
\[
 T_\tau(B,\ell)=\sum_{w\in\Delta^*}
          \overline{\tau(w)}\,\ell B_w,
 \qquad
 h_{\tau,n}(t)=
 \max_{\substack{V^*V\preceq I_n\\\Rnorm{Ve_1}^2=t}}
                  |T_\tau(B,\ell)e_1|^2.
\]
The observation map $Ox=(\ell B_wx)_w$ is contractive by telescoping
the passive energy inequality.  In fact
\begin{equation}
 \sum_w|\ell B_wy|^2
 \le |\ell y|^2+\sum_j\Rnorm{B_jy}^2=\Rnorm{Vy}^2.
 \label{I:R:eq:terminal-refined-energy}
\end{equation}
Consequently the series defining $T_\tau$ converges in operator norm
uniformly over all passive tuples, and
$0\le h_{\tau,n}(t)\le\Rnorm\tau_2^2t$.  Compactness then gives the
maxima in this definition and below.

\begin{lemma}[Universal waiting-transition elimination]
\label{I:R:lem:terminal-elimination}
For every $q\in(0,1)$,
\begin{align}
 P_{\tau,n}(q)
 &:=\max_{\eqref{I:R:eq:terminal-passive}}
       q\Rnorm{T_\tau(B,\ell)(I-rA)^{-1}}^2\notag\\
 &=\max_{0\le t\le1}
       \frac{q h_{\tau,n}(t)}{(1-r\sqrt{1-t})^2}.
 \label{I:R:eq:terminal-elimination}
\end{align}
\end{lemma}
\begin{proof}
For a unit vector $y$, put $D=I-V^*V$ and
$t=\Rnorm{Vy}^2=1-y^*Dy$.  The substitution $x=(I-rA)y$ in the row
operator norm and the passive inequality give
\[
 \Rnorm{(I-rA)y}\ge1-r\Rnorm{Ay}\ge1-r\sqrt{1-t}.
\]
A unitary state change sends $y$ to $e_1$ and preserves the terminal
row and constraints.  This proves the upper direction of
\eqref{I:R:eq:terminal-elimination}.  Conversely, for fixed $V,y$ and
$t<1$, choose
\[
 A=\frac{yy^*D}{\sqrt{1-t}}.
\]
Then $Ay=\sqrt{1-t}\,y$ and
$A^*A=Dyy^*D/(y^*Dy)\preceq D$, by Cauchy--Schwarz in the $D$
metric.  Both denominator inequalities are equalities.  When $t=1$
choose $A=0$.  This proves exact attainability of the scalar
denominator.
\end{proof}


An exact factorization makes the localization transparent.  For
$\rho(x,y)=|x-y|/|1-xy|$ on the real unit interval, put
\begin{equation}
 K_q(t)=\frac{qt}{(1-\sqrt{(1-q)(1-t)})^2}
       =1-\rho\bigl(\sqrt{1-q},\sqrt{1-t}\bigr)^2.
 \label{I:R:eq:terminal-kernel}
\end{equation}
Thus $0\le K_q(t)\le1$ and $K_q(q)=1$.  This is just
$(1-r\sqrt{1-t})^2-qt=(r-\sqrt{1-t})^2$.

\begin{lemma}[The residual dimension-drop coefficient]
\label{I:R:lem:terminal-dimension-drop}
As $t\downarrow0$,
\begin{equation}
 h_{\tau,n}(t)=a_{\tau,n}t+o(t),
 \qquad
 a_{\tau,n}=|\tau(\varepsilon)|^2+
       \max_{K\in\mathfrak K_{n-1}}
             \sum_{j\in\Delta}\Rnorm{P_K\tau_j}_2^2.
 \label{I:R:eq:terminal-leading-coefficient}
\end{equation}
The maximum is attained, and $P_{\tau,n}(q)\to a_{\tau,n}$.
\end{lemma}
\begin{proof}
Write $t=s^2$ and let $N=|\Delta|n+1$.  Every feasible $V$ has the
exact compact parametrization
\begin{equation}
 V=[sz,W_s(z)Z],\qquad
 \Rnorm z=1,\quad Z^*Z\preceq I_{n-1},\qquad
 W_s(z)=I_N-(1-\sqrt{1-s^2})zz^*.
 \label{I:R:eq:terminal-parametrization}
\end{equation}
Indeed this is the Schur complement of $I-V^*V$.  For a polynomial
$\tau$, the quotient $T_\tau e_1/s$ has an analytic extension to
$s=0$.  For general $\tau\in\ell^2$, the tail bound from
\eqref{I:R:eq:terminal-refined-energy}, divided by $|s|$, is at most
the terminal tail's $\ell^2$ norm.  Polynomial approximation therefore
gives a uniform continuous extension on the compact parameter domain.

At $s=0$ write
\[
 B_j(0)=\begin{pmatrix}0&b_j\\0&C_j\end{pmatrix},\qquad
 \ell(0)=(0,l),\qquad z_j=(\alpha_j,\xi_j),
\]
where $z_j$ is the first-column derivative block for $B_j$, and let
$z_o$ be its scalar output block.  The limiting divided terminal row is
\begin{equation}
 \overline{\tau(\varepsilon)}z_o+
 \sum_j R_j(C,l)\xi_j,
 \qquad
 R_j(C,l)=\sum_v\overline{\tau(jv)}\,lC_v.
 \label{I:R:eq:terminal-divided-limit}
\end{equation}
The $\alpha_j$ do not contribute.  Maximizing over unit $z$ gives
$|\tau(\varepsilon)|^2+\sum_j\Rnorm{R_j(C,l)}^2$.
The remaining columns in \eqref{I:R:eq:terminal-parametrization} imply
$\sum_j C_j^*C_j+l^*l\preceq I_{n-1}$.  Conversely every such reduced
tuple is feasible by choosing $b_j=0$.  Uniform convergence of the
compact objective proves
\[
 a_{\tau,n}=|\tau(\varepsilon)|^2+
 \max_{\sum C_j^*C_j+l^*l\preceq I_{n-1}}
                    \sum_j\Rnorm{R_j(C,l)}^2.
\]
For $n=1$ the reduced state space is zero and the sum is zero.

This passive maximum equals the invariant-space maximum in
\eqref{I:R:eq:terminal-leading-coefficient}.  A canonical observation map
on $K$ is its isometric inclusion, so $R_j$ has norm
$\Rnorm{P_K\tau_j}$.  Conversely a passive reduced tuple has a
contractive observation map $O$ whose finite-dimensional range $K$ is
backward invariant.  Since $OO^*\preceq P_K$,
$\Rnorm{O^*\tau_j}\le\Rnorm{P_K\tau_j}$.  This gives the reverse
inequality after summing.  A maximizing passive tuple exists by
compactness and uniform convergence; its range attains the
invariant-space maximum as well.

Finally write \eqref{I:R:eq:terminal-elimination} as
$P_{\tau,n}(q)=\max_t K_q(t)h_{\tau,n}(t)/t$.
The trial $t=q$ tends to $a_{\tau,n}$.  On $t\ge\eta>0$ the kernel
$K_q(t)$ tends uniformly to zero, whereas on $0<t<\eta$ it is at most
one and $h(t)/t\to a_{\tau,n}$.  Splitting at $\eta$ and then sending
$\eta$ to zero proves the upper limit.  This also covers
$a_{\tau,n}=0$.
\end{proof}


\subsection{The exact boundary theorem and an explicit lift}

\begin{theorem}[Simultaneous terminal approximation]
\label{I:R:thm:terminal-boundary}
For every $n\ge1$ and terminal series satisfying the stated boundedness
hypothesis,
\begin{equation}
 \boxed{\displaystyle
 \lim_{r\uparrow1}\delta_n^{\RC}(H_{f_{r,\tau}})^2
       =E_{n-1}(\tau)^2.}
 \label{I:R:eq:terminal-boundary-theorem}
\end{equation}
The infimum in \eqref{I:R:eq:terminal-E} is attained.  For real $\tau$, the
same statement holds separately for real competitors and real
invariant spaces.  Equality of the two fields' constants is not
assumed.
\end{theorem}
\begin{proof}
The unit vector
$u_r=\sqrt q\sum_{m\ge0}r^m e_{a^m}$ is mapped by
$H_{f_{r,\tau}}^{\mathsf T}$ to
\[
 v_r=\sqrt q\sum_{m,w}r^m\tau(w)e_{a^m w},\qquad
 \Rnorm{v_r}=\Rnorm\tau_2.
\]
For a bounded rank-at-most-$n$ competitor, its transposed range $S$ is
backward invariant.  Canonical orthonormal state coordinates yield
\eqref{I:R:eq:terminal-passive} and
\[
 \Rnorm{P_Sv_r}^2
   =q\Rnorm{T_\tau(B,\ell)(I-rA)^{-1}}^2.
\]
Orthogonal projection thus gives
\[
 \delta_n^{\RC}(H_{f_{r,\tau}})^2
   \ge\Rnorm\tau_2^2-P_{\tau,n}(q).
\]
Lemma~\ref{I:R:lem:terminal-dimension-drop} proves the lower boundary
bound in \eqref{I:R:eq:terminal-boundary-theorem}, because
$\Rnorm\tau_2^2=|\tau(\varepsilon)|^2+\sum_j\Rnorm{\tau_j}_2^2$.

For the upper bound choose a maximizing space $K$ in
Lemma~\ref{I:R:lem:terminal-dimension-drop}.  Put
\[
 h(\varepsilon)=\tau(\varepsilon),\qquad
 h(jv)=(P_K\tau_j)(v),\qquad e=\tau-h.
\]
All residuals $x_j=P_K\tau_j$ share $K$.  With $C_j$ the canonical
backward shifts on $K$ and $l$ evaluation at the empty word, the usual
forward realization of the lifted series
$g_r(a^m w)=qr^m h(w)$ uses at most $1+\dim K\le n$ states:
\begin{equation}
 A_a=\begin{pmatrix}r&0\\0&0\end{pmatrix},\qquad
 A_j=\begin{pmatrix}0&x_j^{\mathsf T}\\0&C_j^{\mathsf T}\end{pmatrix},
 \quad\alpha^*=q(1,0),\quad
 \beta=\begin{pmatrix}\tau(\varepsilon)\\l^{\mathsf T}\end{pmatrix}.
 \label{I:R:eq:terminal-lift}
\end{equation}
These are coefficient transposes, not Hermitian adjoints.  Products
give $h(jv)=lC_vx_j$ with the reversed canonical word order.  A letter
$a$ following a terminal letter makes the product zero, as required.
The common terminal tuple is jointly strictly stable: in its
finite-dimensional shift-invariant space, finite-horizon tail energies
tend uniformly to zero.  Thus its lifted Hankel is bounded.

Here is also a direct infinite-operator check.  Let $U_r$ be the
isometry from $\ell^2(\Delta^+)$ given by
$U_re_w=\sqrt q\sum_mr^m e_{a^m w}$, let $W$ embed
$\ell^2(\Delta^*)$ as words without $a$, and let
$K_e[w,v]=e(wv)$ with $w\ne\varepsilon$.  Since
$e(\varepsilon)=0$, direct entry matching gives
\begin{equation}
 H_{f_{r,\tau}}-H_{g_r}
   =u_r(U_r\overline e)^*+\sqrt q\,U_rK_eW^*.
 \label{I:R:eq:terminal-error-decomposition}
\end{equation}
The two output ranges are orthogonal.  Also $K_e$ is bounded: the
terminal Hankel of $\tau$ is bounded by hypothesis, and that of $h$ is
bounded by its stable finite realization.  Consequently
\[
 \Rnorm{H_{f_{r,\tau}}-H_{g_r}}^2
   \le\Rnorm e_2^2+q\Rnorm{K_e}^2
   =E_{n-1}(\tau)^2+q\Rnorm{K_e}^2.
\]
This proves the upper boundary bound and attainment.  The argument
over real coefficients uses orthogonal state changes and real spaces
throughout.
\end{proof}


\begin{proposition}[A fixed optimal lift has no positive linear error]
\label{I:R:prop:terminal-fixed-lift}
For the fixed lift above, suppose $E=\Rnorm e_2>0$ and put
$M=\Rnorm{K_e}$.  Then
\begin{equation}
 \Rnorm{H_{f_{r,\tau}}-H_{g_r}}^2
 =E^2+\frac{\Rnorm{K_e\overline e}^2}{E^2}q^2+O(q^3).
 \label{I:R:eq:terminal-fixed-lift-expansion}
\end{equation}
If $E=0$, the lift equals the target and its error is zero.
\end{proposition}
\begin{proof}
Put $v=U_r\overline e$ and $R=K_e^*K_e$.  Equation
\eqref{I:R:eq:terminal-error-decomposition} gives
$H_{\rm err}^*H_{\rm err}=vv^*+qWRW^*$, while
$\Rnorm v=E$ and $W^*v=\sqrt q\,\overline e$.
In the split $\operatorname{span}\{v\}\oplus v^\perp$, the
upper-left perturbation is
$a_q=q^2\Rnorm{K_e\overline e}^2/E^2$, the off-diagonal block has
norm at most $q^{3/2}M^2$, and the lower block has norm at most $qM^2$.
For $qM^2\le E^2/2$, the Schur complement of the isolated top
eigenvalue yields
\[
 E^2+a_q\le\Rnorm{H_{\rm err}}^2
 \le E^2+a_q+2q^3M^4/E^2.
\]
This proves \eqref{I:R:eq:terminal-fixed-lift-expansion}.  Complex
coefficient conjugation in $K_e\overline e$ is essential.
\end{proof}


This proposition is an upper construction.  A $q$-dependent lift can
have a negative linear correction, as the original two-state example
demonstrates; no matching first-order lower coefficient is inferred
from the fixed lift.

\subsection{An exact mixed-terminal example}

Take $\Delta=\{b\}$, $n=2$, and let
$\tau(b^2)=\beta$, $\tau(b^3)=\gamma$, all other coefficients zero,
where $\beta,\gamma\ge0$ and $\beta^2+\gamma^2=1$.
The reduced tuple in Lemma~\ref{I:R:lem:terminal-dimension-drop} is scalar,
so phase alignment gives
\[
 a_{\tau,2}=\max_{0\le\rho\le1}
          (1-\rho^2)\rho^2(\beta+\gamma\rho)^2.
\]
For $\beta,\gamma>0$ its unique maximizing point satisfies
\[
 \beta(1-2\rho^2)+\gamma\rho(2-3\rho^2)=0.
\]
Uniqueness follows from strict concavity of the logarithm of the
positive objective on $(0,1)$.  When
$\beta=\gamma=1/\sqrt2$,
\begin{gather}
 \rho=\frac{1+\sqrt{13}}6,\qquad
 a_{\tau,2}=\frac{587+143\sqrt{13}}{2916},\notag\\
 \boxed{\displaystyle
 \lim_{r\uparrow1}\delta_2(H_{f_{r,\tau}})^2
 =\frac{2329-143\sqrt{13}}{2916}
 =0.62188140178623097774\ldots .}
 \label{I:R:eq:terminal-mixed-example}
\end{gather}
A real scalar reduced tuple attains the complex maximum here, so the
real and complex limits agree for this example.  This is a boundary
value for a nonmonomial terminal language, not an exact finite-$r$
optimization statement.

\subsection{A sharp two-state correction for every terminal degree}
\label{I:R:subsec:all-terminal-degrees}

The same complex phase calculation also gives a family of sharp
first-order results with a fixed two-state budget and arbitrary
terminal degree.  This extension changes the target family rather than
the state budget.

\begin{theorem}[All terminal degrees with two states]
\label{I:R:thm:two-state-all-degrees}
Fix an integer $d\ge2$ and let
$f_{r,d}(a^j b^d)=qr^j$, zero elsewhere.  Put
\[
 a_d=\frac{(d-1)^{d-1}}{d^d}.
\]
For each field separately,
\begin{equation}
 \boxed{\displaystyle
 \delta_2^{\mathbb F}(H_{f_{r,d}})^2
   =1-a_d-\frac{a_d}{(d-1)^2}q+O(q^2),
 \qquad\mathbb F\in\{\mathbb R,\RC\}.}
 \label{I:R:eq:two-state-all-degrees}
\end{equation}
The constants may depend on fixed $d$.  In particular $d=2$ gives
$3/4-q/4+O(q^2)$, whereas $d=3$ gives
$23/27-q/27+O(q^2)$.
\end{theorem}
\begin{proof}
We first prove the lower bound globally over complex tuples.  Write
\[
 h_{2,d}(t)=\max_{\substack{[B;\ell]^*[B;\ell]\preceq I_2\\
                       \Rnorm{[B;\ell]e_1}^2=t}}
                         |\ell B^d e_1|^2.
\]
In the exact parametrization
$[B;\ell]=[sz,W_s(z)\nu]$ put
$\mathcal J_d=\ell B^de_1/s$.  At $s=0$,
\[
 \mathcal J_d=z_2\nu_3\nu_2^{d-1},\qquad
 \max|\mathcal J_d|=M_0=lc^{d-1}=\sqrt{a_d},\qquad
 c=\sqrt{\frac{d-1}{d}},\quad l=d^{-1/2}.
\]
Equality holds precisely when $z_1=z_3=\nu_1=0$,
$|z_2|=1$, $|\nu_2|=c$, and $|\nu_3|=l$.
The exact symmetries
$B\mapsto e^{i\theta}U^*BU$, $\ell\mapsto e^{i\phi}\ell U$,
$U=\operatorname{diag}(1,e^{i\psi})$, change the phases of
$(z_2,\nu_2,\nu_3)$ by
$(\theta-\psi,\theta,\phi+\psi)$.  This invertible phase map fixes
$z_2,\nu_2,\nu_3$ positive near the maximizing torus without discarding
any complex transverse direction.  The outward
radial derivative in $\nu$ is $dM_0>0$, so maximizers have
$\Rnorm\nu=1$ for small $|s|$.

Let $\theta_0=\arccos c$.  In the resulting seven real coordinates,
take $z_1,z_3,\nu_1\in\RC$, $\xi\in\mathbb R$, and
\[
 z_2=\sqrt{1-|z_1|^2-|z_3|^2},\quad
 \nu_2=\sqrt{1-|\nu_1|^2}\cos(\theta_0+\xi),\quad
 \nu_3=\sqrt{1-|\nu_1|^2}\sin(\theta_0+\xi).
\]
The analytic objective $F_d=|\mathcal J_d|$ has the expansion
\begin{align}
 F_d(s,x)=M_0
 &-\frac{d-1}{2}M_0s^2
   +\frac{M_0}{c}s\Re z_1
   +\frac{(d-2)M_0}{c^2}s\Re\nu_1\notag\\
 &-\frac{M_0}{2}(|z_1|^2+|z_3|^2)
   -\frac{dM_0}{2}|\nu_1|^2-dM_0\xi^2
   +O((|s|+\Rnorm x)^3).
 \label{I:R:eq:degree-d-Hessian}
\end{align}
For clarity, the pure $s^2$ term follows from the exact value
$F_d(s,0)=M_0(1-s^2)^{(d-1)/2}$.  The mixed $z_1$ term is the path
with one extra waiting step in the first state.  There are $d-2$
paths producing the mixed $\nu_1$ term: one return from the second
state to the first, followed immediately by re-entry.  The angular
second logarithmic derivative of
$\cos^{d-1}\theta\sin\theta$ at $\theta_0$ is $-2d$.
These observations give every coefficient in
\eqref{I:R:eq:degree-d-Hessian}.

Its transverse Hessian is negative definite.  Compact localization,
the complete phase gauge, and the analytic implicit function theorem
therefore identify its local maximum with the global maximum for
small $|s|$.  The optimizing first derivatives are
\[
 \Re z_1=s/c+O(s^2),\qquad
 \Re\nu_1=\frac{d-2}{dc^2}s+O(s^2),
\]
with all other transverse coordinates $O(s^2)$.  The maximum modulus's
$s^2$ coefficient is
\[
 \frac{M_0}{2}\left[-(d-1)+\frac1{c^2}
                         +\frac{(d-2)^2}{dc^4}\right]
   =\frac{M_0}{2(d-1)^2}.
\]
Physical evenness removes the odd analytic coefficients.  Hence
\begin{equation}
 h_{2,d}(t)=a_dt+\frac{a_d}{(d-1)^2}t^2+O(t^3),\qquad
 P_{2,d}(q)=a_d+\frac{a_d}{(d-1)^2}q+O(q^2).
 \label{I:R:eq:two-state-all-degree-envelope}
\end{equation}
The second equality follows by scalar localization in
Lemma~\ref{I:R:lem:terminal-elimination}; its leading profile is
$4a_dx/(1+x)^2$ for $t=qx$, with a unique quadratic maximum at $x=1$.
Projection gives the claimed complex lower bound.

For completeness the matching real trial is explicit.  With
$s=\sqrt q$, set
\[
 \lambda=1/c,\quad w=\frac{d-2}{d-1},\quad
 F_s=\sqrt{1-s^2w^2},\quad
 N_s=\begin{pmatrix}sw\\cF_s\\lF_s\end{pmatrix},\quad
 z_s=\begin{pmatrix}s\lambda\\\sqrt{1-s^2\lambda^2}\\0\end{pmatrix}.
\]
For sufficiently small $s>0$, define
\[
 V_s=[sz_s,W_s(z_s)N_s]=\begin{pmatrix}B_s\\\ell_s\end{pmatrix},
 \quad v_s=\begin{pmatrix}r\\-sN_s^{\mathsf T}z_s\end{pmatrix},
 \quad A_s=e_1v_s^{\mathsf T}.
\]
Then $N_s^{\mathsf T}N_s=1$, $\Rnorm{z_s}=1$, and
$A_s^{\mathsf T}A_s+B_s^{\mathsf T}B_s+
\ell_s^{\mathsf T}\ell_s=I_2$ exactly.  Define
\[
 \omega_s=s\ell_sB_s^d(I-rA_s)^{-1},\qquad
 b_s=s^{-1}(I-rA_s)\omega_s^{\mathsf T},\qquad
 g_s(u)=\ell_sT_{\operatorname{rev}(u)}b_s.
\]
Here $T_u$ is the standard ordered product of $A_s,B_s$ along $u$.
This is a real two-state series.  The structured first-order matching
argument applies with terminal power $d$: its required cancellation is
\[
 l c^{d-1}=M_0,\qquad
 l c^d-M_0c=0.
\]
In particular
$\omega_s=(M_0+O(q),O(s^3))$, $b_s=(sM_0+O(s^3),O(q))$.
The exact energy identity and the block orders
\[
 B_s=\begin{pmatrix}O(q)&O(s)\\s+O(s^3)&c+O(q)\end{pmatrix},
 \qquad
 A_s=\begin{pmatrix}r&-sc+O(s^3)\\0&0\end{pmatrix}
\]
give joint stability for small positive $s$: the unique transition
eigenvalue near one is $1-q+O(q^2)$ and all other eigenvalues remain
strictly smaller in modulus.  The reachability Gram has slow entry
$M_0^2+O(q)$, off-diagonal $O(s^3)$ and fast entry $O(q)$.
As in the structured matching proof, interpolation on the dominant
input is exact, the error's squared dominant Rayleigh quotient is
$1-\Rnorm{\omega_s}^2$, its off-diagonal block is $O(q)$, and its
remaining block is $O(q)$.  The fixed spectral gap $1-a_d>0$ gives
\[
 \Rnorm{H_{f_{r,d}}-H_{g_s}}^2
       =1-\Rnorm{\omega_s}^2+O(q^2).
\]
The chosen $\lambda,w$ realize the quadratic maximum in
\eqref{I:R:eq:degree-d-Hessian}, so
$\Rnorm{\omega_s}^2=a_d+a_dq/(d-1)^2+O(q^2)$.
This matches the global complex lower bound by a real construction,
proving \eqref{I:R:eq:two-state-all-degrees} for each field.  Equality of
their full finite-$q$ optimum functions is not asserted.
\end{proof}


% End terminal_languages.tex

\section{Scalar input for the response examples}
\label{I:R:app:Hardy-bridge}
The examples use the single Hardy bridge proved in
Theorem~\ref{I:B:thm:hardy-bridge}; its proof is not repeated here.
The lower bound is a direct H\"older inequality, the equality for
$m\le3$ uses published low-degree extremals, and equality in general
uses the external outer-extremal theorem~\cite{outer2026}.

\section{Balancing and certified dyadic recertification}
\label{I:U:sec:dyadic-recertification}
A spectral-radius margin by itself does not supply the resolvent estimate
needed for a spectral truncation bound. Balancing remains useful, provided
the reduced candidate is recertified. The following statement separates
these two steps.

\begin{theorem}[Balancing and dyadic recertification]
\label{I:U:thm:dyadic-recertification}
\begin{enumerate}[label=(\roman*)]
\item A minimal mean-square-stable realization with positive definite
Gramians admits a balancing transformation
\[
 S^{-1}P_fS^{-*}=S^*Q_fS=\Sigma_H.
\]
A leading-block truncation whose mean-square stability has been verified
is a valid candidate of the stated dimension. Its error can be certified
by Theorem~\ref{I:U:thm:exact}; no additional spectral-tail formula is asserted
here.
\item Let $g=(T_a,b,c)$ have real algebraic coefficients and be
mean-square stable. Choose a positive definite algebraic $X$ with
$X-\sum_aT_aXT_a^*=I$. For perturbed matrices
$\widehat T_a=T_a+E_a$, define
\[
 \theta=\left\|\sum_a(T_aXE_a^*+E_aXT_a^*+E_aXE_a^*)\right\|.
\]
If a certified bound gives $\theta<1$, then
\begin{equation}
 X-\sum_a\widehat T_aX\widehat T_a^*\succeq(1-\theta)I,
 \qquad
 \rho(\Phi_{\widehat T})\le1-\frac{1-\theta}{\|X\|}<1.
 \label{I:U:eq:dyadic-verified-margin}
\end{equation}
An elementary sufficient upper bound for $\theta$ is
\[
 \sum_a\left(2\|T_aX^{1/2}\|\|E_aX^{1/2}\|
                         +\|E_aX^{1/2}\|^2\right).
\]
\item For every positive rational $\varepsilon$, an algorithm can
increase dyadic precision, round all entries of $T_a,b,c$, and return a
stable dyadic realization $\widehat g$ with
$\|H_g-H_{\widehat g}\|\le\varepsilon$. Each candidate is accepted
only after verifying \eqref{I:U:eq:dyadic-verified-margin} and a joint
Gramian error certificate. This search terminates. This statement gives
no uniform polynomial bound on the required precision or search time.
\end{enumerate}
\end{theorem}
\begin{proof}
For (i), diagonalize
$P_f^{1/2}Q_fP_f^{1/2}=U\Sigma_H^2U^*$ and take
$S=P_f^{1/2}U\Sigma_H^{-1/2}$. Direct substitution proves both
identities. Stability of a proposed truncation is an explicit condition,
not a consequence asserted by this proof.

For (ii), subtract the perturbation from the Stein equality. The
operator-norm bound gives the first inequality in
\eqref{I:U:eq:dyadic-verified-margin}. Since $X\succeq I$ and
$I\succeq X/\|X\|$, the positive map $\Phi_{\widehat T}$ contracts
the order-unit norm induced by $X$ by at most the displayed factor.
This bounds its spectral radius. Factoring each cross term through
$X^{1/2}$ proves the sufficient bound.

For (iii), choose dyadic approximations converging entrywise to the
algebraic input. Their $\theta$ tends to zero, so all sufficiently fine
approximations are stable. Form the joint error realization
\[
 B_a=\Udiag(T_a,\widehat T_a),\qquad
 \beta=(b^\top,\widehat b^\top)^\top,\qquad
 \gamma=(c^\top,-\widehat c^\top)^\top.
\]
Its Gramians $P_E,Q_E$ are obtained from the two joint Stein equations.
The block diagonal matrix $\Udiag(X,X)$ supplies a strict Lyapunov
inequality in a neighborhood of the unperturbed tuple. Thus the joint
Stein inverses and $\operatorname{tr}(P_EQ_E)$ depend continuously on
the perturbation there. At zero perturbation the error series vanishes,
so this trace is zero. Consequently it eventually becomes strictly
smaller than $\varepsilon^2$. Algebraic arithmetic can certify this
strict inequality and the Lyapunov margin, and
\[
 \|H_g-H_{\widehat g}\|^2
 \le\|H_g-H_{\widehat g}\|_{\rm HS}^2
 =\operatorname{tr}(P_EQ_E)<\varepsilon^2.
\]
Enumerating increasing precisions with these tests therefore terminates.
\end{proof}


\paragraph{Why a radius margin is not a resolvent norm.}
For $A_t=\left(\begin{smallmatrix}0&t\\0&0\end{smallmatrix}\right)$,
the positive map $\Phi(X)=A_tXA_t^*$ has $\rho(\Phi)=0$, but
\[
 (I-\Phi)^{-1}(I)=I+A_tA_t^*
       =\Udiag(1+t^2,1).
\]
Its induced operator norm is at least $1+t^2$, although the inverse
spectral-radius margin equals one. Thus a bound by
$\max(\gamma^{-1},\gamma_k^{-1})$ is not a general resolvent estimate.
This counterexample to that proof step does not by itself disprove its
separate balanced-truncation inequality. That inequality remains
unestablished here. A uniform spectral-tail guarantee and a complete efficient synthesis theorem remain separate questions.


\subsection{Quantitative truncation with two verified Stein witnesses}
\label{I:U:sec:quantitative-truncation}
The following result supplies one explicit answer to the quantitative
truncation target. Its constants are allowed to deteriorate when the
source approaches the obstruction of Theorem~\ref{I:U:thm:separation}.
It uses full matrix inequalities, not a spectral-radius margin.

\begin{theorem}[Conditional square-root bound for balanced truncation]
\label{I:U:thm:quantitative-truncation}
Let a minimal stable realization $(A_a,b,c)$ be balanced with
$P=Q=\Sigma=\Udiag(\Sigma_1,\Sigma_2)\succ0$. Retain the first $k$
coordinates and put $B_a=A_{a,11}$, $g(w)=c_1^*B_wb_1$.
Assume the tuple $B$ is independently certified mean-square stable.
Suppose positive definite matrices $X,Y$ and $0\le r_X,r_Y<1$ satisfy
\[
 \sum_a A_aXA_a^*\preceq r_X^2X,\qquad
 \sum_a A_a^*YA_a\preceq r_Y^2Y.
\]
Write $\kappa_X=\|X\|/\lambda_{\min}(X)$ and similarly for $Y$.
Then
\begin{equation}
 \|H_f-H_g\|\le\sqrt{\sigma_1\Utr\Sigma_2}
 \left(\frac{\sqrt{\kappa_X}}{1-r_X}
             +\frac{\sqrt{\kappa_Y}}{1-r_Y}\right).
 \label{I:U:eq:quantitative-truncation}
\end{equation}
In particular $\Utr\Sigma_2\le(n-k)\sigma_{k+1}$ gives an explicit
square-root rate with all other displayed parameters fixed.
\end{theorem}
\begin{proof}
Let $E:\mathbb F^k\to\mathbb F^n$ be the coordinate inclusion and
let $C,O$ and $C_B,O_B$ denote reachability and observation maps.
The leading Stein blocks show $P_B,Q_B\preceq\Sigma_1$: for example,
iterate $\Sigma_1\succeq b_1b_1^*+\sum_aB_a\Sigma_1B_a^*$ and
use stability of $B$. Hence $\|O\|=\sqrt{\sigma_1}$ and
$\|C_B\|\le\sqrt{\sigma_1}$.

For $z_w=A_wb-EB_wb_1$, the tree recursion is
\[
 z_\varepsilon=(0,b_2)^\top,\qquad
 z_{aw}=A_az_w+(0,A_{a,21}B_wb_1)^\top.
\]
On the Hilbert sum of state vectors indexed by words, with each state
norm $\|z\|_Y$, the map $(z_w)\mapsto(A_az_w)_{aw}$ has norm at
most $r_Y$. Root and nonroot forcing coordinates are disjoint. Their
squared Euclidean sum is
\[
 \|b_2\|^2+\sum_a\Utr(A_{a,21}P_BA_{a,21}^*)
 \le\Utr\Sigma_2,
\]
by the lower reachability Stein block. The Neumann series on the tree
and equivalence of the state norms therefore give
\[
 \|C-EC_B\|_{\rm HS}
 \le\frac{\sqrt{\kappa_Y\Utr\Sigma_2}}{1-r_Y}.
\]
Apply the identical argument to $(A_a^*,c)$ with witness $X$, reversing
word order, to obtain
$\|O-O_BE^*\|_{\rm HS}\le
\sqrt{\kappa_X\Utr\Sigma_2}/(1-r_X)$. Finally,
\[
 H_f-H_g=O(C-EC_B)+(O-O_BE^*)EC_B.
\]
The operator norm is at most the displayed Hilbert--Schmidt bounds
times $\sqrt{\sigma_1}$, proving the assertion.
\end{proof}

Both witnesses exist for each stable finite tuple, but their quantitative
constants need not be uniform. A witness for the full tuple does not
automatically give a witness for its leading block. For example,
\[
 A=\begin{pmatrix}2&1\\-4&-2\end{pmatrix},\qquad A^2=0,
 \qquad X=I+AA^*,\qquad X-AXA^*=I,
\]
whereas $A_{11}=2$ is unstable. This refutes the proposed inheritance
step. Balanced truncation and singular-value automata have established
antecedents \citep{sva2019}; the theorem above records the precise
conditional operator-norm estimate used here, without a novelty claim.

\begin{theorem}[A sufficient dyadic precision chosen before recertification]
\label{I:U:thm:explicit-dyadic}
Let an $n$-state tuple $(A_a,b,c)$ have witnesses as above. Choose
$r\in(0,1)$ and $\kappa\ge1$ with
$r_X,r_Y\le r$, $\kappa_X,\kappa_Y\le\kappa$, and put $a=1-r$.
Round every scalar entry to a dyadic number with absolute error at most
$h=2^{-b_0}$. For complex entries, replace these scalar-error bounds
by bounds on their complex moduli. Define
\[
 D_b=\frac{2\sqrt{\kappa n}}a+
                  \frac{2\kappa\sqrt d\,n\|b\|}{a^2},\qquad
 D_c=\frac{2\sqrt{\kappa n}}a+
                  \frac{2\kappa\sqrt d\,n\|c\|}{a^2},
\]
\[
 K= D_c\frac{\sqrt\kappa\|b\|}a
       +\frac{2\sqrt\kappa(\|c\|+1)}a D_b.
\]
The choices
\begin{equation}
 h\le\min\left\{n^{-1/2},
       \frac{a}{2\sqrt\kappa\sqrt d\,n},\frac\epsilon K\right\}
 \label{I:U:eq:explicit-dyadic}
\end{equation}
guarantee mean-square stability of the rounded tuple and
$\|H_f-H_{\widehat f}\|\le\epsilon$.
If $K=0$, the last condition is omitted. Rational inputs and rational
upper bounds for these constants give polynomial bit cost in the input
length and $b_0$ for rounding and exact joint-Gramian recertification.
\end{theorem}
\begin{proof}
In either witness norm the tree-operator perturbation has norm at most
$\sqrt\kappa\sqrt d\,nh\le a/2$. Thus both rounded tree operators
have norm at most $1-a/2$, proving stability. The resolvent identity,
$\|\delta b\|,\|\delta c\|\le\sqrt n h$, and the original and
rounded resolvent bounds $a^{-1}$ and $2a^{-1}$ give
\[
 \|\widehat C-C\|_{\rm HS}\le D_bh,\qquad
 \|\widehat O-O\|_{\rm HS}\le D_ch.
\]
Indeed, the two terms for reachability are the changed root vector
and the changed tree operator applied to the original reachability
sequence; their bounds are exactly the two terms of $D_bh$.
The dual argument gives $D_ch$. Also
$\|C\|\le\sqrt\kappa\|b\|/a$ and
$\|\widehat O\|\le2\sqrt\kappa(\|c\|+1)/a$.
Use $\widehat H-H=(\widehat O-O)C+\widehat O(\widehat C-C)$.
This proves the bound $Kh$. Rational Stein systems and their joint
error systems have polynomial dimension; exact linear algebra and
positive-semidefinite tests have polynomial bit cost. This is a
bound for specified witnesses, not for finding a well-conditioned
balancing transform or a uniformly well-conditioned truncation.
\end{proof}


\subsection{A second truncation certificate from positive Stein resolvents}
\label{I:U:sec:resolvent-truncation}
The extended proposal supplies a useful alternative to the tree estimate
above. Its positive-resolvent argument is valid after correcting the
block inequalities. This subsection retains that argument and exploits
the commuting reference Gramians; it does not supply a uniform AAK factor.

\begin{lemma}[Hermitian norm of a positive Stein resolvent]
\label{I:U:lem:positive-resolvent}
Let $\Phi(Z)=\sum_a T_aZT_a^*$ have spectral radius less than one,
and put $\mathcal R=(I-\Phi)^{-1}$ and $X=\mathcal R(I)$.
For every Hermitian $M$,
\[
 -\|M\|X\preceq\mathcal R(M)\preceq\|M\|X,
 \qquad \|\mathcal R(M)\|\le\|X\|\,\|M\|.
\]
Thus the induced spectral norm on the Hermitian subspace is exactly
$\|X\|$, with equality at $M=I$.
\end{lemma}
\begin{proof}
Apply the positive map $\sum_{j\ge0}\Phi^j$ to
$-\|M\|I\preceq M\preceq\|M\|I$. The order bounds imply the norm
bound, and $M=I$ attains it. This is the elementary Hermitian case of
the classical norm-at-identity property for positive maps
\citep{bourin2019}, not a new general operator-theory result.
\end{proof}


\begin{theorem}[Commuting-reference certificate for balanced truncation]
\label{I:U:thm:reference-truncation}
Let $(A_a,b,c)$ be minimal, mean-square stable, and balanced with
$P=Q=\Sigma=\Udiag(\Sigma_1,\Sigma_2)\succ0$, where $1\le k<n$
coordinates are retained. Set $B_a=A_{a,11}$ and
$g(w)=c_1^*B_wb_1$, and independently assume $B$ is mean-square
stable. Write $\delta=\|\Sigma_2\|=\sigma_{k+1}$ and
$\sigma_1=\|\Sigma\|$. Define the four Stein solutions
\[
 X_A-\sum_aA_aX_AA_a^*=I_n,\quad
 X_B-\sum_aB_aX_BB_a^*=I_k,
\]
\[
 Y_A-\sum_aA_a^*Y_AA_a=I_n,\quad
 Y_B-\sum_aB_a^*Y_BB_a=I_k,
\]
and $K_P=\max(\|X_A\|,\|X_B\|)$,
$K_Q=\max(\|Y_A\|,\|Y_B\|)$.
On the joint space with block sizes $(k,n-k,k)$ put
\[
 T_a=\Udiag(A_a,B_a),\quad \beta=(b_1,b_2,b_1)^\top,
 \quad\chi=(c_1,c_2,-c_1)^\top,
\]
\[
 P_0=\begin{pmatrix}\Sigma_1&0&\Sigma_1\\0&\Sigma_2&0\\
 \Sigma_1&0&\Sigma_1\end{pmatrix},\qquad
 Q_0=\begin{pmatrix}\Sigma_1&0&-\Sigma_1\\0&\Sigma_2&0\\
 -\Sigma_1&0&\Sigma_1\end{pmatrix}.
\]
Let
\[
 D_P=\beta\beta^*+\sum_aT_aP_0T_a^*-P_0,\qquad
 D_Q=\chi\chi^*+\sum_aT_a^*Q_0T_a-Q_0,
\]
and take any $e_P\ge K_P\|D_P\|$, $e_Q\ge K_Q\|D_Q\|$.
Then
\begin{equation}
 \|H_f-H_g\|^2\le
 \max\{(\delta+e_P)(\delta+e_Q),\,
                    2\sigma_1\max(e_P,e_Q)+e_Pe_Q\}.
 \label{I:U:eq:reference-truncation}
\end{equation}
In particular, define
\[
 M_P=\Big\|\sum_aA_{a,12}A_{a,12}^*\Big\|,\quad
 M_Q=\Big\|\sum_aA_{a,21}^*A_{a,21}\Big\|,
\]
\[
 c_P=K_P(2M_P+\sqrt{M_P}),\qquad
 c_Q=K_Q(2M_Q+\sqrt{M_Q}).
\]
One may take $e_P=c_P\delta$, $e_Q=c_Q\delta$, giving
\begin{equation}
 \|H_f-H_g\|^2\le
 \max\{\delta^2(1+c_P)(1+c_Q),\,
             2\sigma_1\delta\max(c_P,c_Q)+c_Pc_Q\delta^2\}.
 \label{I:U:eq:reference-square-root}
\end{equation}
This is an $O(\sqrt\delta)$ bound when the other displayed quantities
are bounded uniformly. It contains no explicit tail-dimension factor,
but its coupling and resolvent constants may depend on that dimension.
\end{theorem}
\begin{proof}
Let $P_E,Q_E$ be the exact joint error Gramians. Subtracting the
reference Stein equations gives
$P_E-P_0=(I-\Phi_T)^{-1}D_P$, and the dual identity for $Q_E$.
The identity forcings have solutions
$\Udiag(X_A,X_B)$ and $\Udiag(Y_A,Y_B)$, respectively.
Lemma~\ref{I:U:lem:positive-resolvent} therefore gives
\[
 0\preceq P_E\preceq P_0+e_PI=:\overline P,\qquad
 0\preceq Q_E\preceq Q_0+e_QI=:\overline Q.
\]
For positive semidefinite matrices the quantity
$\lambda_{\max}(P^{1/2}QP^{1/2})$ is separately monotone in $P,Q$:
use congruence for $Q$, and the equality of the nonzero spectra of
$P^{1/2}QP^{1/2}$ and $Q^{1/2}PQ^{1/2}$ for $P$.
Apply Theorem~\ref{I:U:thm:exact}. The reference matrices commute.
On the retained symmetric and antisymmetric subspaces their pairs
of eigenvalues are $(2\sigma_j,0)$ and $(0,2\sigma_j)$; on a tail
coordinate they are $(\sigma_j,\sigma_j)$.
The largest eigenvalue of $\overline P\overline Q$ is exactly the
right-hand side of~\eqref{I:U:eq:reference-truncation}.

For the explicit estimates put
$S=\sum_aA_{a,12}\Sigma_2A_{a,12}^*$ and
$U=\sum_aA_{a,22}\Sigma_2A_{a,12}^*$.
Direct multiplication of the balanced Stein equations gives
\[
 D_P=-\begin{pmatrix}0&0&S\\0&0&U\\S&U^*&S\end{pmatrix}.
\]
Here $\|S\|\le M_P\delta$, and operator Cauchy--Schwarz, together
with $\sum_aA_{a,22}\Sigma_2A_{a,22}^*\preceq\Sigma_2$, gives
$\|U\|\le\sqrt{M_P}\delta$.
Split $D_P$ into its $S$ and $U$ parts; their norms are at most
$2\|S\|$ and $\|U\|$. Hence
$\|D_P\|\le(2M_P+\sqrt{M_P})\delta$.
Apply the same calculation to $A_a^*,c$, followed by changing the
sign of the last joint block, to obtain the asserted $D_Q$ bound.
\end{proof}


\paragraph{Correct balanced block inequalities.}
The identities imply
\[
 b_2b_2^*+\sum_aA_{a,22}\Sigma_2A_{a,22}^*\preceq\Sigma_2,
 \qquad
 \Big\|\sum_aA_{a,12}\Sigma_2A_{a,12}^*\Big\|
                  \le\min\{\sigma_1,M_P\delta\}.
\]
They do not permit replacing the norm of a positive sum by the sum
of its spectral norms, nor multiplying the two independent bounds.
For a concrete balanced counterexample let
\[
 A_0=E_{12},\quad A_1=E_{23},\quad
 b=\sqrt t\,e_3,\quad c=\sqrt t\,e_1,\qquad 0<t<1,
\]
and retain $k=1$. All products of three letters vanish, and the
reachability and observation Gramians are both $tI_3$; in particular
the representation is stable and minimal. The retained tuple is zero.
Nevertheless,
\[
 \|b_2\|^2+\sum_a\|A_{a,22}\Sigma_2^{1/2}\|^2=2t>t,
 \qquad
 \Big\|\sum_aA_{a,12}\Sigma_2A_{a,12}^*\Big\|=t>t^2.
\]
The reference-Gramian theorem uses neither false strengthening.

\section{A rational full-operator upper certificate in the commuting example}\label{I:S:app:commupper}
Use $b=13/24$, $c=197/200$, and the rank-one series $g(b^n)=cb^n$, zero on all words containing $a$. To avoid overloading the letter notation, the scalar $b$ in this sentence is the coefficient of the alphabet letter $b$. Take the orthogonal basis
\[
 e_\Seps,\ e_a,\ e_b,\ e_{aa},\
 h=\sum_{n\ge2}b^ne_{b^n}.
\]
Its metric matrix is $D=\Sdiag(1,1,1,1,d_0)$, $d_0=b^4/(1-b^2)$. The kernel has coordinates $k=(1,0,b,0,1)^\top$. Put $H_5=H\oplus0$ for the four-square target in Theorem~\ref{I:S:thm:commutative}. The error bilinear-form matrix is
$L=DH_5-cDkk^\top D$. The operator vanishes on the orthogonal complement of this five-dimensional space. Thus its norm is less than $s=29/25$ exactly when both $sD+L$ and $sD-L$ are positive definite. Their leading principal minors are:
\begin{center}\small
\begin{tabular}{@{}ccc@{}}
\toprule
Order & $sD+L$ & $sD-L$\\\midrule
1 & $47/40$ & $229/200$\\
2 & $1269/500$ & $229/1250$\\
3 & $3481227/2000000$ & $1037897/4500000$\\
4 & $13775541/100000000$ & $16067927/450000000$\\
5 & $75291604330033/79511520000000000$ & $131840555454959/1073405520000000000$\\
\bottomrule
\end{tabular}
\end{center}
All are strictly positive. These entries are exact rational values, and the displayed inequalities require no numerical search.
